\section{Kodaira--Spencer II: fixed points and assembly}\label{sec:pf:ks2}

\providecommand{\kstwoG}{G(\A_f^{p})}
\providecommand{\kstwoKS}{\mathcal{A}}
\providecommand{\kstwoOF}{\OO_{F}}
\providecommand{\kstwoLaHat}{\hat\Lambda^{p}}
\providecommand{\kstwoGp}{G_{p}}
\providecommand{\kstwoKp}{K_{p}}
\providecommand{\kstwoBall}{\mathrm{ball}}
\providecommand{\kstwoFib}{\mathrm{fib}}
\providecommand{\kstwoloc}{\mathrm{loc}}

% ---------------------------------------------------------------------------
\subsection{Introduction}\label{pf:ks2:intro}

This section proves \cref{pf:ks2:main}: \emph{the holomorphic Lefschetz formula for finite
\'etale correspondences (in the form of \cref{thm:galois:hlf}) implies the proposition
$\sfname{KSAll}$ of \cref{pf:ks2:def:ksall}.} This implication is one of the three inputs of the
proof of the geometric realization theorem \cref{thm:geoone}.

$\sfname{KSAll}$ is a package of six statements about \emph{Kuga--Sato towers}: towers of fibre
powers $\kstwoKS_N^k$ of the universal abelian schemes over the fine moduli schemes $M_N$ of
abelian varieties of Harris--Taylor type (PEL type, unitary of signature $(1,n-1)$ at one
archimedean place and definite at the others), equipped with the Hecke correspondences of double
cosets. The six statements say that standard bases exist ($\sfname{StdEx}$), that Kuga--Sato
towers exist ($\sfname{KSEx}$), that the associated Hecke towers are smooth, projective and have
finite \'etale correspondences ($\sfname{KSNice}$), that they have good reduction at places
where the Hecke data are unramified ($\sfname{KSGood}$), that they satisfy a pointwise
Eichler--Shimura congruence relation at such places ($\sfname{KSCong}$), and that a Hecke
eigensystem which occurs in algebraic modular forms on the definite inner forms occurs in the
Hodge cohomology $H^b(\kstwoKS_N^k,\Omega^{h})$ in prescribed degrees $h$ ($\sfname{KSHodge}$).
The letters ``KS'' in the formal names refer to Kuga--Sato varieties, not to Kodaira--Spencer
maps.

The last statement is proved in \lref{UpperTheory.Gal.LocalSmoothExtension.ksHodge_of_geo}
from the holomorphic Lefschetz formula and from the explicit description of the fixed points of
twisted Hecke correspondences on Kuga--Sato varieties, the proposition $\sfname{FixK}$
(\cref{pf:ks2:def:fixk}). The proof of $\sfname{FixK}$ is the main geometric content of this
section (\cref{pf:ks2:sec:fix}): the fixed points of a Hecke correspondence composed with fibrewise
multiplications are CM points; they are parametrized, through the complex uniformization of $M_N$,
by pairs consisting of a point of a double coset space of a definite-at-infinity similitude group
and a rational element fixing it; and their local Lefschetz contributions are computed from the
eigenvalues of the CM endomorphism on the tangent spaces of the base and of the fibres. The
weighted fixed-point count is thereby expressed as a finite sum of orbital type over the classes of
hermitian forms. The remaining five statements are deduced in
\cref{pf:ks2:sec:std,pf:ks2:sec:min,pf:ks2:sec:cong,pf:ks2:sec:good} from an integral input
($\sfname{MIn}$, \cref{pf:ks2:min}): integral models, their properness in the compact case,
\'etale-level models at auxiliary primes, and the pointwise congruence relation on Kuga--Sato
varieties over them, whose combinatorial heart (a count of subspaces of a vector space over a
finite field containing, or not containing, a distinguished line) is proved in full in
\cref{pf:ks2:sec:cong}.

\paragraph{What is proved in full, and what is cited.}
The top layers of all arguments are proved here: the assembly of $\sfname{KSAll}$, the
statements $\sfname{StdEx}$ and $\sfname{KSNice}$, the construction of Kuga--Sato towers from
Hecke witnesses, the reduction of the congruence relation to the combinatorial identity and the
identity itself, the deduction of $\sfname{KSGood}$ and $\sfname{KSCong}$ from sites, covers and
lying-over statements, and, for $\sfname{FixK}$, the reduction to the fine moduli schemes, the
change of level, the assembly of the fixed-point formula from the complex uniformization, and the
computation of the local contribution over a fixed point of the base. We also prove the top
layers of a number of intermediate results: the splitting and Galois properties of tower maps,
the comparison of generic models with generic fibres of models, normal $T_j$-correspondences,
quotients by subgroup schemes, the reading of Hecke correspondences at ordinary points and the
resulting embedding criterion, the extension of level-forgetting witnesses to \'etale-level
models, the $p$-adic and real dictionaries and the period coordinates, the extension of objects
over discrete valuation rings, and the classification of complex points of models. Deeper lemmas of
the formal proof (the analytic theory of complex tori, generic ordinariness of the special fibre,
existence of \'etale-level models and of the lying-over maps, the local inputs of the period
coordinates) are cited by cross-reference
to the formal development, with their content stated as we use it. Some of them are classical theorems formalized in
the development; for those we also give literature references
(\cref{pf:ks2:sec:classical}).

\paragraph{Organization.}
\Cref{pf:ks2:sec:setting} recalls the objects and states $\sfname{KSAll}$;
\cref{pf:ks2:sec:main} proves the main theorem modulo the components;
\cref{pf:ks2:sec:std} treats standard bases and the existence and niceness of Kuga--Sato towers;
\cref{pf:ks2:sec:min} the integral input; \cref{pf:ks2:sec:cong} the congruence relation;
\cref{pf:ks2:sec:good} good reduction and congruences on the towers; \cref{pf:ks2:sec:fix} the
fixed-point theorem $\sfname{FixK}$; \cref{pf:ks2:sec:classical} lists the classical inputs; and
\cref{pf:ks2:sec:remarks} collects remarks on the formal statements.

% ---------------------------------------------------------------------------
\subsection{The setting}\label{pf:ks2:sec:setting}

All schemes are schemes in the usual sense; $\mathrm{Sp}\,R=\Spec R$. For a morphism $f$ and an
integer $d\ge0$ we say that $f$ is \emph{finite \'etale of degree $d$} if $f$ is finite and
\'etale and every geometric point $x\colon\Spec\Omega\to X$ ($\Omega$ algebraically closed) of the
target has exactly $d$ lifts $y\colon\Spec\Omega\to Y$ with $f\circ y=x$.

\subsubsection{PEL data of Harris--Taylor type}

\begin{definition}[Data, groups and links]\label{pf:ks2:def:datum}
Let $p$ be a prime.
\begin{enumerate}
\item A \emph{datum} $D$ for $p$ consists of: a CM field
  $F$ with complex conjugation $c$; an element $\omega\in\kstwoOF$ and integers $\mathrm{tr},\mathrm{nm}$
  with $\omega+c(\omega)=\mathrm{tr}$, $\omega c(\omega)=\mathrm{nm}$, $\mathrm{tr}^2-4\mathrm{nm}<0$,
  $p\nmid\mathrm{tr}^2-4\mathrm{nm}$ and $p\mid r^2-\mathrm{tr}\,r+\mathrm{nm}$ for some $r\in\Z$
  (so $\Q(\omega)$ is imaginary quadratic and $p$ splits in it); an $\kstwoOF$-algebra $O$, free
  of rank $m^2[F:\Q]$ over $\Z$ ($m\ge1$), with centre $\kstwoOF$ up to nonzero integer multiples,
  an involution ${}^*$ inducing $c$ on $\kstwoOF$, and an isomorphism
  $O\otimes\Z_p\cong M_m(\kstwoOF\otimes\Z_p)$; an $O$-module $\Lambda$, free of rank
  $mn[F:\Q]$ over $\Z$ ($n\ge1$), with a nondegenerate alternating form $E_\Lambda$ satisfying
  $E_\Lambda(bx,y)=E_\Lambda(x,b^*y)$ and perfect up to an integer prime to $p$; an integer $a$, a root
  $\rho_c\in\C$ of $X^2-\mathrm{tr}X+\mathrm{nm}$ and distinct embeddings
  $\tau_1,\dots,\tau_a\colon F\to\C$ with $\tau_i(\omega)=\rho_c$; and a complex structure on
  $\Lambda\otimes\R$ compatible with $E$ and $O$ satisfying Kottwitz's determinant condition.
  Put $g=mn[F:\Q]/2$. $D$ is of \emph{compact type} if $2a<[F:\Q]$, or $O\otimes\Q$ is a division
  algebra and $n<2m$.
\item An \emph{HT-group} $G$ for $p$ consists of a CM
  field (again denoted $F$) with an imaginary quadratic subfield, a central simple $F$-algebra $C$
  of dimension $n^2$ with an involution of the second kind, archimedean splittings
  $\psi_v\colon C\otimes_{F,v}\C\cong M_n(\C)$ under which the involution becomes adjunction for
  a hermitian Gram matrix $H_v$, and local data at $p$. Its group of prime-to-$p$ finite adelic
  points is
  \[
   \kstwoG=\{(g,\nu)\in(\A_f^p\otimes_\Q C)^\times\times(\A_f^p)^\times:\ g^*g=\nu\},
  \]
  where $\A_f^p$ is the ring of finite adeles of $\Q$ away
  from $p$. The integer $G.a$ is the number of archimedean places $v$ at which $H_v$ is
  indefinite. We write $\kstwoGp$ for the $p$-adic group $\GL_1(\Q_p)\times\prod_{i<m}\GL_n(K_i)$
  attached to $G$ (with the local fields $K_i$ of the data of
  $G$) and $\kstwoKp=\Z_p^\times\times\prod_i\GL_n(\OO_{K_i})$ for its standard maximal compact
  subgroup.
\item A \emph{link} $L$ between $D$ and $G$ identifies the CM
  fields, sets $n_G=n_D$, realizes $C$ as the centralizer of $O$ in $\End_\Q(\Lambda\otimes\Q)$
  compatibly with the involutions, and matches the archimedean and $p$-adic data. Through $L$ the
  group $\kstwoG$ acts on $\A_f^p\otimes\Lambda$ (we write $h\cdot v$); put
  $\kstwoLaHat=\hat\Z^p\otimes\Lambda$. An element $h$ is \emph{integral} if $h\cdot\kstwoLaHat\subseteq\kstwoLaHat$. For
  $N\ge1$, $K_N\subseteq\kstwoG$ is the group of $g$ such
  that $g$ and $g^{-1}$ act trivially on $\kstwoLaHat/N\kstwoLaHat$. For $h\in\kstwoG$ and levels
  $N',N$, $h$ is \emph{admissible from $N'$ to $N$} if $hK_{N'}h^{-1}\subseteq K_N$.
\item The \emph{standing conditions} $\mathrm{Sc}(D,G)$ are:
  $a_D=1$, $m=1$, $G.a=1$, $n$ odd, and $D$ of compact type.
\end{enumerate}
A \emph{level} is an integer $N\ge3$ prime to $p$.
\formal{GlobalTheory.UPel.HT.Datum}\formal{GlobalTheory.UPel.HT.Datum.g}\formal{GlobalTheory.UPel.HT.Datum.CompactB}\formal{GlobalTheory.LKTrip.HTGroup}\formal{GlobalTheory.LKTrip.HTGroup.GAf}\formal{GlobalTheory.LKRes.GHT}\formal{GlobalTheory.LKRes.KHT}\formal{GlobalTheory.LKGeo.Link}\formal{UpperTheory.Gal.Int.Geo.Integ}\formal{GlobalTheory.LKGeo.Link.KN}\formal{UpperTheory.Gal.G3.AdmK}\formal{UpperTheory.Gal.Int.Geo.Sc}\formal{UpperTheory.Gal.Int.Geo.Lv}
\end{definition}

\begin{definition}[Bases]\label{pf:ks2:def:base}
Let $D$ be a datum for $p$ and $R$ a commutative ring.
\begin{enumerate}
\item A \emph{base} $b$ of $D$ over $R$ consists of an
  element $\rho\in R$ with $\rho^2-\mathrm{tr}\rho+\mathrm{nm}=0$ and ring maps
  $t_1,\dots,t_a\colon\kstwoOF\to R$ with $t_i(\omega)=\rho$, such that every integer prime to
  $p$ is a unit in $R$, the $t_i$ are pairwise comaximal (the elements $t_i(x)-t_j(x)$ generate
  the unit ideal for $i\ne j$), and for $i\neq j$ and every ring map $\beta$ from $R$ to a field of
  characteristic $p$ the kernels of $\beta\circ t_i$ and $\beta\circ t_j$ differ.
\item If $R$ is a $\Q$-algebra, a \emph{characteristic-zero base} is the same without the two conditions involving $p$.
\item For a number field $E$ and a ring isomorphism $e\colon F\to E$, a characteristic-zero base
  $g_b$ over $E$ is \emph{standard for $e$} if
  $\rho=e(\omega)$ and $t_i(x)=e(x)$ for all $i$ and $x\in\kstwoOF$. Over $\C$, $g_b$ is
  \emph{standard} if $\rho=\rho_c$ and
  $t_i=\tau_i|_{\kstwoOF}$.
\end{enumerate}
\formal{GlobalTheory.UPel.HT.Base}\formal{GlobalTheory.UPel.HT.Gen.GBase}\formal{UpperTheory.Gal.Int.Geo.IsStd}\formal{UpperTheory.Gal.LocalSmoothExtension.StdC}
\end{definition}

\begin{definition}[Moduli problems and models]\label{pf:ks2:def:model}
Let $b$ be a base over $R$ and $N\ge1$. For an $R$-scheme $s\colon S\to\Spec R$, an
\emph{object of level $N$ over $S$} consists of an abelian
scheme $A/S$ with an action of $O$, a symmetric line bundle on $A\times_SA$, rigidified along both
zero sections, relatively ample on the diagonal, whose polarization has kernel killed by an
integer prime to $p$ and whose Rosati involution induces ${}^*$ on $O$, Kottwitz's condition on
$\Lie A$ with respect to $t_1,\dots,t_a$, and an $O$-linear level-$N$ structure
$\eta\colon\Lambda\to A(S)$ which at every geometric point is the reduction of a prime-to-$p$
trivialization of the Tate modules compatible with $O$ and with the Weil pairing up to a unit.
Morphisms of objects over a morphism $f\colon S'\to S$ are cartesian squares respecting all
structures. A \emph{model of level $N$} is a quasi-projective $R$-scheme $M$, smooth of relative
dimension $a(n-1)$, with an object $\mathcal A$ over $M$ which is universal: every object over
every $R$-scheme $S$ is the pullback of $\mathcal A$ along a unique morphism $S\to M$. For a
characteristic-zero base $g_b$ the analogous notions are \emph{generic objects} and \emph{generic models}; the definition of a generic model requires only the
universal property (no quasi-projectivity or smoothness).
\formal{GlobalTheory.UPel.HT.Obj}\formal{GlobalTheory.UPel.HT.ObjMap}\formal{GlobalTheory.UPel.HT.Model}\formal{GlobalTheory.UPel.HT.Gen.GObj}\formal{GlobalTheory.UPel.HT.Gen.GenModel}
\end{definition}

\begin{definition}[Hecke isogenies and witnesses]\label{pf:ks2:def:wit}
Let $g_b$ be a characteristic-zero base over $R$, $h\in\kstwoG$, and let $o',o$ be generic
objects over the same $S$, of levels $N',N$. A homomorphism $\psi\colon A_{o'}\to A_o$ is a
\emph{Hecke isogeny for $h$} if it is $O$-linear, has an
$O$-linear quasi-inverse $\psi'$ with $\psi'\psi=d$ and $\psi\psi'=d$ for an integer $d$ prime to
$p$, pulls the polarization bundle of $o$ back to the $c$-th power of that of $o'$ for an integer
$c$ prime to $p$ such that the similitude factor of $h$ is $c$ times a unit of $\hat\Z^p$, and at
every geometric point maps the level structures according to the action of $h$: for $n$ prime to
$p$ and $x,y\in\Lambda$, $\psi$ maps the image of $x$ in the $n$-torsion to the image of $y$ iff
$h\cdot x\equiv y\bmod n\kstwoLaHat$. For generic models $M',M$ of levels $N',N$, a \emph{Hecke
witness} for $h$ is a pair $(\pi,\Phi)$ of a morphism
$\pi\colon M'\to M$ over $R$ and a morphism $\Phi\colon\mathcal A'\to\mathcal A$ of universal
abelian schemes covering $\pi$ which factors as a Hecke isogeny for $h$ from $\mathcal A'$ to a
generic object $o$ over $M'$, followed by a morphism of objects $o\to\mathcal A$ over $\pi$.
\formal{UpperTheory.Gal.Int.Geo.IsTrans}\formal{UpperTheory.Gal.G3.Wit}
\end{definition}

\subsubsection{Kuga--Sato towers and Hecke towers}

\begin{definition}[Towers]\label{pf:ks2:def:tower}
Let $B$ be a scheme and $\Gamma$ a group. A \emph{tower} over $B$ for $\Gamma$ consists of an index set $I$, subgroups $K_i\subseteq\Gamma$
and a submonoid $\Delta\supseteq\bigcup K_i$, $B$-schemes $q_i\colon Z_i\to B$, and morphisms
$\mathrm{tr}(i,j,h)\colon Z_i\to Z_j$ over $B$ for $h\in\Delta$ with $hK_ih^{-1}\subseteq K_j$,
such that: (i) for every $j$ and every finite set of pairs $(i,x)$ there is $m$ with
$K_m\subseteq K_j$ normalized by $K_j$ and $xK_mx^{-1}\subseteq K_i$ for all the pairs;
(ii) all relative indices $[K_j:K_i\cap K_j]$ are finite; (iii) $\mathrm{tr}(i,i,k)=\mathrm{id}$ for
$k\in K_i$ and $\mathrm{tr}(j,m,h')\circ\mathrm{tr}(i,j,h)=\mathrm{tr}(i,m,h'h)$; (iv) every
$\mathrm{tr}(i,j,1)$ is finite, \'etale and surjective; (v) if $K_i$ is normalized by $K_j$, two
geometric points of $Z_i$ with the same image under $\mathrm{tr}(i,j,1)$ differ by
$\mathrm{tr}(i,i,k)$ for some $k\in K_j$; (vi) if $K_i\subseteq K_j$, $k\in K_j$ and
$\mathrm{tr}(i,i,k)$ fixes a geometric point, then $k\in K_i$.
\formal{UpperTheory.Gal.G3.Tower}
\end{definition}

\begin{definition}[Kuga--Sato towers]\label{pf:ks2:def:ks}
Let $D,G,L$ be as above, $R$ a field of characteristic $0$, $g_b$ a characteristic-zero base over
$R$ and $k\in\N$. A \emph{Kuga--Sato tower} for $(L,g_b,k)$
is a tower $T$ over $\Spec R$ for $\kstwoG$, indexed by the levels $N$, with $K_N$ the principal
congruence subgroups and $\Delta$ the monoid of integral elements, together with generic models
$M_N$ of level $N$ with universal abelian schemes $\kstwoKS_N\to M_N$, and isomorphisms
$Z_N\cong\kstwoKS_N^k=\kstwoKS_N\times_{M_N}\dots\times_{M_N}\kstwoKS_N$ ($k$ factors) over
$\Spec R$, such that every transition map $\mathrm{tr}(N',N,h)$ corresponds to the $k$-th fibre
power $\Phi^{k}$ of the second component of a Hecke witness $(\pi,\Phi)$ for $h$ from $M_{N'}$
to $M_N$.
\formal{UpperTheory.Gal.Int.Geo.KS}
\end{definition}

\begin{definition}[Hecke data and Hecke towers]\label{pf:ks2:def:heckedata}
Let $N$ be a level.
\begin{enumerate}
\item A \emph{Hecke datum of level $N$} is a triple
  $(K_g,h_g,M_g)$ of a subgroup $K_N\subseteq K_g\subseteq K_1$, an integral element $h_g$, and a
  level $M_g$ divisible by $N$ such that $h_gxK_{M_g}x^{-1}h_g^{-1}\subseteq K_N$ for all
  $x\in K_g$.
\item A \emph{Hecke tower} over a field $k$ consists of an
  index set $\Gamma$, a $k$-scheme $q\colon Z\to\Spec k$ and correspondences
  $Z\xleftarrow{s_\gamma}Y_\gamma\xrightarrow{t_\gamma}Z$ over $\Spec k$. It is \emph{nice} if $q$ is smooth and projective (a closed immersion into
  some $\mathbb P^N_k$ over $k$ followed by the projection) and every $s_\gamma$ is finite \'etale
  of some degree.
\item For a Kuga--Sato tower $T$ and a family $(g_j)_{j\in J}$ of Hecke data of level $N$, the
  Hecke tower $T.\mathrm{tower}(N,g)$ has $\Gamma=J$,
  $Z=Z_N$, $Y_j=\coprod_{c\in K_{g_j}/(K_N\cap K_{g_j})}Z_{M_{g_j}}$, $s_j$ equal to
  $\mathrm{tr}(M_{g_j},N,1)$ on every copy, and $t_j$ equal to $\mathrm{tr}(M_{g_j},N,h_{g_j}x_c)$
  on the copy indexed by $c$, where $x_c\in K_{g_j}$ is a chosen representative of $c$.
\end{enumerate}
\formal{UpperTheory.Gal.Int.Geo.Mem}\formal{GlobalTheory.GalV.HTower}\formal{UpperTheory.Gal.Int.Geo.Nice}\formal{UpperTheory.Gal.Int.Geo.KS.tower}
\end{definition}

\subsubsection{Integral structures}

\begin{definition}[Covers, lying over, good covers]\label{pf:ks2:def:cover}
Let $H$ be a Hecke tower over a field $k$, $\theta\colon\Gamma\to A$ a map to a field, and $O\to
k$ a ring map. A \emph{cover} consists of a smooth
projective $p_Y\colon Y\to\Spec O$ of some relative dimension and a morphism $\pi\colon Z\to Y$
compatible with $q$ and $\Spec k\to\Spec O$, such that the induced morphism $Z\to Y\times_O k$ is
finite \'etale and Galois for a finite group $\Delta$ acting on $Z$ (the action map
$\coprod_{\Delta}Z\to Z\times_{Y_k}Z$ is an isomorphism), each $\delta\in\Delta$ being realized by
a correspondence $\gamma$ of $H$ with $s_\gamma$ an isomorphism, $t_\gamma=\delta\circ s_\gamma$
and $\theta(\gamma)=1$. For a correspondence $Y\xleftarrow{s}T\xrightarrow{t}Y$ and $m\ge1$,
$\gamma$ \emph{lies over} $(s,t)$ with multiplicity $m$
if there is $\kappa\colon Y_\gamma\to T$ with $s\kappa=\pi s_\gamma$ and $t\kappa=\pi t_\gamma$
such that the induced morphism $Y_\gamma\to Z\times_YT$ is finite \'etale of degree $m$. The cover
is \emph{good} if every $\gamma$ lies over some
correspondence $(s,t)$ on $Y$ over $O$ with $s$ finite \'etale of some degree.
\formal{UpperTheory.Gal.Int.Geo.Cover}\formal{UpperTheory.Gal.Int.Geo.Cover.Lies}\formal{UpperTheory.Gal.Int.Geo.Cover.Good}
\end{definition}

\begin{definition}[Pointwise congruence relation]\label{pf:ks2:def:ptid}
Let $\ell,q,n\in\N$, $Y$ a scheme and $\varpi\in\Gamma(Y,\OO_Y)$.
\begin{enumerate}
\item A \emph{correspondence of $Y$ away from $\ell$}
  is $Y\xleftarrow{p}Z\xrightarrow{r}Y$ with $p$ finite, \'etale over $Y[1/\ell]$, and with a
  constant number $d$ of geometric points in the fibres of $p$ over geometric points where
  $\ell$ is invertible.
\item A point $x\in Y$ is a \emph{generic point of $\{\varpi=0\}$} if $\varpi(x)=0$ and $x$ has no proper
  generization in $\{\varpi=0\}$.
\item For a set $P$, a map $\phi\colon P\to P$, finite index sets $I_0,\dots,I_n$ and maps
  $v_j\colon I_j\to P$, the \emph{alternating identity} $\mathrm{Alt}(q,n,\phi,v)$ is the equality of multisets
  \begin{equation}\label{pf:ks2:eq:altid}
  \sum_{j\ \mathrm{even}}q^{\binom j2}\bigl\{\phi^{\,n-j}(v_j(i)):i\in I_j\bigr\}
  =\sum_{j\ \mathrm{odd}}q^{\binom j2}\bigl\{\phi^{\,n-j}(v_j(i)):i\in I_j\bigr\}.
  \end{equation}
\item For correspondences $T_0,\dots,T_n$ of $Y$ away from $\ell$, the \emph{pointwise congruence
  relation} $\mathrm{PC}(\ell,q,n,\varpi,T)$ holds if: for every valuation ring $V$ with
  algebraically closed fraction field and $\ell\ne0$ in $V$, every $y\colon\Spec V\to Y$ mapping
  the closed point to a generic point of $\{\varpi=0\}$, the $q$-power Frobenius $\phi$ of the
  residue field $\kappa$ of $V$, and every enumeration of the sections $z$ of $T_j$ over $y$
  (lifts $z\colon\Spec V\to Z_j$ of $y$), the identity $\mathrm{Alt}$ holds for the residue
  points $\bar z=r_j\circ z|_{\Spec\kappa}\in Y(\kappa)$, with $\phi$ acting on $Y(\kappa)$ by
  precomposition.
\end{enumerate}
\formal{UpperTheory.Gal.crystallineComparison.Corr}\formal{UpperTheory.Gal.crystallineComparison.GenPt}\formal{UpperTheory.Gal.crystallineComparison.AltId}\formal{UpperTheory.Gal.crystallineComparison.PtId}
\end{definition}

\begin{definition}[Good reduction and congruence at a place]\label{pf:ks2:def:goodatw}
Let $E$ be a number field, $w$ a finite place with valuation ring $O_w$, $H$ a Hecke tower over
$E$ and $\theta\colon\Gamma\to A$.
\begin{enumerate}
\item $(H,\theta)$ has \emph{good reduction at $w$}, written $\mathcal{G}(H,\theta,w)$, if there is a
  good cover of $(H,\theta)$ over $O_w$.
\item For $n\in\N$ and $\vartheta\colon\N\to A$, $(H,\theta)$ satisfies the \emph{congruence relation
  at $w$ with coefficients $\vartheta$}, written $\mathcal{C}(H,\theta,w,n,\vartheta)$, if there are a
  cover $(Y,\pi)$ over $O_w$, a prime
  $q\in w$ with $\Nm w=q^f$, a uniformizer $\varpi$ of $O_w$, correspondences $T_0,\dots,T_n$ of
  $Y$ away from $q$ over $O_w$ with $\mathrm{PC}(q,q^f,n,\varpi,T)$, and elements
  $\gamma_0,\dots,\gamma_n\in\Gamma$ and integers $m_j\ge1$ such that $\gamma_j$ lies over $T_j$
  with multiplicity $m_j$ and $\theta(\gamma_j)=m_j\vartheta_j$.
\end{enumerate}
\formal{UpperTheory.Gal.Int.Geo.GoodAtW}\formal{UpperTheory.Gal.Int.Geo.CongAtW}\formal{UpperTheory.Gal.Int.Geo.Cover.Cong}
\end{definition}

\begin{definition}[Group-theoretic conditions at a prime]\label{pf:ks2:def:goodgrp}
Let $D,G,L$ be as above, $N$ a level, $(g_j)_{j\in J}$ Hecke data of level $N$ and
$\theta\colon J\to A$. For a prime $q$ and $r\in\Z$ put $\mathfrak u=(q,\omega-r)\subseteq\kstwoOF$.
\begin{enumerate}
\item $(q,r,\mathfrak w)$ is \emph{admissible}, written $\mathrm{Adm}(q,r,\mathfrak w)$, if $q\ne p$ is a
  prime, $q\nmid\mathrm{tr}^2-4\mathrm{nm}$, $q\mid r^2-\mathrm{tr}\,r+\mathrm{nm}$, $E_\Lambda$ is
  perfect up to an integer prime to $q$, and $\mathfrak w\supseteq\mathfrak u$ is a maximal ideal.
  An \emph{auxiliary level} for $(q,r,\mathfrak w)$ is a pair $(N_0,\mathfrak a)$ with $N_0\ge3$ prime
  to $pq$, $\mathfrak a+\mathfrak w=\kstwoOF$ and $\mathfrak a\supseteq\mathfrak u^e$ for some $e$; an
  \emph{auxiliary ideal} is an ideal of the form $N_0\mathfrak a$ for an auxiliary level.
\item For an ideal $\mathfrak b$, $K(\mathfrak b)\subseteq\kstwoG$, the \emph{level group of $\mathfrak b$},
  consists of the $g\in K_1$ acting trivially modulo $\mathfrak b$, i.e.\ with $g\cdot v-v\in\mathfrak
  b\kstwoLaHat$ for $v\in\kstwoLaHat$.
\item $(g,\theta)$ satisfies the \emph{deck condition} for $K$
  if for every $x\in K$ some $j$ has $K_{g_j}=K_N$, $M_{g_j}=N$, $h_{g_j}=x$ and $\theta_j=1$.
\item $h$ is a \emph{$T_i$-element at $\mathfrak w$} if
  $h$ is trivial modulo every $n$ prime to $q$ and modulo every ideal $\mathfrak a$ prime to
  $\mathfrak w$ containing a power of $\mathfrak u$, maps $\kstwoLaHat$ into itself with
  $h\kstwoLaHat+\mathfrak w\kstwoLaHat$ of index $\Nm\mathfrak w^{\,i}$ in $\kstwoLaHat$ and
  $\Nm\mathfrak w^{\,i}\kstwoLaHat\subseteq h\kstwoLaHat$, kills $\mathfrak w$ modulo an element
  $s\notin\mathfrak w$ (precisely: for $x\in\mathfrak w$, $sx\kstwoLaHat\subseteq h\kstwoLaHat$),
  and has similitude factor $\Nm\mathfrak w^{\,i}$ times a unit.
\item The Hecke data are \emph{unramified at $\mathfrak w$}, written $\mathcal{U}(g,\theta,\mathfrak w)$,
  if there are $q,r$ with $\mathrm{Adm}(q,r,\mathfrak w)$ and an auxiliary level $(N_0,\mathfrak a)$ such
  that $K_0=K(N_0\mathfrak a)$ contains $K_N$, $N\in N_0\mathfrak a$, the deck condition holds for $K_0$,
  and for every $j$ there are $\kappa\in\N$, an auxiliary level $(N_0\kappa,\mathfrak a')$ with
  $\mathfrak a'\subseteq\mathfrak a$, integral elements $h_1,h_2$ with $h_{g_j}=h_1h_2$, and a level
  $N_1$, such that, with $K_0'=K(N_0\kappa\mathfrak a')$: $K_{M_{g_j}}\subseteq K_0'$, $M_{g_j}\in
  N_0\kappa\mathfrak a'$, every $u\in K_0$ is $xu'$ with $x\in K_{g_j}\cap K_0$ and $u'\in K_0'$,
  $\kappa\kstwoLaHat\subseteq h_2\kstwoLaHat$, $N_1\in N_0\mathfrak a'$, $h_2xK_{M_{g_j}}(h_2x)^{-1}\subseteq
  K_{N_1}$ for $x\in K_{g_j}$, $h_1K_{N_1}h_1^{-1}\subseteq K_N$, and either $h_1=1$ or $h_1$ is a
  $T_i$-element at a maximal ideal $\mathfrak w'\supseteq\mathfrak u$ with $\mathfrak a'\subseteq\mathfrak
  a\mathfrak w'$, $i\le n$ and $\Nm\mathfrak w'^{\,i}\mid N_1$.
\item The Hecke data \emph{contain the operators $T_i$ at $\mathfrak w$ with eigenvalues $\vartheta$},
  written $\mathcal{H}(g,\theta,\mathfrak w,\vartheta)$, if
  there are $q,r$ with $\mathrm{Adm}(q,r,\mathfrak w)$, an auxiliary ideal $\mathfrak n$ with
  $K_0=K(\mathfrak n)\supseteq K_N$ and
  $N\in\mathfrak n$, the deck condition for $K_0$, and for every $0\le i\le n$ an index $j$ with
  $K_{g_j}=K_0$, $h_{g_j}$ a $T_i$-element at $\mathfrak w$, $\Nm\mathfrak w^{\,i}\mid M_{g_j}$,
  $[K_0:\Gamma_j]$ equal to the number of $i$-dimensional subspaces of
  $(\kstwoOF/\mathfrak w)^n$, and $\theta_j=[\Gamma_j:\Gamma_j\cap K_{M_{g_j}}]\,\vartheta_i$, where
  $\Gamma_j=K_0\cap h_{g_j}^{-1}K_0h_{g_j}$.
\item A place $w$ of a number field $E$ \emph{corresponds} to $\mathfrak w$ under $e\colon F\cong
  E$ if $x\in\mathfrak w\iff e(x)\in w$ for
  $x\in\kstwoOF$.
\end{enumerate}
\formal{UpperTheory.Gal.Int.Geo.uId}\formal{UpperTheory.Gal.Int.Geo.AtPrime}\formal{UpperTheory.Gal.Int.Geo.IsLev}\formal{UpperTheory.Gal.Int.Geo.CongId}\formal{UpperTheory.Gal.Int.Geo.Deck}\formal{UpperTheory.Gal.Int.Geo.IsTjElt}\formal{UpperTheory.Gal.Int.Geo.GoodGrp}\formal{UpperTheory.Gal.Int.Geo.LevAt}\formal{UpperTheory.Gal.Int.Geo.CongGrp}\formal{UpperTheory.Gal.Int.Geo.LevOK}\formal{UpperTheory.Gal.Int.Geo.IsPl}
\end{definition}

\subsubsection{The statement}

\begin{definition}[$\sfname{KSAll}$]\label{pf:ks2:def:ksall}
$\sfname{KSAll}$ is the conjunction of the following
statements, in which $p$ is a prime and $D,G,L$ are as in \cref{pf:ks2:def:datum}.
\begin{enumerate}[label=(\alph*)]
\item $\sfname{StdEx}$: if $a_D=1$, then for every number
  field $E$ and every ring isomorphism $e\colon F\to E$ there is a characteristic-zero base over $E$
  which is standard for $e$.
\item $\sfname{KSEx}$: if $\mathrm{Sc}(D,G)$, then for every
  field $R$ of characteristic $0$, every characteristic-zero base $g_b$ over $R$ and every $k$,
  a Kuga--Sato tower for $(L,g_b,k)$ exists.
\item $\sfname{KSNice}$: if $\mathrm{Sc}(D,G)$, then
  $T.\mathrm{tower}(N,g)$ is nice for every Kuga--Sato tower $T$ over a characteristic-zero field,
  every level $N$ and every family $g$ of Hecke data of level $N$.
\item $\sfname{KSGood}$: if $\mathrm{Sc}(D,G)$, $E$ is a
  number field, $g_b$ is standard for $e\colon F\cong E$, $T$ is a Kuga--Sato tower for
  $(L,g_b,k)$, $N$ a level, $g$ a family of Hecke data with values $\theta$ in a field, $w$ a place
  of $E$ corresponding to $\mathfrak w$, and $\mathcal{U}(g,\theta,\mathfrak w)$, then
  $\mathcal{G}(T.\mathrm{tower}(N,g),\theta,w)$.
\item $\sfname{KSCong}$: in the same situation,
  $\mathcal{H}(g,\theta,\mathfrak w,\vartheta)$ implies
  $\mathcal{C}(T.\mathrm{tower}(N,g),\theta,w,n,\vartheta)$.
\item $\sfname{KSHodge}$: the Hodge occurrence statement
  recalled in \lref{UpperTheory.Gal.LocalSmoothExtension.ksHodge_of_geo}: for every
  $\mathrm{Sc}(D,G)$ and every antitone weight $\lambda$ there is $k$ such that, for all Kuga--Sato
  towers $T$ for $(L,g_b,k)$ over characteristic-zero fields, eigensystems $\theta$ of Hecke data
  satisfying the spectral condition of loc.\ cit., ring maps $\tau$ from the base field to $\C$
  with $\tau(\rho)=\rho_c$ and $\tau\circ t_i=\tau_i$, $\sigma\in\Aut(\C)$ and Hodge exponent vectors
  $h$ of $\sigma\lambda$, the eigensystem $\sigma\circ\theta$ occurs in the Hodge cohomology of the
  base change of $T.\mathrm{tower}(N,g)$ along $\tau$ in each degree $h_i$.
\end{enumerate}
\formal{UpperTheory.Gal.Int.Geo.KSAll}\formal{UpperTheory.Gal.Int.Geo.StdEx}\formal{UpperTheory.Gal.Int.Geo.KSEx}\formal{UpperTheory.Gal.Int.Geo.KSNice}\formal{UpperTheory.Gal.Int.Geo.KSGood}\formal{UpperTheory.Gal.Int.Geo.KSCong}\formal{UpperTheory.Gal.Int.Geo.KSHodge}
\end{definition}

The holomorphic Lefschetz formula enters through the proposition $\sfname{HLF}$: the statement of
\cref{thm:galois:hlf} for a smooth projective $Z$ of relative dimension $\delta$ over $\C$,
finite \'etale $c_1,c_2\colon Y\rightrightarrows Z$, a finite set $\mathcal F$ which is exactly the
set of complex points $y$ with $c_1(y)=c_2(y)$, and matrices $A_y$ given by the tangent data of
\cref{pf:ks2:def:istan} below with $\det(1-A_y)\ne0$.

% ---------------------------------------------------------------------------
\subsection{The main theorem}\label{pf:ks2:sec:main}

\begin{theorem}\label{pf:ks2:main}
$\sfname{HLF}\Rightarrow\sfname{KSAll}$: the holomorphic Lefschetz formula for finite \'etale
correspondences implies $\sfname{KSAll}$.
\formallabel{UpperTheory.Gal.Int.Geo.kodaira_spencer_of_weight_one}
\end{theorem}

\begin{proof}
Assume the holomorphic Lefschetz formula. The six components of $\sfname{KSAll}$ are:
$\sfname{StdEx}$ by \cref{pf:ks2:stdex}; $\sfname{KSEx}$ by \cref{pf:ks2:ksex}; $\sfname{KSNice}$
by \cref{pf:ks2:ksnice}; $\sfname{KSGood}$ by \cref{pf:ks2:ksgood}; $\sfname{KSCong}$ by
\cref{pf:ks2:kscong}; and $\sfname{KSHodge}$ by
\lref{UpperTheory.Gal.LocalSmoothExtension.ksHodge_of_geo}, which states that
$\sfname{HLF}$ and $\sfname{FixK}$ together imply $\sfname{KSHodge}$,
applied to the assumption and to \cref{pf:ks2:fixk}. Only the last component uses the
assumption.
\end{proof}

% ---------------------------------------------------------------------------
\subsection{Standard bases, Kuga--Sato towers and niceness}\label{pf:ks2:sec:std}

\begin{proposition}\label{pf:ks2:stdex}
$\sfname{StdEx}$ holds.
\formallabel{UpperTheory.standard_datum_exists}
\end{proposition}

\begin{proof}
Let $a_D=1$ and $e\colon F\to E$ an isomorphism. Put $\varphi=e|_{\kstwoOF}$, $\rho=\varphi(\omega)$
and $t_1=\varphi$. Applying $\varphi$ to $\omega^2-\mathrm{tr}\,\omega+\mathrm{nm}=
\omega^2-(\omega+c\omega)\omega+\omega\,c\omega=0$ shows that $\rho$ is a root of
$X^2-\mathrm{tr}X+\mathrm{nm}$; $t_1(\omega)=\rho$; and comaximality is vacuous because there is
only one index. This base is standard for $e$ by construction.
\end{proof}

In the formal proof this proposition is derived from the integral input $\sfname{MIn}$ of
\cref{pf:ks2:min} by a lemma which does not use that hypothesis (\cref{pf:ks2:rem:unused}).

\subsubsection{Generic models}

We first record two facts on fibre powers and the comparison of generic models with generic fibres
of models.

\begin{lemma}[Smoothness of fibre powers]\label{pf:ks2:powsmooth}
Let $A\to X$ be a morphism of schemes which is smooth of relative dimension $g$, and $k\in\N$. Then
the $k$-fold fibre power $A^k=A\times_X\dots\times_XA\to X$ is smooth of relative dimension $kg$.
\formallabel{GlobalTheory.FibPow.smoothOfRelativeDimension_pow}
\end{lemma}

\begin{proof}
Induction on $k$. For $k=0$, $A^0=X$, and the identity is smooth of relative dimension $0$. For
$k\ge0$, $A^{k+1}=A^k\times_XA$, and the projection $A^{k+1}\to A^k$ is the base change of $A\to X$
along $A^k\to X$, hence smooth of relative dimension $g$. Its composite with $A^k\to X$, which is
smooth of relative dimension $kg$ by induction, is smooth of relative dimension $kg+g$, since
relative dimensions add under composition of smooth morphisms. (In the formal proof this is
\lref{GlobalTheory.smoothOfRelativeDimension_powObj} for the iterated fibre product, transported
along its canonical isomorphism with the fibre power.)
\end{proof}

\begin{lemma}[Projectivity of fibre powers]\label{pf:ks2:powproj}
Let $R$ be a Noetherian ring, $X\to\Spec R$ projective, $A\to X$ an abelian scheme and $L$ a line
bundle on $A$ which is relatively ample over $X$. Then for every $k\in\N$ the composite
$A^k\to X\to\Spec R$ is projective.
\formallabel{GlobalTheory.KS.isProjOver_pow}
\end{lemma}

\begin{proof}
Over a Noetherian ring, a morphism to $\Spec R$ is projective if and only if it is proper and its
source carries an ample line bundle
(\lref{GlobalTheory.isProjOver_iff_isProper_and_exists_isAmpleLB}). So $X\to\Spec R$ is proper
and $X$ carries an ample line bundle. Let $\mathrm{pr}_1,\dots,\mathrm{pr}_k\colon A^k\to A$ be the
projections. The external tensor product $\bigotimes_i\mathrm{pr}_i^*L$ is relatively ample over $X$
(\lref{GlobalTheory.relAmple_boxLB}: external tensor products of relatively ample line bundles on
the factors of a fibre power are relatively ample). Hence $A^k$ carries an ample line bundle
(\lref{GlobalTheory.RelAmpleQProj.isProper_and_exists_isAmpleLB}: a scheme with a relatively ample
line bundle over a proper base with an ample line bundle carries an ample line bundle). The morphism
$A^k\to X$ is proper, as an iterated fibre product of copies of the proper morphism $A\to X$ (proper
morphisms are stable under base change and composition), so $A^k\to\Spec R$ is proper. By the
criterion above it is projective.
\end{proof}

\begin{proposition}[Generic models and generic fibres]\label{pf:ks2:genmodel}
Let $R$ be a Noetherian ring with a base $b$ of $D$, $\varphi\colon R\to R'$ a ring map to a
$\Q$-algebra, and $g_b$ a characteristic-zero base of $D$ over $R'$ with $\rho=\varphi(\rho_b)$ and
$t_i=\varphi\circ t_{b,i}$ for all $i$. Let $N$ be a level, $M$ a model of level $N$ over $R$ and $X$
a generic model of level $N$ over $R'$, and put $M_{R'}=M\times_{\Spec R}\Spec R'$ with projections
$\mathrm{pr}_1,\mathrm{pr}_2$. Then there is a morphism $j\colon X\to M_{R'}$ which is an open and a
closed immersion, satisfies $\mathrm{pr}_2\circ j=s_X$ (the structure map of $X$), and under which
the universal generic object $\mathcal A_X$ of $X$ is the pullback of the universal object
$\mathcal A_M$ of $M$ along $\mathrm{pr}_1\circ j$.
\formallabel{GlobalTheory.UPel.HT.GenCmp.genModel_eq_genericFibre}
\formallabel{GlobalTheory.T.exists_open_closed_immersion_genModel_to_fibre}
\end{proposition}

\begin{proof}
We show that $j$ can be chosen to be an isomorphism. Two remarks on the moduli problems come
first. (a) Let $S$ be an $R'$-scheme, viewed as an $R$-scheme through $\varphi$. Since
$\rho=\varphi(\rho_b)$ and $t_i=\varphi\circ t_{b,i}$, objects of level $N$ for $b$ over $S$ are the
same as objects for $g_b$, viewed as a base over the $\Q$-algebra $R'$. Such an object is a generic
object, because at every geometric point it admits a trivialization of the full Tate module,
compatible with $O$ and with the Weil pairing up to a unit, which induces its level structure
(\lref{GlobalTheory.UPel.HT.GenCmp.geomFull}: in characteristic $0$ the $p$-adic Tate module,
with its $O$-action and Weil pairing, is isomorphic to $\Lambda\otimes\Z_p$ as well). (b)
Conversely, a generic object over $S$ is an object for $b$, the only additional condition being
that the kernel of its polarization is killed by an integer prime to $p$
(\lref{GlobalTheory.UPel.HT.GenCmp.genericObject_kernelCondition}: the trivialization of the
full Tate module identifies the Weil pairing with $E_\Lambda$ up to a unit, and $E_\Lambda$ is perfect
up to an integer prime to $p$).

Let $o_S$ be the pullback of $\mathcal A_M$ along $\mathrm{pr}_1$ (objects pull back along arbitrary
morphisms, all their structures and conditions being stable under base change). By (a) it is a generic object over the
$R'$-scheme $M_{R'}$, so the universal property of $X$ gives a unique $f\colon M_{R'}\to X$ over $R'$
with a morphism of objects $o_S\to\mathcal A_X$ over $f$. By (b), $\mathcal A_X$ is an object for $b$
over $X$, so the universal property of $M$ gives a unique $g\colon X\to M$ over $R$ with a morphism
$\mathcal A_X\to\mathcal A_M$ over $g$. Let $j=(g,s_X)\colon X\to M_{R'}$. Then
$\mathrm{pr}_2\circ j=s_X$ and $\mathrm{pr}_1\circ j=g$, so $\mathcal A_X$ is the pullback of
$\mathcal A_M$ along $\mathrm{pr}_1\circ j$.

We show that $f$ is inverse to $j$. First, composing $o_S\to\mathcal A_X$ over $f$ with $\mathcal
A_X\to\mathcal A_M$ over $g$ gives a morphism $o_S\to\mathcal A_M$ over $g\circ f$, and the defining
morphism $o_S\to\mathcal A_M$ lies over $\mathrm{pr}_1$; by the uniqueness in the universal property
of $M$, $g\circ f=\mathrm{pr}_1$. Second, since $o_S$ is the pullback of $\mathcal A_M$ along
$\mathrm{pr}_1$ and $\mathrm{pr}_1\circ j=g$, the morphism $\mathcal A_X\to\mathcal A_M$ over $g$
factors through a morphism $\mathcal A_X\to o_S$ over $j$; composed with $o_S\to\mathcal A_X$ it
gives a morphism $\mathcal A_X\to\mathcal A_X$ over $f\circ j$, and the identity is one over
$\mathrm{id}_X$; by the uniqueness in the universal property of $X$, $f\circ j=\mathrm{id}_X$. Third,
$\mathrm{pr}_1\circ j\circ f=g\circ f=\mathrm{pr}_1$ and $\mathrm{pr}_2\circ j\circ f=s_X\circ
f=\mathrm{pr}_2$ ($f$ lies over $R'$), so $j\circ f=\mathrm{id}$. Hence $j$ is an isomorphism, in
particular an open and a closed immersion.
\end{proof}

\begin{lemma}\label{pf:ks2:gensm}
Let $R$ be a field of characteristic $0$, $g_b$ a characteristic-zero base of $D$ over $R$, $N$ a
level and $M_N$ a generic model of level $N$. Then $M_N\to\Spec R$ is smooth, separated and
quasi-compact, and projective if $D$ is of compact type. For every $k$, the fibre power
$\kstwoKS_N^k\to\Spec R$ of the universal abelian scheme is smooth and separated, and projective
if $D$ is of compact type.
\formallabel{UpperTheory.Gal.G3.pGenSm}\formallabel{UpperTheory.Gal.G3.pPowSm}
\end{lemma}

\begin{proof}
Since $R$ is a $\Q$-algebra, $g_b$ is also a base of $D$ over $R$ in the sense of
\cref{pf:ks2:def:base}(1) (integers prime to $p$ are units, and there is no ring map from $R$ to
a field of characteristic $p$). By \lref{GlobalTheory.UPel.HT.model_nonempty} (models of
level $N$ exist over Noetherian rings with a base) there is a model $M$ of level $N$ over $R$; by
\cref{pf:ks2:genmodel} (applied with $\varphi=\mathrm{id}_R$) there is a morphism
$j\colon M_N\to M$ over $R$ which is an open and a closed immersion. As $M$ is quasi-projective and
smooth over $R$, $M_N$ is smooth, separated and quasi-compact; if $D$ is of compact type, $M$ is
projective over $R$ by \cref{pf:ks2:proj}, hence so is the closed subscheme $M_N$. The abelian
scheme $\kstwoKS_N\to M_N$ is smooth and proper, so its $k$-fold fibre power is smooth and proper
over $M_N$ (\cref{pf:ks2:powsmooth}, and stability of properness under base change and
composition), hence smooth and separated over $R$. In the compact case the universal object carries
a relatively ample line bundle (the restriction of the polarization bundle to the diagonal), and
\cref{pf:ks2:powproj}, applied to the projective morphism $M_N\to\Spec R$, gives projectivity of
$\kstwoKS_N^k$ over $R$.
\end{proof}

\subsubsection{Existence of Kuga--Sato towers}

The construction of the towers rests on the following properties of Hecke witnesses
(\cref{pf:ks2:def:wit}). Throughout, $R$ is a field of characteristic $0$, $g_b$ a
characteristic-zero base over $R$, and for every level $N$ a generic model $M_N$ is fixed.

\begin{lemma}[Principal congruence subgroups]\label{pf:ks2:kn}
Let $N\ge1$. (1) Every $k\in K_1$ normalizes $K_N$. (2) For every $h\in\kstwoG$ there is an integer
$c$ prime to $p$ with $hK_{Nc}h^{-1}\subseteq K_N$. (3) For $M,N>0$ the index $[K_N:K_N\cap K_M]$ is
finite.
\formallabel{UpperTheory.Gal.G3.pKNBasic}
\end{lemma}

\begin{proof}
Recall that $u\in K_N$ iff $(u-1)\kstwoLaHat\subseteq N\kstwoLaHat$ and $(u^{-1}-1)\kstwoLaHat\subseteq
N\kstwoLaHat$. (1) For $k\in K_1$, $k$ and $k^{-1}$ preserve $\kstwoLaHat$, so
$(kuk^{-1}-1)\kstwoLaHat=k(u-1)k^{-1}\kstwoLaHat\subseteq k(u-1)\kstwoLaHat\subseteq kN\kstwoLaHat=
N\kstwoLaHat$, and similarly for $(kuk^{-1})^{-1}=ku^{-1}k^{-1}$. (2) Since $\kstwoLaHat$ is open and
compact in $\A_f^p\otimes\Lambda$, there are integers $j,j'\ge1$ prime to $p$ with
$jh\kstwoLaHat\subseteq\kstwoLaHat$ and $j'h^{-1}\kstwoLaHat\subseteq\kstwoLaHat$. Put $c=jj'$. For $u\in
K_{Nc}$: $(huh^{-1}-1)\kstwoLaHat=h(u-1)h^{-1}\kstwoLaHat\subseteq h(u-1)j'^{-1}\kstwoLaHat\subseteq
hNj\kstwoLaHat\subseteq N\kstwoLaHat$, and the same for $u^{-1}$. (3) $K_{NM}\subseteq K_N\cap K_M$, and
$K_{NM}$ has finite index in $K_1\supseteq K_N$, being the kernel of the action of $K_1$ on the finite
set $\kstwoLaHat/NM\kstwoLaHat$.
\end{proof}

\begin{lemma}[Hecke witnesses]\label{pf:ks2:wit}
Let $N',N$ be levels and $h\in\kstwoG$.
\begin{enumerate}
\item\label{pf:ks2:wit:kn} (\cref{pf:ks2:kn}) Every $k\in K_1$ normalizes
  $K_N$; for every $h$ there is $c$ prime to $p$ with $hK_{Nc}h^{-1}\subseteq K_N$; and all
  relative indices between the $K_N$ are finite.
\item\label{pf:ks2:wit:ex} (\lref{UpperTheory.Gal.G3.witEx_holds}) If $h$ is integral and
  admissible from $N'$ to $N$, a Hecke witness $(\pi,\Phi)$ for $h$ from $M_{N'}$ to $M_N$ exists.
\item\label{pf:ks2:wit:uniq} (\lref{UpperTheory.Gal.G3.witUniq_holds}) Under the same
  hypotheses it is unique.
\item\label{pf:ks2:wit:comp} (\lref{UpperTheory.Gal.G3.witComp_holds}) For $u\in K_{N'}$,
  $(\mathrm{id},\mathrm{id})$ is a witness for $u$ from $M_{N'}$ to itself; and if $(\pi,\Phi)$ is a
  witness for $h$ from $M_{N_1}$ to $M_{N_2}$ and $(\pi',\Phi')$ one for an integral $h'$,
  admissible from $N_2$ to $N_3$, from $M_{N_2}$ to $M_{N_3}$, then $(\pi'\pi,\Phi'\Phi)$ is a
  witness for $h'h$.
\item\label{pf:ks2:wit:cart} (\lref{UpperTheory.Gal.G3.witCart_holds}) If $u\in K_1$, the
  square formed by $\Phi$, $\pi$ and the two structure maps of the universal abelian schemes is
  cartesian.
\item\label{pf:ks2:wit:one} (\lref{UpperTheory.Gal.G3.witOne_holds}) If $1$ is admissible from
  $N'$ to $N$ (i.e.\ $K_{N'}\subseteq K_N$), the map $\pi$ of a witness for $1$ is finite, \'etale
  and surjective.
\item\label{pf:ks2:wit:trans} (\lref{UpperTheory.Gal.G3.witTrans_holds}) If $(\pi,\Phi)$ is a
  witness for $1$ and two geometric points $x,y$ of $M_{N'}$ have $\pi x=\pi y$, then there is
  $u\in K_N$ such that the witness $(\pi_u,\Phi_u)$ for $u$ from $M_{N'}$ to itself satisfies
  $\pi_u x=y$ and $\Phi\Phi_u=\Phi$ on the fibre over $x$.
\item\label{pf:ks2:wit:free} (\lref{UpperTheory.Gal.G3.witFree_holds}) If $K_{N'}\subseteq
  K_N$, $u\in K_N$, and the map $\pi$ of the witness for $u$ from $M_{N'}$ to itself fixes a
  geometric point, then $u\in K_{N'}$.
\end{enumerate}
\end{lemma}

\begin{proof}
Part (1) is \cref{pf:ks2:kn}; the other parts are the cited formal results, whose arguments we
indicate. (2): by \lref{UpperTheory.Gal.G3.witDeep_holds} there are a level $M$ divisible by
$N'$, with $h$ admissible from $M$ to $N$, and a witness for $h$ from $M_M$ to $M_N$ (it rests on
\lref{UpperTheory.Gal.G3.transDeep_holds}: for integral $h$ there is an integer $c$ prime to
$p$ such that every generic object of level $N_1c$ over a locally Noetherian base admits a Hecke
isogeny for $h$ to a generic object of level $N_1$); the level-forgetting witness $M_M\to M_{N'}$ is
finite, \'etale and surjective (\lref{UpperTheory.Gal.G3.witForget_holds}); the morphism of
moduli schemes descends along it (\lref{UpperTheory.Gal.G3.descMap_holds}), and the Hecke
isogeny property descends along finite \'etale surjective base changes and can be cancelled against
a Hecke isogeny for $1$ (\lref{UpperTheory.Gal.G3.transDesc_holds},
\lref{UpperTheory.Gal.G3.transCancel_holds}). (3): $M_N$ is separated and $M_{N'}$ is reduced
(\cref{pf:ks2:gensm}), so two morphisms $M_{N'}\to M_N$ agree if they agree on geometric points; at a
geometric point both witnesses give Hecke isogenies for $h$ out of the same object, which coincide;
the same argument applies to $\Phi$ on the reduced scheme $\kstwoKS_{N'}$. (4) composes Hecke
isogenies. (5): a Hecke isogeny for $u\in K_1$ is an isomorphism. (6): with $M=NN'$, the composite
of the level-forgetting witness $M_M\to M_{N'}$ (finite, \'etale and surjective) with $\pi$ is, by (3)
and (4), the level-forgetting witness $M_M\to M_N$, which is finite, \'etale and surjective; hence so
is $\pi$, as $\pi$ is separated and locally of finite type. (7) and (8) are statements about
geometric points: two level-$N'$ objects with the same image at level $N$ differ by an element of
$K_N$ acting on the level structure, and a point fixed by the action of $u$ carries an automorphism
inducing $u$ on the level-$N'$ structure, which forces $u\in K_{N'}$.\end{proof}

\begin{proposition}\label{pf:ks2:ksex}
$\sfname{KSEx}$ holds: for $\mathrm{Sc}(D,G)$, every field $R$ of characteristic $0$, every
characteristic-zero base $g_b$ over $R$ and every $k\in\N$, a Kuga--Sato tower for
$(L,g_b,k)$ exists.
\formallabel{UpperTheory.Gal.G3.kSEx_holds}
\end{proposition}

\begin{proof}
For every level $N$ choose a generic model $M_N$ of level $N$. It exists: $g_b$ is a base over $R$, so a model $M$ of
level $N$ exists by \lref{GlobalTheory.UPel.HT.model_nonempty}, and $M$ with its universal
object, viewed as generic, is a generic model.
For integral
$h$ admissible from $N'$ to $N$ let $(\pi_{N',N,h},\Phi_{N',N,h})$ be the Hecke witness of
\cref{pf:ks2:wit}\ref{pf:ks2:wit:ex}. Define the tower $T$: the
index set is the set of levels, $K_N$ the principal congruence subgroups, $\Delta$ the integral
elements, $Z_N=\kstwoKS_N^k$ with its structure map to $\Spec R$, and
$\mathrm{tr}(N',N,h)=\Phi_{N',N,h}^{k}$, the $k$-th fibre power of $\Phi_{N',N,h}$ over
$\pi_{N',N,h}$. We verify the axioms of \cref{pf:ks2:def:tower}.

$K_N\subseteq\Delta$ since elements of $K_N$ preserve $\kstwoLaHat$. Axiom (i): given $j=N$ and
finitely many pairs $(N_i,x_i)$, choose by \cref{pf:ks2:wit}\ref{pf:ks2:wit:kn} integers $c_i$
prime to $p$ with $x_iK_{N_ic_i}x_i^{-1}\subseteq K_{N_i}$ and put $m=N\prod_iN_ic_i$; then
$K_m\subseteq K_N$, $K_m$ is normalized by $K_1\supseteq K_N$, and $x_iK_mx_i^{-1}\subseteq
K_{N_i}$ because $K_m\subseteq K_{N_ic_i}$.
Axiom (ii) is the finiteness of relative indices. The maps $\mathrm{tr}$ lie over $\Spec R$
because $\pi$ does. Axiom (iii): for $u\in K_N$, both $(\pi_{N,N,u},\Phi_{N,N,u})$ and
$(\mathrm{id},\mathrm{id})$ are witnesses for $u$, so they coincide by uniqueness, and the fibre
power of the identity is the identity; for composable $h,h'$, the composite of the witnesses for
$h$ and $h'$ is a witness for $h'h$ (\cref{pf:ks2:wit}\ref{pf:ks2:wit:comp}), hence equals the
chosen one, and fibre powers are compatible with composition. Axiom (iv): by
\cref{pf:ks2:wit}\ref{pf:ks2:wit:cart} the square formed by $\Phi_{N',N,1}$ and $\pi_{N',N,1}$ is
cartesian, hence so is the square formed by its $k$-th fibre power; so $\mathrm{tr}(N',N,1)$ is a base change of
$\pi_{N',N,1}$, which is finite, \'etale and surjective by \cref{pf:ks2:wit}\ref{pf:ks2:wit:one},
and these properties are stable under base change. Axiom (v): if two geometric points $x,y$ of
$\kstwoKS_{N'}^k$ have the same image, their images $\bar x,\bar y$ in $M_{N'}$ have the same image
in $M_N$; \cref{pf:ks2:wit}\ref{pf:ks2:wit:trans} gives $u\in K_N$ with $\pi_{N',N',u}\bar x=\bar y$
and $\Phi_{N',N,1}\circ\Phi_{N',N',u}=\Phi_{N',N,1}$ on the fibre over $\bar x$; as the square of
$\Phi^k_{N',N,1}$ is cartesian, $y$ and $\mathrm{tr}(N',N',u)(x)$ are two points over $\bar y$
with the same image, hence equal. Axiom (vi) follows from
\cref{pf:ks2:wit}\ref{pf:ks2:wit:free} applied to the image of the fixed point in $M_{N'}$.
Finally the structure of a Kuga--Sato tower is given by the identity isomorphisms
$Z_N=\kstwoKS_N^k$, and the compatibility of $\mathrm{tr}(N',N,h)$ with a Hecke witness holds by
construction.
\end{proof}

\subsubsection{Niceness}

\begin{lemma}\label{pf:ks2:twfet}
Let $X$ be a tower over $B$ for $\Gamma$ (\cref{pf:ks2:def:tower}) and $K_i\subseteq K_j$. Then
$\mathrm{tr}(i,j,1)\colon Z_i\to Z_j$ is finite \'etale of degree $[K_j:K_i]$.
\formallabel{UpperTheory.Gal.G3.pTwFEt}
\end{lemma}

\begin{proof}
Finiteness and \'etaleness are axiom (iv); we count geometric fibres. First assume that $K_i$ is
normalized by $K_j$. Let $x$ be a geometric point of $Z_j$; by surjectivity it has a lift $y$.
By axiom (v) every lift is $\mathrm{tr}(i,i,k)(y)$ for some $k\in K_j$, and by axiom (iii) these
are lifts; $\mathrm{tr}(i,i,k)(y)=\mathrm{tr}(i,i,k')(y)$ iff $\mathrm{tr}(i,i,k'^{-1}k)$ fixes
$y$ iff $k'^{-1}k\in K_i$ by axiom (vi). Hence the fibre is in bijection with $K_j/K_i$. In
general, axiom (i) gives $m$ with $K_m\subseteq K_i$
and $K_m$ normalized by $K_j$, hence also by $K_i$. Since
$\mathrm{tr}(m,j,1)=\mathrm{tr}(i,j,1)\circ\mathrm{tr}(m,i,1)$, the fibre of $\mathrm{tr}(m,j,1)$
over $x$ is the disjoint union, over the lifts $y$ of $x$ to $Z_i$, of the fibres of
$\mathrm{tr}(m,i,1)$ over $y$. By the normal case,
$[K_j:K_m]=\#\{y\}\cdot[K_i:K_m]$, and $[K_j:K_m]=[K_j:K_i][K_i:K_m]$ with $[K_i:K_m]\neq0$; so
the number of lifts is $[K_j:K_i]$.
\end{proof}

\begin{definition}[Splittings and Galois covers]\label{pf:ks2:def:splits}
Let $s\colon Y\to Z$ and $c\colon W\to Z$ be morphisms of schemes and $(\sigma_x)_{x\in J}$ a family
of morphisms $W\to Y$ with $s\circ\sigma_x=c$. The family \emph{splits $s$ over $c$} if the morphism
$\coprod_{x\in J}W\to Y\times_ZW$ which is $(\sigma_x,\mathrm{id}_W)$ on the copy indexed by $x$ is an
isomorphism. For a group $\Gamma$ and a map $a\colon\Gamma\to\End(W)$, a morphism $c\colon W\to X$ is
\emph{Galois with group $\Gamma$ acting by $a$} if $c$ is finite, \'etale and surjective,
$a(1)=\mathrm{id}_W$, $a(gh)=a(h)\circ a(g)$ for $g,h\in\Gamma$, and the family $(a(g))_{g\in\Gamma}$
splits $c$ over $c$.
\formal{GlobalTheory.GalV.Splits}\formal{GlobalTheory.GalV.IsGalois}
\end{definition}

The \emph{splitting criterion} is the following statement: if $s\colon Y\to Z$ is finite and \'etale,
$c\colon W\to Z$ is a morphism and $(\sigma_x)_{x\in J}$ is a finite family of morphisms $W\to Y$ with
$s\circ\sigma_x=c$ such that, for every algebraically closed field $\Omega$ and all $w\in W(\Omega)$,
$y\in Y(\Omega)$ with $s(y)=c(w)$, there is exactly one $x\in J$ with $\sigma_x(w)=y$, then
$(\sigma_x)$ splits $s$ over $c$. It is proved in the formal development
(\lref{UpperTheory.Gal.G3.splitsPts_holds}: the morphism in question is a morphism of finite
\'etale $W$-schemes which is bijective on geometric points, hence an isomorphism).

\begin{lemma}[Splitting of tower maps]\label{pf:ks2:twsplit}
Assume the splitting criterion. Let $X$ be a tower over $B$ for $\Gamma$, and let $m,i,j$ be indices
with $K_m\subseteq K_i\subseteq K_j$ and $K_m$ normalized by $K_j$. Let $J$ be a finite set and
$r\colon J\to K_j$ a map with $r_xK_mr_x^{-1}\subseteq K_i$ for all $x$, such that for every $k\in
K_j$ there is exactly one $x\in J$ with $kr_x^{-1}\in K_i$. Then the family
$(\mathrm{tr}(m,i,r_x))_{x\in J}$ splits $\mathrm{tr}(i,j,1)$ over $\mathrm{tr}(m,j,1)$.
\formallabel{UpperTheory.Gal.G3.pTwSplit}
\end{lemma}

\begin{proof}
All the transition maps used below are defined, because $K_m$ is normalized by $K_j\supseteq K_i$.
By axiom (iii) of \cref{pf:ks2:def:tower} and $r_x\in K_j$,
$\mathrm{tr}(i,j,1)\circ\mathrm{tr}(m,i,r_x)=\mathrm{tr}(m,j,r_x)=\mathrm{tr}(j,j,r_x)\circ
\mathrm{tr}(m,j,1)=\mathrm{tr}(m,j,1)$. The map $\mathrm{tr}(i,j,1)$ is finite and \'etale by axiom
(iv). We verify the hypothesis of the splitting criterion. Let $w$ and $y$ be geometric points of
$Z_m$ and $Z_i$ (over the same algebraically closed field) with
$\mathrm{tr}(i,j,1)(y)=\mathrm{tr}(m,j,1)(w)$.

\emph{Existence.} The map $\mathrm{tr}(m,i,1)$ is finite and surjective, so
$y=\mathrm{tr}(m,i,1)(z)$ for a geometric point $z$ of $Z_m$. Then
$\mathrm{tr}(m,j,1)(z)=\mathrm{tr}(i,j,1)(y)=\mathrm{tr}(m,j,1)(w)$, and by axiom (v) (for $K_m$,
which is normalized by $K_j$) $z=\mathrm{tr}(m,m,k)(w)$ for some $k\in K_j$. Let $x_0$ be the index
with $kr_{x_0}^{-1}\in K_i$. Since $\mathrm{tr}(i,i,u)$ is the identity for $u\in K_i$,
\begin{align*}
y&=\mathrm{tr}(m,i,1)\,\mathrm{tr}(m,m,k)(w)=\mathrm{tr}(m,i,k)(w)\\
&=\mathrm{tr}(i,i,kr_{x_0}^{-1})\,\mathrm{tr}(m,i,r_{x_0})(w)=\mathrm{tr}(m,i,r_{x_0})(w).
\end{align*}

\emph{Uniqueness.} Suppose that $\mathrm{tr}(m,i,r_{x_1})(w)=\mathrm{tr}(m,i,r_{x_2})(w)$. Then the
geometric points $\mathrm{tr}(m,m,r_{x_1})(w)$ and $\mathrm{tr}(m,m,r_{x_2})(w)$ of $Z_m$ have the same
image under $\mathrm{tr}(m,i,1)$, and $K_m$ is normalized by $K_i$; by axiom (v) there is $k\in K_i$
with $\mathrm{tr}(m,m,r_{x_2})(w)=\mathrm{tr}(m,m,kr_{x_1})(w)$. Applying $\mathrm{tr}(m,m,r_{x_2}^{-1})$
shows that $\mathrm{tr}(m,m,u)$ fixes $w$ for $u=r_{x_2}^{-1}kr_{x_1}\in K_j$, so $u\in K_m$ by axiom
(vi). The element $k'=kr_{x_1}=r_{x_2}u$ of $K_j$ satisfies $k'r_{x_1}^{-1}=k\in K_i$ and
$k'r_{x_2}^{-1}=r_{x_2}ur_{x_2}^{-1}\in r_{x_2}K_mr_{x_2}^{-1}\subseteq K_i$. By the uniqueness in the
hypothesis on $r$, $x_1=x_2$.

The splitting criterion now shows that $(\mathrm{tr}(m,i,r_x))_x$ splits $\mathrm{tr}(i,j,1)$ over
$\mathrm{tr}(m,j,1)$.
\end{proof}

\begin{lemma}[Galois property of tower maps]\label{pf:ks2:twgal}
Assume that the conclusion of \cref{pf:ks2:twsplit} holds for all towers (by that lemma, this
follows from the splitting criterion). Let $X$ be a tower and $m,j$ indices with $K_m\subseteq K_j$
and $K_m$ normalized by $K_j$. Then there are a finite group $\Gamma$ and a map $r\colon\Gamma\to K_j$
with $r_cK_mr_c^{-1}\subseteq K_m$ for all $c$, such that for every $k\in K_j$ there is exactly one
$c\in\Gamma$ with $kr_c^{-1}\in K_m$, and $\mathrm{tr}(m,j,1)$ is Galois with group $\Gamma$ acting by
$c\mapsto\mathrm{tr}(m,m,r_c)$.
\formallabel{UpperTheory.Gal.G3.pTwGal}
\end{lemma}

\begin{proof}
$K_m$ is a normal subgroup of $K_j$; let $\Gamma=K_j/K_m$, which is finite by axiom (ii). Choose a
representative $\bar c\in K_j$ of every $c\in\Gamma$ and put $r_c=\bar c^{-1}$; then
$r_cK_mr_c^{-1}=K_m$. For $k\in K_j$, $kr_c^{-1}=k\bar c$ lies in $K_m$ if and only if $c$ is the class
of $k^{-1}$, so $c$ exists and is unique. Put $a(c)=\mathrm{tr}(m,m,r_c)$. The map $\mathrm{tr}(m,j,1)$
is finite, \'etale and surjective by axiom (iv). Since $\mathrm{tr}(m,m,u)=\mathrm{id}$ for $u\in K_m$,
axiom (iii) shows that $\mathrm{tr}(m,m,h)$ depends only on the coset $K_mh$ for $h\in K_j$. Hence
$a(1)=\mathrm{id}$, as $\bar1\in K_m$; and for $c,d\in\Gamma$ we have $\overline{cd}\in\bar c\bar dK_m$,
so $r_{cd}\in K_mr_dr_c$ and $a(cd)=\mathrm{tr}(m,m,r_dr_c)=a(d)\circ a(c)$. Finally, the conclusion
of \cref{pf:ks2:twsplit} for $i=m$, $J=\Gamma$ and this $r$ (whose hypotheses were just verified)
states that $(a(c))_{c\in\Gamma}$ splits $\mathrm{tr}(m,j,1)$ over itself. So $\mathrm{tr}(m,j,1)$ is
Galois with group $\Gamma$ acting by $a$.
\end{proof}

\begin{lemma}\label{pf:ks2:sigmafet}
If $s\colon Y\to Z$ is finite \'etale of degree $d$ and $I$ is a finite set, the map
$\coprod_{i\in I}Y\to Z$ equal to $s$ on every copy is finite \'etale of degree $\#I\cdot d$.
\formallabel{UpperTheory.Gal.G3.pSigmaFEt}
\end{lemma}

\begin{proof}
The map factors as the codiagonal $\coprod_IY\to Y$, which is finite, followed by $s$; so it is
finite. \'Etaleness is local on the source and holds on each copy. A geometric point of $Z$ has $d$
lifts to each copy, hence $\#I\cdot d$ lifts in total.
\end{proof}

\begin{proposition}\label{pf:ks2:ksnice}
$\sfname{KSNice}$ holds.
\formallabel{UpperTheory.Gal.G3.kSNice_holds}
\end{proposition}

\begin{proof}
Let $T$ be a Kuga--Sato tower over a characteristic-zero field $R$ under $\mathrm{Sc}(D,G)$, $N$ a
level and $(g_j)$ Hecke data of level $N$. The structure map
$Z_N\to\Spec R$ is identified with $\kstwoKS_N^k\to\Spec R$, which is smooth, and projective because
$D$ is of compact type (part of $\mathrm{Sc}$), by \cref{pf:ks2:gensm}. For $j\in J$, the map
$s_j\colon\coprod_{c}Z_{M_{g_j}}\to Z_N$ is $\mathrm{tr}(M_{g_j},N,1)$ on each of the
$r_j=[K_{g_j}:K_N]$ copies (the index set $K_{g_j}/(K_N\cap K_{g_j})$ has $r_j$ elements, finite
because $[K_1:K_N]=[K_1:K_{g_j}][K_{g_j}:K_N]$ is finite). Since $K_{M_{g_j}}\subseteq K_N$
(as $N\mid M_{g_j}$), \cref{pf:ks2:twfet} shows that $\mathrm{tr}(M_{g_j},N,1)$ is finite \'etale of
degree $d_j=[K_N:K_{M_{g_j}}]$. By \cref{pf:ks2:sigmafet}, $s_j$ is
finite \'etale of degree $r_jd_j$.
\end{proof}

% ---------------------------------------------------------------------------
\subsection{The integral input}\label{pf:ks2:sec:min}

\subsubsection{\'Etale-level models}

Let $D$ be a datum for $p$, $\mathfrak a\subseteq\kstwoOF$ a nonzero ideal and $r\in\N$. Put
$\Lambda_{\mathfrak a,r}=\Lambda/\mathfrak a^r\Lambda$ and
let $\Aut_O(\Lambda_{\mathfrak a,r})$ be its group of $O$-linear automorphisms. For an object $o$ over $S$, an \emph{\'etale level
structure of type $(\mathfrak a,r)$} is an $O$-linear
homomorphism $\lambda\colon\Lambda\to A_o(S)$ vanishing on $\mathfrak a^r\Lambda$ such that at every
geometric point the induced map $\Lambda_{\mathfrak a,r}\to A_o[\mathfrak a^r]$ is injective and
surjective onto the points killed by $\mathfrak a^r$. An \emph{\'etale-level model} of level
$(N,\mathfrak a,r)$ is a scheme $Y\to\Spec R$ with an
object of level $N$ and an \'etale level structure of type $(\mathfrak a,r)$ which represent
pairs (object, \'etale level structure). A maximal ideal $\mathfrak w_0$ is an \emph{\'etale place}
of $b$ if $p\in\mathfrak w_0$, each
$t_i(\mathfrak w_0)$ generates the unit ideal of $R$, and for all ring maps $\beta\colon R\to k$ and
$\varphi\colon\kstwoOF\to k$ to a field of characteristic $p$ with $\ker\varphi=\mathfrak w_0$ one has
$\varphi(\omega)=\beta(\rho)$; the ideal $\mathfrak a$ is \emph{\'etale} for $b$ if it is nonzero and every maximal ideal
containing it is an \'etale place of $b$.

\begin{theorem}[The integral input]\label{pf:ks2:min}
The proposition $\sfname{MIn}$ holds, i.e.\ the
following seven statements hold for every prime $p$ and datum $D$ for $p$.
\begin{enumerate}
\item (Existence of models.) For every Noetherian ring $R$ with a base $b$ and every level $N$,
  a model of level $N$ exists.
\item (Projectivity.) If moreover $R$ is a domain of characteristic $0$ and $D$ is of
  compact type, every model of level $N$ is projective over $R$.
\item (Generic models.) If $R$ is Noetherian with a base $b$, $\varphi\colon R\to R'$ is a ring
  map to a $\Q$-algebra and $g_b$ is a characteristic-zero base over $R'$ with $\rho=\varphi(\rho_b)$
  and $t_i=\varphi\circ t_{b,i}$, then for every model $M$ of level $N$ over $R$ and every generic
  model $X$ of level $N$ over $R'$ there is a morphism $X\to M\times_RR'$ over $R'$ which is an
  open and a closed immersion and under which the universal generic object is the pullback of the
  universal object of $M$.
\item (\'Etale-level models.) For a Noetherian ring $R$ with a base $b$, a level $N$, an \'etale ideal
  $\mathfrak a$, $r\in\N$ and a model $M$ of level $N$, there are an \'etale-level model $Y$ of
  level $(N,\mathfrak a,r)$, a morphism $\pi\colon Y\to M$ under which the universal object of $Y$
  is the pullback of that of $M$, and an action of $\Aut_O(\Lambda_{\mathfrak a,r})$ on $Y$ making
  $\pi$ a finite \'etale surjective Galois cover with group $\Aut_O(\Lambda_{\mathfrak a,r})$
  (the group acts over $M$, and two geometric points with
  the same image differ by a unique group element).
\item (\'Etale ideals.) Let $O$ be a discrete valuation ring of characteristic $0$ with $p$ in its
  maximal ideal, $b$ a base over $O$ with $a_D=1$, $r\in\Z$ and $\mathfrak w\supseteq(p,\omega-r)$ a
  maximal ideal with $t_i(\mathfrak w)$ contained in the maximal ideal of $O$. Then an ideal
  $\mathfrak a$ with $\mathfrak a+\mathfrak w=\kstwoOF$ is \'etale for $b$ iff it contains a power
  of $(p,\omega-r)$.
\item (Congruence relation.) The congruence relation of \cref{pf:ks2:congks}.
\item (Degrees of $T_j$-correspondences.) The $T_j$-correspondences of \cref{pf:ks2:def:tj} are, after any base
  change along $Y\to M$, finite, \'etale after inverting $p$, and have, over every geometric point
  where $p$ is invertible, exactly as many geometric points as there are $j$-dimensional subspaces
  of $(\kstwoOF/\mathfrak w)^n$.
\end{enumerate}
\formallabel{UpperTheory.Gal.Int.Geo.CovMain.mIn}
\end{theorem}

\begin{proof}
(1) is \lref{GlobalTheory.UPel.HT.model_nonempty}, which states that models of every level
$N\ge3$ prime to $p$ exist over every Noetherian ring with a base. (2) is \cref{pf:ks2:proj}. (3) is
\cref{pf:ks2:genmodel}. (4) is
\cref{pf:ks2:finmodel}, whose hypothesis is
\lref{UpperTheory.Gal.FiniteEtaleModels.finCover_holds}. (5) is
\lref{UpperTheory.Gal.FiniteEtaleModels.placeOf_holds}. (6) is
\cref{pf:ks2:congks}. (7) is \lref{UpperTheory.Gal.crystallineComparison.Mod.corrBC_holds},
applied to \cref{pf:ks2:quotobj} (existence of
quotients of objects by finite \'etale $O$-stable subgroup schemes killed by $\mathfrak w$ as
$T_j$-isogenies); it computes the number of geometric points in the fibres through the bijection
between sections of the $T_j$-correspondence and $j$-dimensional $\kstwoOF/\mathfrak w$-subspaces
of the $\mathfrak w$-torsion, which is free of rank $n$ in characteristic $0$.
\end{proof}

\begin{proposition}\label{pf:ks2:proj}
Let $R$ be a Noetherian domain of characteristic $0$ with a base $b$, $N$ a level, $D$ of compact
type and $M$ a model of level $N$. Then $M$ is projective over $R$.
\formallabel{GlobalTheory.UPel.HT.Model.isProjOver_of_compactB}
\end{proposition}

\begin{proof}
By \lref{GlobalTheory.T.harrisTaylor_integralModel_isProper_of_compact}, $M\to\Spec R$ is
proper (this is the classical properness of compact PEL moduli spaces, see
\cref{pf:ks2:classical:proper}). A quasi-projective morphism which is proper is projective: the
immersion into $\mathbb P^N_R$ is then a closed immersion, since a proper morphism to a separated
scheme has closed image and a proper immersion is closed.
\end{proof}

\begin{proposition}\label{pf:ks2:finmodel}
Assume the following relative statement: for every Noetherian ring $R$ with a base, every level $N$,
every \'etale ideal $\mathfrak a$, every $r$, and every object $o$ of level $N$ over a locally
Noetherian $R$-scheme $S$, there are a finite \'etale Galois cover $\pi\colon Y\to S$ with group
$\Aut_O(\Lambda_{\mathfrak a,r})$, an object $o'$ over $Y$ with a morphism $o'\to o$ over $\pi$, and an
\'etale level structure $\lambda$ on $o'$ of type $(\mathfrak a,r)$, such that for every $t\colon
T\to S$, every object $o_1$ over $T$ with a morphism $o_1\to o$ over $t$, and every \'etale level
structure on $o_1$, there is a unique $f\colon T\to Y$ over $S$ carrying $(o',\lambda)$ to $o_1$ with
its level structure. Then statement (4) of \cref{pf:ks2:min} holds: \'etale-level models exist.
\formallabel{UpperTheory.Gal.FiniteEtaleModels.finModel_holds}
\end{proposition}

\begin{proof}
Let $M$ be a model of level $N$ over $R$ and $\mathfrak a$ \'etale for $b$. $M$ is smooth of finite
type over $R$, hence locally Noetherian. Apply the hypothesis to the universal object
$\mathcal A_M$ over $S=M$, obtaining $\pi\colon Y\to M$, $o'$ and $\lambda$. Then $Y$ represents
pairs: given an object $o_1$ over an
$R$-scheme $T$ with an \'etale level structure, the universal property of $M$ gives a unique
$c\colon T\to M$ with a morphism $o_1\to\mathcal A_M$ over $c$, and the relative universal property
gives a unique $f\colon T\to Y$ over $c$; uniqueness of $f$ over $R$ follows because any $f'$
compatible with the structures has $\pi f'$ compatible with the object (morphisms of objects
compose), hence $\pi f'=c$ by the uniqueness in $M$, and then $f'=f$. So $(Y,o',\lambda)$ is an
\'etale-level model and $\pi$ has the required properties.
\end{proof}

The hypothesis of \cref{pf:ks2:finmodel} is proved in the development: it is
\lref{UpperTheory.Gal.FiniteEtaleModels.finCover_holds}.

% ---------------------------------------------------------------------------
\subsection{The congruence relation on Kuga--Sato varieties}\label{pf:ks2:sec:cong}

\subsubsection{Statement}

\begin{definition}[$T_j$-isogenies and $T_j$-correspondences]\label{pf:ks2:def:tj}
Let $b$ be a base over $R$, $r\in\Z$, $\mathfrak w\supseteq\mathfrak u=(p,\omega-r)$ a maximal ideal of
$\kstwoOF$, $0\le j\le n$ and $c$ a power of $p$.
\begin{enumerate}
\item A homomorphism $\psi\colon A_o\to A_{o'}$ between objects of the same level over $S$ is a
  \emph{$p$-isogeny of degree $c$} if it is $O$-linear,
  compatible with the level structures, pulls the polarization bundle of $o'$ back to the $c$-th
  power of that of $o$, and has a quasi-inverse $\psi'$ with $\psi'\psi=c$ and $\psi\psi'=c$. It is
  \emph{of type $T_j$ at $\mathfrak w$} if at every geometric
  point of characteristic $0$ every point of $\ker\psi$ killed by a power of $\mathfrak u$ is killed
  by $\mathfrak w$, and the number of points of $\ker\psi$ killed by $\mathfrak w$ is
  $\Nm\mathfrak w^{\,mj}$.
\item For a model $M$ of level $N$ over $R$, a \emph{$T_j$-correspondence of degree $c$} is $M\xleftarrow{s}Y_j\xrightarrow{t}M$
  with $s$ finite and \'etale over $M[1/p]$, together with objects $o,o'$ over $Y_j$ mapping to the
  universal object over $s$ and over $t$, and a $p$-isogeny $\psi\colon A_o\to A_{o'}$ of degree $c$
  of type $T_j$ at $\mathfrak w$, which is universal among such isogenies over schemes where $p$ is
  invertible.
\item For an \'etale-level model $Y\to M$ (with a morphism of universal objects over $\pi$) and a
  $T_j$-correspondence $T$ on $M$, a \emph{lift of $T$ to $Y$} is a correspondence
  $Y\xleftarrow{s}Z\xrightarrow{t}Y$ with $s$ the base change of $T.s$ along $\pi$, $\pi t$
  factoring through $T.t$, and a $T_j$-isogeny over $Z$ compatible with that of $T$ and carrying
  the \'etale level structure along. Its $k$-th Kuga--Sato power is the correspondence
  $\kstwoKS_Y^k\leftarrow\kstwoKS_Z^k\to\kstwoKS_Y^k$ given by the fibre power of the pullback map and
  of the isogeny.
\end{enumerate}
\formal{GlobalTheory.UPel.HT.ES.uIdeal}\formal{GlobalTheory.UPel.HT.ES.IsPIsog}\formal{GlobalTheory.UPel.HT.ES.IsTj}\formal{UpperTheory.Gal.crystallineComparison.Mod.TjD}\formal{UpperTheory.Gal.crystallineComparison.Mod.Univ}\formal{UpperTheory.Gal.crystallineComparison.Mod.TjE}
\end{definition}

The local hypotheses are collected in $\mathrm{Std}(f,b,N,r,\mathfrak w)$: $O$ is a discrete valuation ring with
$p$ in its maximal ideal and residue field of order $p^f$, $N$ is a level, $m=a=1$,
$(p,\omega-r)\subseteq\mathfrak w$, $\Nm\mathfrak w=p^f$, and $t_1(\mathfrak w)$ lies in the maximal
ideal of $O$.

\begin{theorem}[Congruence relation]\label{pf:ks2:congks}
Let $O$ be a discrete valuation domain of characteristic $0$, $b$ a base over $O$, and assume
$\mathrm{Std}(f,b,N,r,\mathfrak w)$ with $\mathfrak w$ maximal. Let $M$ be a model of level $N$,
$\varpi$ a uniformizer of $O$, $\mathfrak a$ an ideal with $\mathfrak a+\mathfrak w=\kstwoOF$
containing a power of $(p,\omega-r)$, $Y$ an \'etale-level model of level $(N,\mathfrak a,\rho)$
with a flat morphism $\pi\colon Y\to M$ compatible with the universal objects, and $k\in\N$. Then
there are $T_j$-correspondences $T_0,\dots,T_n$ on $M$ of degrees $p^{fj}$, lifts $L_j$ of $T_j$ to
$Y$, and integers $d_j$ such that the $k$-th Kuga--Sato powers of the $L_j$ are correspondences of
$\kstwoKS_Y^k$ away from $p$ of degree $d_j$, lying over $\Spec O$, and
$\mathrm{PC}(p,p^f,n,\varpi,(L_j^{(k)})_j)$ holds on $\kstwoKS_Y^k$.
\formallabel{UpperTheory.Gal.crystallineComparison.Mod.congKS_link}
\end{theorem}

\begin{proof}
The theorem follows from four statements:
\begin{enumerate}[label=(\roman*)]
\item (existence of flat $T_j$-correspondences,
  \lref{UpperTheory.Gal.crystallineComparison.Mod.tjExF_holds}) for every
  $0\le j\le n$ a $T_j$-correspondence $T_j$ of degree $p^{fj}$ on $M$ with $T_j.s$ flat over $O$
  exists; it is assembled from two statements: a correspondence $M\leftarrow Y_j\to M$ with both legs
  finite, $Y_j$ with integrally closed local rings and flat over $O$, carrying a universal
  $T_j$-isogeny of degree $p^{fj}$ (\cref{pf:ks2:esnorm}),
  and a correspondence of the same kind whose legs factor through finite \'etale maps to $M[1/p]$
  (\lref{UpperTheory.Gal.crystallineComparison.Mod.esGen_link}), the latter resting on the
  representability of $T_j$-isogenies (\cref{pf:ks2:classical:tj});
\item (lifts, \lref{UpperTheory.Gal.crystallineComparison.Mod.tjEEx_holds}) every such
  $T_j$ has a lift $L_j$ to $Y$;
\item (ordinary reading, \cref{pf:ks2:ordkse}) the $k$-th Kuga--Sato powers of the lifts are
  correspondences away from $p$ of the degrees $d_j=\#\mathrm{Gr}_j((\kstwoOF/\mathfrak w)^n)$,
  lying over $\Spec O$, which satisfy the \emph{ordinary reading} $\mathrm{Ord}$ of
  \cref{pf:ks2:def:ordread} at the generic points of the special fibre;
\item (combinatorics, \cref{pf:ks2:ptidoford}) the ordinary reading implies the pointwise
  identity $\mathrm{PC}$.
\end{enumerate}
Choosing $T_j$ by (i), $L_j$ by (ii), the integers $d_j$ and the properties of the Kuga--Sato
powers by (iii), and applying (iv) gives the theorem.
\end{proof}

\subsubsection{The combinatorial identity}

\begin{definition}[Ordinary reading]\label{pf:ks2:def:ordread}
Let $q,n\in\N$, $\phi\colon P\to P$ a map of sets and $v_j\colon I_j\to P$ ($0\le j\le n$) maps from
finite sets. An \emph{ordinary reading} of $v$
consists of a finite field $k$ with $q$ elements, an $n$-dimensional $k$-vector space $H$, a line
$C\subseteq H$, a map $A$ from the set of subspaces of $H/C$ to $P$, and bijections $e_j$ from
$I_j$ to the set $\mathrm{Gr}_j(H)$ of $j$-dimensional subspaces of $H$, such that for all $j$ and
$i\in I_j$, with $W=e_j(i)$ and $\bar W$ its image in $H/C$:
\[
 v_j(i)=\phi\bigl(A(\bar W)\bigr)\ \text{if } C\subseteq W,\qquad v_j(i)=A(\bar W)\ \text{if }
 C\not\subseteq W.
\]
For correspondences $T_0,\dots,T_n$ of $Y$ away from $\ell$ and $\varpi\in\Gamma(Y,\OO_Y)$,
$\mathrm{Ord}(\ell,q,n,\varpi,T)$ means:
in the situation of \cref{pf:ks2:def:ptid}(4) (a valuation ring $V$, a point $y$ over a generic point
of $\{\varpi=0\}$, the $q$-Frobenius $\phi$ of the residue field), the maps $z\mapsto\bar z$ from the
sections of $T_j$ over $y$ to $Y(\kappa)$ admit an ordinary reading with respect to $\phi$.
\formal{UpperTheory.Gal.crystallineComparison.OrdRead}\formal{UpperTheory.Gal.crystallineComparison.OrdPt}
\end{definition}

\begin{lemma}[Subspaces and a line]\label{pf:ks2:countin}
Let $k$ be a field, $H$ a finite-dimensional $k$-vector space, $C\subseteq H$ a line and
$\pi\colon H\to H/C$ the projection.
\begin{enumerate}
\item $W\mapsto\pi(W)$ is a bijection from the subspaces containing $C$ onto the subspaces of $H/C$,
  and $\dim W=\dim\pi(W)+1$ for $W\supseteq C$;
\item if $C\not\subseteq W$, then $\dim W=\dim\pi(W)$;
\item if $k$ is finite with $q$ elements and $E\subseteq H/C$ is a subspace, the number of subspaces
  $W$ with $C\not\subseteq W$ and $\pi(W)=E$ is $q^{\dim E}$.
\end{enumerate}
\formallabel{UpperTheory.Gal.crystallineComparison.countIn_of}\formallabel{UpperTheory.Gal.crystallineComparison.countOut_of}
\end{lemma}

\begin{proof}
(1) is the correspondence theorem: the inverse is $E\mapsto\pi^{-1}(E)$, and $\ker(\pi|_W)=C$ for
$W\supseteq C$. (2): $\ker(\pi|_W)=W\cap C=0$ since $C$ is a line not contained in $W$. (3): Let
$\tilde E=\pi^{-1}(E)$, of dimension $\dim E+1$. A subspace $W$ with $C\not\subseteq W$ and
$\pi(W)=E$ satisfies $W\subseteq\tilde E$, $W\cap C=0$ and $\dim W=\dim E$ by (2), i.e.\ $W$ is a
complement of $C$ in $\tilde E$; conversely every complement $W$ of $C$ in $\tilde E$ has
$\pi(W)=E$ and $C\not\subseteq W$. Fix one complement $W_0$; then $\pi|_{W_0}\colon W_0\to E$ is an
isomorphism, and the complements of $C$ in $\tilde E=W_0\oplus C$ are exactly the graphs
$\{w+\varphi(w):w\in W_0\}$ of the linear maps $\varphi\colon W_0\to C$, distinct maps giving distinct
graphs. Their number is $\#\Hom_k(W_0,C)=q^{\dim E}$.
\end{proof}

\begin{lemma}[Decomposition]\label{pf:ks2:orddecomp}
Let $(k,H,C,A,e)$ be an ordinary reading of $v$ as in \cref{pf:ks2:def:ordread}, and for $i\ge0$
let $\mathcal A_i$ be the multiset $\{A(E):E\subseteq H/C,\ \dim E=i\}$. Then $\mathcal A_n=0$,
$\{v_0(i):i\in I_0\}=\mathcal A_0$, and for $0\le j<n$
\begin{equation}\label{pf:ks2:eq:decomp}
\{v_{j+1}(i):i\in I_{j+1}\}=\phi(\mathcal A_j)+q^{\,j+1}\mathcal A_{j+1},
\end{equation}
where $\phi(\cdot)$ is applied elementwise and $q^{j+1}\mathcal A$ is the multiset $\mathcal A$
repeated $q^{j+1}$ times.
\formallabel{UpperTheory.Gal.crystallineComparison.ordDecomp_of}
\end{lemma}

\begin{proof}
(The formal statement takes the two counting statements of \cref{pf:ks2:countin} as hypotheses.)
$\dim H/C=n-1$, so $H/C$ has no $n$-dimensional subspace and $\mathcal A_n=0$. Transporting along
the bijections $e_j$, the multiset $\{v_j(i)\}$ is the sum over $W\in\mathrm{Gr}_j(H)$ of
$\{\phi(A(\pi W))\}$ for $W\supseteq C$ and of $\{A(\pi W)\}$ for $W\not\supseteq C$. For $j=0$ the
only subspace is $W=0\not\supseteq C$, with $\pi W=0$, giving $\mathcal A_0$. For $j+1$, by
\cref{pf:ks2:countin}(1) the subspaces $W\supseteq C$ of dimension $j+1$ correspond bijectively to the
$E=\pi W$ of dimension $j$, contributing $\phi(\mathcal A_j)$; by (2) and (3) the subspaces
$W\not\supseteq C$ of dimension $j+1$ have $\pi W$ of dimension $j+1$, and each $E$ of dimension $j+1$
arises from exactly $q^{j+1}$ of them, contributing $q^{j+1}\mathcal A_{j+1}$.
\end{proof}

\begin{lemma}[Telescoping]\label{pf:ks2:telescope}
Let $\phi\colon P\to P$, let $\mathcal T_0,\dots,\mathcal T_n$ and $\mathcal A_0,\mathcal A_1,\dots$ be
finite multisets on $P$ with $\mathcal A_n=0$, $\mathcal T_0=\mathcal A_0$ and $\mathcal
T_{j+1}=\phi(\mathcal A_j)+q^{j+1}\mathcal A_{j+1}$ for $0\le j<n$. Then
\[
\sum_{j\ \mathrm{even}}q^{\binom j2}\phi^{\,n-j}(\mathcal T_j)=\sum_{j\ \mathrm{odd}}q^{\binom
j2}\phi^{\,n-j}(\mathcal T_j)\qquad(0\le j\le n).
\]
\formallabel{UpperTheory.Gal.crystallineComparison.telescope_of}
\end{lemma}

\begin{proof}
Put $\mathcal A_{-1}=0$ and $X_i=q^{\binom{i+1}2}\phi^{\,n-i}(\mathcal A_i)$ for $-1\le i\le n$; thus
$X_{-1}=X_n=0$. Since $\binom j2+j=\binom{j+1}2$ and $\phi$ commutes with repetition, the hypotheses
give, for $0\le j\le n$ (for $j=0$ use $\mathcal T_0=\mathcal A_0=\phi(\mathcal A_{-1})+q^0\mathcal
A_0$),
\[
q^{\binom j2}\phi^{\,n-j}(\mathcal T_j)=q^{\binom j2}\phi^{\,n-j+1}(\mathcal A_{j-1})+q^{\binom
j2+j}\phi^{\,n-j}(\mathcal A_j)=X_{j-1}+X_j .
\]
Each $X_i$ with $0\le i\le n-1$ occurs exactly twice in $\sum_{j=0}^n(X_{j-1}+X_j)$, namely for $j=i$
and $j=i+1$, which have opposite parities; $X_{-1}$ and $X_n$ vanish. Hence both sides of the
claimed identity equal $\sum_{i=0}^{n-1}X_i$.
\end{proof}

\begin{proposition}\label{pf:ks2:ordcount}
An ordinary reading of $v$ implies $\mathrm{Alt}(q,n,\phi,v)$, i.e.\ \eqref{pf:ks2:eq:altid}.
\formallabel{UpperTheory.Gal.crystallineComparison.ordCount_holds}
\end{proposition}

\begin{proof}
Apply \cref{pf:ks2:telescope} to $\mathcal T_j=\{v_j(i):i\in I_j\}$ and the multisets $\mathcal A_i$
of \cref{pf:ks2:orddecomp}; the hypotheses of the former are the conclusions of the latter, and the
conclusion is \eqref{pf:ks2:eq:altid}.
\end{proof}

\begin{corollary}\label{pf:ks2:ptidoford}
For all $\ell,q,n$, $Y$, $\varpi$ and correspondences $T_0,\dots,T_n$ of $Y$ away from $\ell$,
$\mathrm{Ord}(\ell,q,n,\varpi,T)$ implies $\mathrm{PC}(\ell,q,n,\varpi,T)$.
\formallabel{UpperTheory.Gal.crystallineComparison.ptIdOfOrd_holds}
\end{corollary}

\begin{proof}
In the situation of $\mathrm{PC}$ we are given bijections $z_j\colon I_j\to\{\text{sections of }
T_j\text{ over }y\}$. $\mathrm{Ord}$ provides an ordinary reading $(k,H,C,A,e)$ of the maps
$z\mapsto\bar z$ on the sets of sections; composing $e_j$ with $z_j$ gives an ordinary reading of
$i\mapsto\overline{z_j(i)}$, and \cref{pf:ks2:ordcount} gives $\mathrm{Alt}$.
\end{proof}

\subsubsection{The ordinary reading on Kuga--Sato varieties}

\begin{proposition}\label{pf:ks2:ordkse}
In the situation of \cref{pf:ks2:congks}, let $T_0,\dots,T_n$ be $T_j$-correspondences of degrees
$p^{fj}$ on $M$ and $L_j$ lifts to $Y$ (with $\pi$ flat). Then, with
$d_j=\#\mathrm{Gr}_j((\kstwoOF/\mathfrak w)^n)$, the $k$-th Kuga--Sato powers $L_j^{(k)}$ are
correspondences of $\kstwoKS_Y^k$ away from $p$ of degree $d_j$, both legs of $L_j^{(k)}$ have the same
composite to $\Spec O$, and $\mathrm{Ord}(p,p^f,n,\varpi,(L_j^{(k)}))$ holds.
\formallabel{UpperTheory.Gal.crystallineComparison.Mod.ordKSE_holds}
\end{proposition}

\begin{proof}
The formal statement of this proposition takes generic ordinariness and the existence of
quotients by finite \'etale subgroup schemes as hypotheses; they are supplied by
\cref{pf:ks2:classical:ord} and \cref{pf:ks2:quotobj}. The proof combines
four facts.
First, $L_j^{(k)}$ is the base change of $T_j.s$ along $\kstwoKS_Y^k\to M$, so the finiteness,
\'etaleness over $M[1/p]$ and the degree count follow from
\cref{pf:ks2:min}(7). Second, both legs of $L_j$ have the same composite to $\Spec O$ because
those of $T_j$ do. Third, a morphism into $\kstwoKS_Y^k$ is determined by its image in $M$ together
with its compositions with the $k$ projections to $\kstwoKS_Y$ and the \'etale level sections,
followed by the map to $\mathcal A_M$; and $L_j^{(k)}$ acts on these
coordinates through the isogeny of $T_j$.
Fourth, these are the hypotheses of the embedding criterion \cref{pf:ks2:ordemb} below, applied to
the flat composite $h\colon\kstwoKS_Y^k\to Y\to M$ ($\pi$ is flat and $\kstwoKS_Y^k\to Y$ is smooth),
the family of sections of the third fact, and the correspondences $C_j=L_j^{(k)}$ with their maps to
the sources of the $T_j$; it gives $\mathrm{Ord}(p,p^f,n,\varpi,(L_j^{(k)}))$.
\end{proof}

\subsubsection{Normal $T_j$-correspondences, quotients and the reading at ordinary points}
\label{pf:ks2:sec:ordproofs}

For a $T_j$-correspondence $T$ on $M$ (\cref{pf:ks2:def:tj}) and a morphism $y\colon S\to Y_j$, the
isogeny of $T$, pulled back along $y$, maps a point $a\colon S\to\mathcal A_M$ of $\mathcal A_M$ over
$s\circ y$ to a point $\psi_y(a)$ of $\mathcal A_M$ over $t\circ y$; we say that $T$ \emph{acts on
points} in this way.

\begin{proposition}[Normal $T_j$-correspondences]\label{pf:ks2:esnorm}
Let $R$ be a discrete valuation ring of characteristic $0$ with $p$ in its maximal ideal, $b$ a base
of $D$ over $R$ with $m=1$, $N$ a level, $r\in\Z$, $\mathfrak w\supseteq(p,\omega-r)$ a maximal ideal
of $\kstwoOF$, $0\le j\le n$, and $c=p^e$ with $e>0$ if $j>0$. For every model $M$ of level $N$ over
$R$ there are a scheme $Y$, finite morphisms $s,t\colon Y\to M$ with the same composite to $\Spec R$,
objects $o,o'$ of level $N$ over $Y$ with morphisms to $\mathcal A_M$ over $s$ and over $t$, and a
$p$-isogeny $\psi\colon A_o\to A_{o'}$ of degree $c$ and of type $T_j$ at $\mathfrak w$, such that
$Y$ is flat over $R$, all local rings of $Y$ are integrally closed, and $\psi$ is universal among
such isogenies over schemes on which $p$ is invertible (as in \cref{pf:ks2:def:tj}(2)).
\formallabel{UpperTheory.Gal.crystallineComparison.Mod.esNorm_link}
\end{proposition}

\begin{proof}
By \lref{UpperTheory.Gal.crystallineComparison.Mod.esGen_link} the corresponding data exist over
the open subscheme $M[1/p]$ where $p$ is invertible: there are a scheme $Y_0$, morphisms
$s_0,t_0\colon Y_0\to M$ with the same composite to $\Spec R$ which factor through finite \'etale
morphisms $Y_0\to M[1/p]$, objects $o_0,o_0'$ over $Y_0$ with morphisms to $\mathcal A_M$ over $s_0$
and $t_0$, and a universal $p$-isogeny $\psi_0\colon A_{o_0}\to A_{o_0'}$ of degree $c$ and type $T_j$
(they come from the representability of $T_j$-isogenies, \cref{pf:ks2:classical:tj}, applied to the
universal object over $M[1/p]$). We extend these data over $M$.
\begin{enumerate}
\item (Normalization.) By \lref{GlobalTheory.UPel.HT.ES.NormCover.exists_finite_normal_model}
  there are a locally Noetherian scheme $Y$, a finite morphism $s\colon Y\to M$ and a dominant open
  immersion $i\colon Y_0\to Y$ with $s\circ i=s_0$, such that $Y$ is flat over $R$, its local rings
  are integrally closed, and every morphism to $Y$ from a scheme on which $p$ is invertible factors
  through $i$ (the normalization of $M$ in $Y_0$).
\item (The second leg.) By \lref{GlobalTheory.ExtT.exists_finite_extension} there is a finite
  morphism $t\colon Y\to M$ with $t\circ i=t_0$ and with the same composite to $\Spec R$ as $s$ (the
  target of $\psi_0$ extends over every discrete valuation ring mapping to $Y$, because an abelian
  variety isogenous to one with good reduction has good reduction).
\item (Local rings.) $Y$ is reduced
  (\lref{GlobalTheory.UPel.HT.ES.isReduced_normalModel_of_generic}), so every local ring $A$ of
  $Y$ is a reduced Noetherian local ring which is integrally closed in its total ring of fractions
  $Q(A)$. The ring $Q(A)$ is the product of the fields of fractions of $A/\mathfrak q$ over the
  finitely many minimal primes $\mathfrak q$; its idempotents are integral over $A$, hence lie in the
  local ring $A$, so they are $0$ and $1$. Thus $A$ has a single minimal prime, which is $0$ as $A$ is
  reduced: $A$ is a domain.
\item (Objects.) Let $o,o'$ be the pullbacks of $\mathcal A_M$ along $s$ and $t$. As $s\circ
  i=s_0$ and $t\circ i=t_0$, the
  morphisms $o_0\to\mathcal A_M$ and $o_0'\to\mathcal A_M$ factor through morphisms $o_0\to o$ and
  $o_0'\to o'$ over $i$.
\item (The isogeny.) By \lref{GlobalTheory.UPel.HT.ES.PIsogExt.exists_isPIsog_extend} (over a
  locally Noetherian scheme whose local rings are integrally closed domains, a $p$-isogeny between
  the restrictions of two objects to a dense open subscheme extends to a $p$-isogeny of the same
  degree) there is a $p$-isogeny $\psi\colon A_o\to A_{o'}$ of degree $c$ restricting to $\psi_0$.
\item (Type and universality.) The type-$T_j$ condition only concerns geometric points of
  characteristic $0$; these factor through $Y_0$ by (1), and there $\psi$ restricts to $\psi_0$, for
  which the condition holds. Likewise universality only concerns schemes on which $p$ is invertible,
  whose morphisms to $Y$ factor uniquely through the open immersion $i$; so it follows from the
  universality of $\psi_0$.
\end{enumerate}
\end{proof}

\begin{proposition}[Quotients by subgroup schemes]\label{pf:ks2:quotobj}
Let $R$ be a ring with a base $b$ of $D$, $m=1$, $N$ a level, $r\in\Z$, $\mathfrak w\supseteq
\mathfrak u=(p,\omega-r)$ a maximal ideal, $0\le j\le n$ and $c=p^e$ with $e>0$ if $j>0$. Let $S$ be
a locally Noetherian $R$-scheme on which $p$ is invertible, $o$ an object of level $N$ over $S$, and
$K\subseteq A_o$ a closed subgroup scheme which is finite and \'etale of degree $\Nm\mathfrak w^{\,j}$
over $S$, stable under $O$ and killed by $\mathfrak w$. Then there are an object $o'$ of level $N$
over $S$ and a $p$-isogeny $\psi\colon A_o\to A_{o'}$ of degree $c$ and type $T_j$ at $\mathfrak w$
such that a point $x$ of $A_o$ (over any $S$-scheme) lies in $K$ if and only if $\psi(x)=0$ and $x$
is killed by a power of $\mathfrak u$.
\formallabel{UpperTheory.Gal.crystallineComparison.Mod.quotObj_link}
\end{proposition}

\begin{proof}
By \lref{UpperTheory.exists_noetherian_base_link} we may assume that $R$ is Noetherian and that
$p$ is a unit in $R$ (the structure map of $S$ factors through such a ring carrying a base of $D$, and
objects over $S$ for the two bases correspond). By \cref{pf:ks2:classical:tj} there are a finite
\'etale $S$-scheme $Y$ and a $p$-isogeny $\psi_Y$ of degree $c$ and type $T_j$ out of the pullback of
$o$ to $Y$, universal among such isogenies out of pullbacks of $o$. For a morphism $\sigma\colon
S'\to Y$ let $K_\sigma$ be the group of points of $A_o$ over $S'$ killed by $\sigma^*\psi_Y$ and by a
power of $\mathfrak u$.
\begin{enumerate}
\item By \lref{UpperTheory.exists_clopen_subgroup_locus} there is an open and closed subscheme
  $W\subseteq Y$ such that $K_\sigma$ is the pullback of $K$ whenever $\sigma$ factors through $W$, and
  a geometric point of $Y$ lies in $W$ if and only if this holds at it (both subgroups are finite and
  \'etale, and the locus where two finite \'etale subgroup schemes agree is open and closed).
\item Over every geometric point $x$ of $S$ there is a geometric point of $W$: the quotient of
  $A_{o,x}$ by $K_x$, with its induced structures, is a $p$-isogeny of degree $c$ and type $T_j$ with
  $K_x$ as the $\mathfrak u$-part of its kernel
  (\lref{UpperTheory.repData_geometric_points_of_subData}); by universality it is the fibre of
  $\psi_Y$ at a geometric point of $Y$ over $x$, which lies in $W$ by (1).
\item There is at most one: two geometric points of $W$ over $x$ carry $p$-isogenies of degree $c$
  and type $T_j$ out of $A_{o,x}$ with the same $\mathfrak u$-part $K_x$ of the kernel, which are
  isomorphic, so the two points coincide
  (\lref{UpperTheory.subData_points_correspond_repData}).
\item Hence $W\to S$ is finite and \'etale with exactly one geometric point over every geometric
  point, i.e.\ an isomorphism; its inverse is a section $\sigma\colon S\to Y$ with image in $W$
  (\lref{UpperTheory.exists_section_of_etale_clopen_geomPoint_unique}).
\item Pulling $\psi_Y$ back along $\sigma$ gives an object $o'$ and a $p$-isogeny $\psi\colon
  A_o\to A_{o'}$ of degree $c$ and type $T_j$
  (\lref{UpperTheory.exists_isPIsog_isTj_of_section_repData}); as $\sigma$ factors through $W$,
  $K_\sigma$ is $K$, which is the assertion.
\end{enumerate}
\end{proof}

The next four lemmas are elementary.

\begin{lemma}\label{pf:ks2:resclosed}
Let $V$ be a valuation ring whose fraction field is algebraically closed. Then the residue field of
$V$ is algebraically closed.
\formallabel{UpperTheory.Gal.crystallineComparison.Mod.resClosed_holds}
\end{lemma}

\begin{proof}
Let $\bar f$ be a monic polynomial of positive degree over the residue field $\kappa$, and $f\in V[X]$
a monic lift. Over the fraction field, $f=\prod_i(X-a_i)$. Each $a_i$ is integral over $V$, hence lies
in $V$, since valuation rings are integrally closed. Reducing, $\bar f=\prod_i(X-\bar a_i)$ has a root
in $\kappa$.
\end{proof}

\begin{lemma}\label{pf:ks2:gendown}
Let $O$ be a discrete valuation ring with uniformizer $\varpi$, $s\colon X\to\Spec O$ a morphism and
$h\colon Y\to X$ a flat morphism. If $y\in Y$ is a generic point of $\{\varpi=0\}$ (for $\varpi$ pulled
back to $Y$), then $h(y)$ lies over the closed point of $\Spec O$ and has no proper generization lying
over the closed point.
\formallabel{UpperTheory.Gal.crystallineComparison.Mod.genDown_holds}
\end{lemma}

\begin{proof}
Since $\varpi$ vanishes at $y$, the image of $y$ in $\Spec O$ lies in $V(\varpi)$, the closed point.
Let $z'\in X$ lie over the closed point and generize $h(y)$. Flat morphisms satisfy going down, so
there is a generization $y'$ of $y$ with $h(y')=z'$. Then $y'$ lies over the closed point, i.e.\
$\varpi(y')=0$, and $y'=y$ because $y$ has no proper generization in $\{\varpi=0\}$; hence $z'=h(y)$.
\end{proof}

\begin{lemma}[Reduction of a subspace]\label{pf:ks2:redpts}
Let $k$ be a finite field, $H$ a finite-dimensional $k$-vector space, $C\subseteq H$ a subspace with
projection $\pi\colon H\to H/C$, and $\rho\colon H\to P$ a map to a set with $\rho(a)=\rho(b)$ iff
$a-b\in C$. Then for every subspace $D\subseteq H$, $\#\rho(D)=(\#k)^{\dim\pi(D)}$, and
$\rho(D)=\rho(D')$ for every subspace $D'$ with $\pi(D')=\pi(D)$.
\formallabel{UpperTheory.Gal.crystallineComparison.Mod.redPts_holds}
\end{lemma}

\begin{proof}
By hypothesis $\rho=\iota\circ\pi$ for an injective map $\iota\colon H/C\to P$. Hence
$\rho(D)=\iota(\pi(D))$ is in bijection with the $k$-vector space $\pi(D)$, which has
$(\#k)^{\dim\pi(D)}$ elements, and it depends only on $\pi(D)$.
\end{proof}

\begin{lemma}[Constructing ordinary readings]\label{pf:ks2:ordreadof}
Let $q,n\in\N$, $\phi\colon P\to P$ and $v_j\colon I_j\to P$ ($0\le j\le n$). Let $k$ be a field with
$q$ elements, $H$ an $n$-dimensional $k$-vector space, $C\subseteq H$ a line with projection
$\pi\colon H\to H/C$, and $e_j\colon I_j\to\mathrm{Gr}_j(H)$ bijections such that:
\begin{enumerate}[label=(\alph*)]
\item if $i,i'\in I_j$, $C\not\subseteq e_j(i)$, $C\not\subseteq e_j(i')$ and $\pi e_j(i)=\pi
  e_j(i')$, then $v_j(i)=v_j(i')$;
\item if $i\in I_{j+1}$, $i'\in I_j$, $C\subseteq e_{j+1}(i)$, $C\not\subseteq e_j(i')$ and $\pi
  e_{j+1}(i)=\pi e_j(i')$, then $v_{j+1}(i)=\phi(v_j(i'))$.
\end{enumerate}
Then $v$ admits an ordinary reading (\cref{pf:ks2:def:ordread}) with these $k,H,C,e$.
\formallabel{UpperTheory.Gal.crystallineComparison.Mod.ordReadOf_holds}
\end{lemma}

\begin{proof}
Let $E\subseteq H/C$ be a subspace and $d=\dim E\le n-1$. By \cref{pf:ks2:countin}(3) there is a
subspace $W'$ with $C\not\subseteq W'$ and $\pi(W')=E$, and $\dim W'=d$ by \cref{pf:ks2:countin}(2);
put $A(E)=v_d(e_d^{-1}(W'))$, which does not depend on the choice of $W'$ by (a). Let $i\in I_j$ and
$W=e_j(i)$. If $C\not\subseteq W$, then $W$ is an admissible choice for $E=\pi(W)$, so
$v_j(i)=A(\pi W)$. If $C\subseteq W$, then $\dim\pi(W)=j-1$ by \cref{pf:ks2:countin}(1), and
$A(\pi W)=v_{j-1}(i')$ for some $i'\in I_{j-1}$ with $C\not\subseteq e_{j-1}(i')$ and $\pi
e_{j-1}(i')=\pi W$; by (b), $v_j(i)=\phi(v_{j-1}(i'))=\phi(A(\pi W))$. These are the defining
identities of an ordinary reading.
\end{proof}

\begin{proposition}[The reading at ordinary points]\label{pf:ks2:ew}
Assume the conclusion of \cref{pf:ks2:quotobj}. Let $O$ be a discrete valuation domain of
characteristic $0$ and assume $\mathrm{Std}(f,b,N,r,\mathfrak w)$ with $\mathfrak w$ maximal; let $M$
be a model of level $N$ over $O$, $T_j$ ($0\le j\le n$) $T_j$-correspondences of degree $p^{fj}$ on
$M$, and $k=\kstwoOF/\mathfrak w$. Let $V$ be a valuation ring with algebraically closed fraction
field and algebraically closed residue field $\kappa$, with $p\ne0$ in $V$, let $x\colon\Spec V\to
M$ map the closed point over the closed point of $\Spec O$, assume that the geometric point $\bar
x=x|_{\Spec\kappa}$ is $\mathfrak w$-ordinary, and let $\phi$ be the $p^f$-power Frobenius of
$\kappa$. Then there are an $n$-dimensional $k$-vector space $H$, a line $C\subseteq H$ with
projection $\pi\colon H\to H/C$, and for every $j$ a bijection $e_j$ from the set of sections of
$T_j$ over $x$ (the $z\colon\Spec V\to Y_j$ with $s\circ z=x$) onto $\mathrm{Gr}_j(H)$, such that,
writing $\bar z=z|_{\Spec\kappa}$:
\begin{enumerate}[label=(\alph*)]
\item if $z,z'$ are sections of $T_j$ with $C\not\subseteq e_j(z)$, $C\not\subseteq e_j(z')$ and
  $\pi e_j(z)=\pi e_j(z')$, then $t(\bar z)=t(\bar z')$ and $\psi_{\bar z}(a)=\psi_{\bar z'}(a)$ for
  all points $a$ of $\mathcal A_M$ over $\bar x$;
\item if $z$ is a section of $T_{j+1}$ and $z'$ one of $T_j$ with $C\subseteq e_{j+1}(z)$,
  $C\not\subseteq e_j(z')$ and $\pi e_{j+1}(z)=\pi e_j(z')$, then $t(\bar z)=t(\bar
  z')\circ\Spec\phi$ and $\psi_{\bar z}(a)=\psi_{\bar z'}(a)\circ\Spec\phi$ for all points $a$ of
  $\mathcal A_M$ over $\bar x$.
\end{enumerate}
\formallabel{UpperTheory.Gal.crystallineComparison.Mod.eW_holds}
\end{proposition}

\begin{proof}
Let $H$ be the set of points of $\mathcal A_M$ over $x$ killed by $\mathfrak w$, a module over
$\kstwoOF$ killed by $\mathfrak w$, i.e.\ a $k$-vector space, of dimension $n$
(\lref{UpperTheory.Gal.crystallineComparison.Mod.rankN_holds}: as $p\ne0$ in $V$ and $m=1$, the
$\mathfrak w$-torsion of the abelian variety over the fraction field is free of rank $n$ over $k$, and
it extends over $V$ by properness). Let $\rho\colon H\to\mathcal A_M(\bar x)$ be the reduction to the
closed point. As $\bar x$ is $\mathfrak w$-ordinary, there is a line $C\subseteq H$ with
$\rho(a)=\rho(a')$ iff $a-a'\in C$
(\lref{UpperTheory.Gal.crystallineComparison.Mod.lineC_holds}: the kernel of reduction on $H$
has $\#k$ elements, since the $\mathfrak w$-torsion of the fibre at $\bar x$ has $\Nm\mathfrak
w^{\,n-1}$ points). Using the existence of quotients (\cref{pf:ks2:quotobj}) and the universality of
$T_j$, \lref{UpperTheory.Gal.crystallineComparison.Mod.secEquiv_holds} gives bijections $e_j$
from the sections of $T_j$ over $x$ onto $\mathrm{Gr}_j(H)$ such that $a\in e_j(z)$ iff the point $a$
is killed by the isogeny of $T_j$ at $z$ (the $\mathfrak u$-part of the kernel of a $T_j$-isogeny over
the characteristic-$0$ point is a $j$-dimensional subspace of $H$, and every such subspace arises
from exactly one section).

Let $o_1$ be the pullback of $\mathcal A_M$ along $\bar x$, an object over $\kappa$ whose points are
the points of $\mathcal A_M$ over $\bar x$. For a section $z$ of $T_j$ over $x$, the isogeny of $T_j$
at $\bar z$ pulls back to an isogeny $\psi_z$ out of $A_{o_1}$; let $U(z)$ be the subgroup
scheme of points of $\ker\psi_z$ killed by a power of $\mathfrak u$. By
\lref{UpperTheory.Gal.crystallineComparison.Mod.kerRead_holds}, $U(z)$ is finite of order
$p^{fj}$, and its $\kappa$-points are the reductions of the points of $e_j(z)$:
$U(z)(\kappa)=\rho(e_j(z))$. By \cref{pf:ks2:redpts}, $U(z)(\kappa)$ has $p^{f\dim\pi e_j(z)}$
elements and depends only on $\pi e_j(z)$.

(a) Here $\dim\pi e_j(z)=\dim\pi e_j(z')=j$ (\cref{pf:ks2:countin}(2)), so $U(z)$ and $U(z')$ have the
same $p^{fj}$ points, as many as their order. By
\lref{UpperTheory.Gal.crystallineComparison.Mod.caseE_holds} (two isogenies of the type of $T_j$
out of the same abelian variety over an algebraically closed field, whose $\mathfrak u$-parts of the
kernels have the same points, as many as their order, have the same target point and act in the
same way on points), $t(\bar z)=t(\bar z')$ and $\psi_{\bar z}=\psi_{\bar z'}$ on points.

(b) Here $U(z')$ has $p^{fj}$ points, as many as its order, while $U(z)$ has order $p^{f(j+1)}$ and the
same points as $U(z')$, because $\pi e_{j+1}(z)=\pi e_j(z')$. By
\lref{UpperTheory.Gal.crystallineComparison.Mod.caseW_holds} (in this situation, in residue
characteristic $p$, the target of $\psi_z$ is the $p^f$-Frobenius twist of the target of
$\psi_{z'}$, compatibly with the maps on points), $t(\bar z)=t(\bar z')\circ\Spec\phi$ and
$\psi_{\bar z}(a)=\psi_{\bar z'}(a)\circ\Spec\phi$.
\end{proof}

\begin{proposition}[An embedding criterion for the ordinary reading]\label{pf:ks2:ordemb}
Assume generic ordinariness (\cref{pf:ks2:classical:ord}) and the conclusion of
\cref{pf:ks2:quotobj}. Let $O$, $b$, $N$, $r$, $\mathfrak w$, $M$ and the $T_j$ be as in
\cref{pf:ks2:ew}, and let $\varpi$ be a uniformizer of $O$. Let $h\colon Y\to M$ be flat and
$(a_i)_{i\in I}$ a family of morphisms $a_i\colon Y\to\mathcal A_M$ lifting $h$ such that any two
morphisms $u,u'\colon X'\to Y$ with $hu=hu'$ and $a_iu=a_iu'$ for all $i$ are equal. Let
$C_j=(Y\xleftarrow{p_j}Z_j\xrightarrow{r_j}Y)$ ($0\le j\le n$) be correspondences of $Y$ away from
$p$ with morphisms $g_j\colon Z_j\to Y_j$ such that the square formed by $p_j$, $g_j$, $h$ and $s$ is
cartesian, $h\circ r_j=t\circ g_j$, and $a_i\circ r_j=\psi_{g_j}(a_i\circ p_j)$ for all $i$. Then
$\mathrm{Ord}(p,p^f,n,\varpi,C)$ holds, $\varpi$ being viewed as a function on $Y$.
\formallabel{UpperTheory.Gal.crystallineComparison.Mod.ordEmb_holds}
\end{proposition}

\begin{proof}
Let $V$, $y\colon\Spec V\to Y$ (with closed point over a generic point of $\{\varpi=0\}$), the
residue field $\kappa$ and the $p^f$-Frobenius $\phi$ be as in the definition of $\mathrm{Ord}$
(\cref{pf:ks2:def:ordread}). By \cref{pf:ks2:resclosed}, $\kappa$ is algebraically closed. Put
$x=h\circ y$. By \cref{pf:ks2:gendown}, the closed point of $x$ lies over the closed point of $\Spec
O$ and has no proper generization in the special fibre of $M$. The ideal $\mathfrak w$ is the
inverse image of the maximal ideal of $O$ under $t_1$ (it is contained in it by $\mathrm{Std}$, and
it is maximal), so generic ordinariness shows that $\bar x$ is $\mathfrak w$-ordinary. Let $H$, $C$,
$\pi$ and $e_j$ be as in \cref{pf:ks2:ew} for $x$.

Sections of $C_j$ over $y$ correspond bijectively to sections of $T_j$ over $x$ through $z\mapsto
g_j\circ z$, as the square is cartesian. For such $z$ and $i\in I$ the hypotheses give, after
restriction to the closed point,
\[
a_i\circ r_j\circ\bar z=\psi_{g_j\bar z}(a_i\circ\bar y),\qquad h\circ r_j\circ\bar z=t\circ
g_j\circ\bar z .
\]
We apply \cref{pf:ks2:ordreadof} to the maps $z\mapsto r_j\circ\bar z$ from the sections of $C_j$
over $y$ to $Y(\kappa)$, with the bijections $z\mapsto e_j(g_j\circ z)$ and with $\phi$ acting on
$Y(\kappa)$ by precomposition with $\Spec\phi$. For (a) of that lemma, let $z,z'$ be sections of
$C_j$ whose subspaces do not contain $C$ and have the same image in $H/C$. By \cref{pf:ks2:ew}(a),
$t(g_j\bar z)=t(g_j\bar z')$ and $\psi_{g_j\bar z}=\psi_{g_j\bar z'}$ on points over $\bar x$; by the
displayed identities, $r_j\bar z$ and $r_j\bar z'$ have the same composites with $h$ and with all
$a_i$, so they are equal by the separation hypothesis. For (b), \cref{pf:ks2:ew}(b) shows in the same
way that $r_{j+1}\bar z$ and $(r_j\bar z')\circ\Spec\phi$ have the same composites with $h$ and with
all $a_i$, hence are equal. \Cref{pf:ks2:ordreadof} gives an ordinary reading, which is the
assertion of $\mathrm{Ord}(p,p^f,n,\varpi,C)$ at $y$.
\end{proof}

% ---------------------------------------------------------------------------
\subsection{Good reduction and congruences on Kuga--Sato towers}\label{pf:ks2:sec:good}

Throughout this subsection $\mathrm{Sc}(D,G)$ holds, $E$ is a number field, $e\colon F\cong E$, $g_b$ is
a characteristic-zero base over $E$ standard for $e$, $T$ is a Kuga--Sato tower for $(L,g_b,k)$, $N$
is a level, $(g_j)_{j\in J}$ is a family of Hecke data of level $N$ with values
$\theta\colon J\to A$ in a field, $w$ is a place of $E$ corresponding to the maximal ideal
$\mathfrak w$ of $\kstwoOF$, and $O_w$ is the valuation ring of $w$.

\subsubsection{Sites, class maps and covers}

\begin{definition}[Sites]\label{pf:ks2:def:site}
Let $q\ne p$ be a prime, $r\in\Z$ and $\mathfrak n$ an ideal of $\kstwoOF$. A \emph{site} at
$(w,q,r,\mathfrak w,\mathfrak n)$ consists of:
$\mathrm{Adm}(q,r,\mathfrak w)$; an isomorphism $O\otimes\Z_q\cong M_m(\kstwoOF\otimes\Z_q)$
(which makes $D$ a datum $D_q$ for the prime $q$); an
integer $f$ and a base $b$ of $D_q$ over $O_w$ matching $g_b$ under $O_w\to E$; a level $N_0$ prime
to $pq$ and an ideal $\mathfrak a$ with $\mathfrak n=N_0\mathfrak a$, $\mathfrak a+\mathfrak
w=\kstwoOF$ and $\mathfrak a\supseteq\mathfrak u^e$ for some $e$; the conditions
$\mathrm{Std}(f,b,N_0,r,\mathfrak w)$, $q\in w$ and $\Nm w=q^f$; a model $M$ of level $N_0$ of $D_q$
over $O_w$; an \'etale-level model $Y$ of level $(N_0,\mathfrak a,1)$ with a finite \'etale morphism
$\pi\colon Y\to M$ compatible with the universal objects; and, for every $k$, projectivity and
smoothness of some relative dimension of $\kstwoKS_Y^k\to\Spec O_w$.
\formal{UpperTheory.Gal.Int.Geo.Site}\formal{UpperTheory.Gal.Int.Geo.datAt}
\end{definition}

The following statements are proved in the formal development; we use them as stated.
\begin{enumerate}[label=(S\arabic*)]
\item\label{pf:ks2:S:site} (existence of sites, \lref{UpperTheory.Gal.Int.Geo.CovMain.siteEx_of})
  If $w$ corresponds to $\mathfrak w$, $\mathrm{Adm}(q,r,\mathfrak w)$ holds and $\mathfrak n$ is
  an auxiliary ideal, a site at
  $(w,q,r,\mathfrak w,\mathfrak n)$ exists. It is built from $\sfname{MIn}$: $D_q$ is a datum for $q$
  because $q$ splits in $\Q(\omega)$, $m=1$ and $E_\Lambda$ is perfect at $q$ up to a unit; $M$ and $Y$ come
  from \cref{pf:ks2:min}(1),(4), \'etaleness of $\mathfrak a$ from (5), and projectivity from (2)
  and the projectivity of Kuga--Sato powers.
\item\label{pf:ks2:S:num} (numbering, \lref{UpperTheory.Gal.Int.Geo.CovMain.numEx_of}) If
  $N\in\mathfrak n$, there are $c\in\N$ and $a\in\kstwoOF$ with $N=N_0c$ and
  $(a,N)\mathfrak a=(N)$; this property is stable under
  replacing $(N,c,a)$ by $(Nt,ct,ta)$.
\item\label{pf:ks2:S:cls} (class maps, \lref{UpperTheory.Gal.Int.Geo.CovMain.clsEx_of}) For
  such $(N,c,a)$ and every generic object $o$ of level $N$ over an $E$-scheme $X$ there is a
  \emph{class map}: a morphism $j\colon X\to Y$ over $\Spec O_w$ together with a morphism $\gamma$ of
  abelian schemes over $j$ which identifies $A_o$, with its $O$-action and polarization, with the
  pullback of the universal abelian scheme of $Y$, and carries $\eta_o(cy)$ to the level-$N_0$
  structure and $\eta_o(ay)$ to the \'etale level sections of $Y$ ($y\in\Lambda$); it is unique.
\item\label{pf:ks2:S:cover} (covers, \lref{UpperTheory.Gal.Int.Geo.CovMain.ksCover_holds})
  Let $K_0$ be the level group of $\mathfrak n$, $K_N\subseteq K_0$, $C$ the class map of the
  universal generic object of $M_N$, and assume the deck condition for $K_0$. Then the $k$-th fibre
  power of $C$ composed with $Z_N\cong\kstwoKS_N^k$, $\pi_C\colon Z_N\to\kstwoKS_Y^k$, is the map of a
  cover of $(T.\mathrm{tower}(N,g),\theta)$ over $O_w$ with $\Delta=K_0/K_N$.
\item\label{pf:ks2:S:lies} (lying over $T_i$, \lref{UpperTheory.Gal.Int.Geo.CovMain.tjLies_of})
  In the situation of \ref{pf:ks2:S:cover}, let $j\in J$ and $i\le n$ with $K_{g_j}=K_0$, $h_{g_j}$ a
  $T_i$-element at $\mathfrak w$, $\Nm\mathfrak w^{\,i}\mid M_{g_j}$, and $[K_0:\Gamma_j]=\#\mathrm{Gr}_i
  ((\kstwoOF/\mathfrak w)^n)$ with $\Gamma_j=K_0\cap h_{g_j}^{-1}K_0h_{g_j}$. Then for every
  $T_i$-correspondence $T_i$ of degree $q^{fi}$ on $M$ and every lift $L_i$ to $Y$, the index
  $m_j=[\Gamma_j:\Gamma_j\cap K_{M_{g_j}}]$ is positive and the correspondence $j$ of
  $T.\mathrm{tower}(N,g)$ lies over the $k$-th Kuga--Sato power of $L_i$ with multiplicity $m_j$.
\end{enumerate}

\subsubsection{Congruences}

\begin{theorem}\label{pf:ks2:kscong}
$\sfname{KSCong}$ holds. That is, in the situation above, if the Hecke data contain the operators
$T_i$ at $\mathfrak w$ with eigenvalues $\vartheta$, i.e.\ $\mathcal{H}(g,\theta,\mathfrak w,\vartheta)$,
then $T.\mathrm{tower}(N,g)$ satisfies the congruence relation at $w$ with coefficients $\vartheta$,
i.e.\ $\mathcal{C}(T.\mathrm{tower}(N,g),\theta,w,n,\vartheta)$.
\formallabel{UpperTheory.Gal.Int.Geo.CovMain.ksCong_holds}
\end{theorem}

\begin{proof}
By $\mathcal{H}$ choose $q,r$ with $\mathrm{Adm}(q,r,\mathfrak w)$, an auxiliary ideal
$\mathfrak n$ with level group $K_0\supseteq K_N$ and $N\in\mathfrak n$, such that the deck condition
holds for $K_0$, and for every $i\le n$ an index $\gamma_i\in J$ with the properties listed in
\cref{pf:ks2:def:goodgrp}(6). By \ref{pf:ks2:S:site} choose a site at $(w,q,r,\mathfrak w,\mathfrak n)$,
with $f$, $b$, $M$, $Y$, $\pi$. By \ref{pf:ks2:S:num} write $N=N_0c$ with $a$ as there, by
\ref{pf:ks2:S:cls} take the class map $C$ of the universal generic object of $M_N$, and by
\ref{pf:ks2:S:cover} the cover $(\kstwoKS_Y^k,\pi_C)$ of $(T.\mathrm{tower}(N,g),\theta)$ over $O_w$.
Choose a uniformizer $\varpi$ of $O_w$. The congruence relation \cref{pf:ks2:congks}, applied to the
datum $D_q$ for the prime $q$, the discrete valuation ring $O_w$ (of characteristic $0$, with $q$ in
its maximal ideal), the base $b$, the model $M$, the ideal $\mathfrak a$ and the flat (indeed
\'etale) morphism $\pi$, gives $T_i$-correspondences $T_i$ of degrees $q^{fi}$ on $M$, lifts $L_i$
to $Y$ and integers $d_i$ such that the $k$-th Kuga--Sato powers $L_i^{(k)}$ are correspondences of
$\kstwoKS_Y^k$ away from $q$, both of whose legs have the same composite to $\Spec O_w$, satisfying
$\mathrm{PC}(q,q^f,n,\varpi,(L_i^{(k)})_i)$. For each $i$, \ref{pf:ks2:S:lies} applied to $\gamma_i$,
$T_i$ and $L_i$ shows that $\gamma_i$ lies over $L_i^{(k)}$ with multiplicity
$m_i=[\Gamma_{\gamma_i}:\Gamma_{\gamma_i}\cap K_{M_{g_{\gamma_i}}}]>0$, and $\mathcal{H}$ gives
$\theta(\gamma_i)=m_i\vartheta_i$. Together with $q\in w$ and $\Nm w=q^f$ this is
$\mathcal{C}(T.\mathrm{tower}(N,g),\theta,w,n,\vartheta)$
(\cref{pf:ks2:def:goodatw}(2)).
\end{proof}

\subsubsection{Good reduction}

For good reduction one needs, besides the cover at the level $N_0\mathfrak a$, maps between the
Kuga--Sato powers of \'etale-level models at two auxiliary levels which extend the level-forgetting
maps and the Hecke translations of the generic fibre. The formal proof uses the following
statements, all proved in the development.
\begin{enumerate}[label=(G\arabic*)]
\item\label{pf:ks2:G:ext} (extension data, \lref{UpperTheory.Gal.FiniteEtaleModels.extEx}) For a
  site and an auxiliary level $(N_1,\mathfrak b)$ (\cref{pf:ks2:def:goodgrp}(1)),
  an \emph{extension datum} exists: a model of level $N_1$ and an
  \'etale-level model of level $(N_1,\mathfrak b,1)$ for the base $b$ of the site, with a finite
  \'etale morphism between them compatible with the universal objects, such that all Kuga--Sato
  powers of the \'etale-level model are projective and smooth over $O_w$.
\item\label{pf:ks2:G:quot} (quotients, \lref{UpperTheory.Gal.FiniteEtaleModels.quotOf}) The map
  $\pi_C\colon Z_M\to\kstwoKS_Y^k$ attached to a class map at level $M$ exhibits $\kstwoKS_Y^k$ as the
  quotient of the tower at level $M$ by $K_0$: it
  lies over $\Spec E\to\Spec O_w$, $K_M\subseteq K_0$ with finite index, $\pi_C$ is invariant under
  the action of $K_0$ on $Z_M$, two geometric points with the same image differ by an element of
  $K_0$, an element of $K_0$ fixing a geometric point lies in $K_M$, and every geometric point of
  $\kstwoKS_Y^k$ over the generic point of $\Spec O_w$ is in the image.
\item\label{pf:ks2:G:steps} (extension of witnesses; \cref{pf:ks2:stepf} below,
  \lref{UpperTheory.Gal.FiniteEtaleModels.stepA},
  \lref{UpperTheory.Gal.FiniteEtaleModels.stepB}) For extension data $X_1,X_2$ at auxiliary
  levels $N_0\kappa\mathfrak b_1$ and $N_0\mathfrak b_2$ with $\mathfrak b_1\subseteq\mathfrak b_2$, and
  class maps at levels $M_1,M_2$: a Hecke witness for $1$ from $M_{M_1}$ to $M_{M_2}$ extends to a
  finite \'etale morphism between the \'etale-level models of $X_1$ and $X_2$ with cartesian square of
  universal abelian schemes (case $h=1$); a Hecke witness for an integral $h$ with
  $\kappa\kstwoLaHat\subseteq h\kstwoLaHat$ extends to a morphism between them
  (prime-to-$q$ case); and, when $\mathfrak b_1\subseteq\mathfrak b_2\mathfrak w'$ for a prime
  $\mathfrak w'\supseteq\mathfrak u$ and $h$ is a $T_i$-element at $\mathfrak w'$ with
  $\Nm\mathfrak w'^{\,i}\mid M_1$, a Hecke witness for $h$ extends as well ($T_i$ case).
\item\label{pf:ks2:G:deg} (degrees, \lref{UpperTheory.Gal.FiniteEtaleModels.degGen_holds},
  \lref{UpperTheory.Gal.FiniteEtaleModels.fEtExt_holds}) If $\pi'\circ s$ exhibits
  $Y$ as the quotient by $K_0$ and $\pi'$ exhibits $Y'$ as the quotient by $K_0'\subseteq K_0$, then
  $s$ has $[K_0:K_0']$ geometric points over every geometric point of the generic fibre; and a finite
  \'etale morphism to a flat $O$-scheme with constant degree $d$ over the generic fibre is finite
  \'etale of degree $d$ (over $O$, $O$ a domain).
\item\label{pf:ks2:G:lies} (a lying-over criterion, \lref{UpperTheory.Gal.FiniteEtaleModels.liesAbs_holds})
  Let $(Y,\pi)$ be a cover of $(T.\mathrm{tower}(N,g),\theta)$ over $O$, $j\in J$, $K_0'\subseteq K_0$
  with $K_N\subseteq K_0$, and let $\pi'\colon Z_{M_{g_j}}\to Y'$ and a finite \'etale $s\colon Y'\to Y$
  of some degree be such that $\pi'\circ s$ and $\pi'$ exhibit $Y$ and $Y'$ as quotients by $K_0$
  and $K_0'$, $\pi\circ\mathrm{tr}(M_{g_j},N,1)=s\circ\pi'$, every $u\in K_0$ can be written $u=xu'$
  with $x\in K_{g_j}\cap K_0$ and $u'\in K_0'$, and for every $x\in K_{g_j}$ there is $t\colon Y'\to Y$
  over $O$ with $t\circ\pi'=\pi\circ\mathrm{tr}(M_{g_j},N,h_{g_j}x)$. Then $j$ lies over some
  correspondence $(s',t')$ of $Y$ over $O$ with $s'$ finite \'etale.
\end{enumerate}

The case $h=1$ of \ref{pf:ks2:G:steps} is the following proposition.

\begin{proposition}[Extension of level-forgetting witnesses]\label{pf:ks2:stepf}
Let a site at $(w,q,r,\mathfrak w,\mathfrak n)$ with level $N_0$ be given, let $\kappa\in\N$, and let
$X_1$, $X_2$ be extension data (\ref{pf:ks2:G:ext}) at the auxiliary levels $(N_0\kappa,\mathfrak
b_1)$ and $(N_0,\mathfrak b_2)$ with $\mathfrak b_1\subseteq\mathfrak b_2$; write $Y_1,Y_2$ for their
\'etale-level models, with structure maps $s_1,s_2$ to $\Spec O_w$ and universal objects $\mathcal
B_1,\mathcal B_2$. Let $M_1,M_2$ be levels with $M_2\mid M_1$, $(c_1,a_1)$ and $(c_2,a_2)$ numberings
of $M_1$ and $M_2$ with respect to $X_1$ and $X_2$ (\ref{pf:ks2:S:num}), $(j_1,\gamma_1)$ and
$(j_2,\gamma_2)$ the class maps of the universal generic objects of $M_{M_1}$ and $M_{M_2}$ into $Y_1$
and $Y_2$ with these numberings (\ref{pf:ks2:S:cls}), and $(\pi,\Phi)$ a Hecke witness for $1$ from
$M_{M_1}$ to $M_{M_2}$, with $\Phi=\varphi\circ\psi$ for a Hecke isogeny $\psi$ and a morphism of
objects $\varphi$. Then there are a finite \'etale morphism $d\colon Y_1\to Y_2$ with $s_2\circ
d=s_1$ and a morphism $\delta$ from the abelian scheme of $\mathcal B_1$ to that of $\mathcal B_2$ over
$d$, with cartesian square, such that $d\circ j_1=j_2\circ\pi$ and
$\delta\circ\gamma_1=\gamma_2\circ\varphi\circ\psi$.
\formallabel{UpperTheory.Gal.FiniteEtaleModels.O.stepF}
\end{proposition}

\begin{proof}
By \ref{pf:ks2:S:num}, $M_1=N_0\kappa c_1$ and $M_2=N_0c_2$ with $(a_1,M_1)\mathfrak b_1=(M_1)$ and
$(a_2,M_2)\mathfrak b_2=(M_2)$; write $M_1=M_2t$ with $t\in\N$. Let $\eta$ denote level structures.
\begin{enumerate}
\item By \lref{UpperTheory.Gal.FiniteEtaleModels.O.epsEx} (elementary ideal arithmetic in the
  Dedekind domain $\kstwoOF$) there is $\varepsilon\in\kstwoOF$ with $M_1\mid\varepsilon a_1-ta_2$ and
  $((\varepsilon)+\mathfrak b_1)\mathfrak b_2=\mathfrak b_1$.
\item By \lref{UpperTheory.Gal.FiniteEtaleModels.O.downEx} there are a finite \'etale morphism
  $d\colon Y_1\to Y_2$ with $s_2\circ d=s_1$ and a morphism of objects $\delta$, over $d$, from
  $\mathcal B_1$ with its level structure restricted to level $N_0$ to $\mathcal B_2$, which carries
  the \'etale level sections $\lambda_1(\varepsilon y)$ of $Y_1$ to the \'etale level sections
  $\lambda_2(y)$ of $Y_2$ ($y\in\Lambda$). (By the identity in (1), multiplication by $\varepsilon$
  turns an \'etale level structure of type $(\mathfrak b_1,1)$ into one of type $(\mathfrak b_2,1)$, and
  $Y_2$ classifies the resulting pair; finiteness and \'etaleness come from the finite \'etale maps
  from $Y_1,Y_2$ to the models of the extension data and from the level-forgetting map between these
  models, which is finite, \'etale and surjective,
  \lref{GlobalTheory.T.harrisTaylor_levelForgetting_finite_etale_surjective}.)
\item Composing $(j_1,\gamma_1)$ with $(d,\delta)$ gives a pair $(d\circ j_1,\delta\circ\gamma_1)$ of
  the type of a class map of the universal generic object of $M_{M_1}$ into $Y_2$ with the numbering
  $(c_1\kappa,ta_2)$. Indeed, $\gamma_1$ carries $\eta(c_1\kappa y)$ to the level-$N_0$ structure of
  $\mathcal B_1$, which $\delta$ carries to that of $\mathcal B_2$; and $\eta(ta_2y)=\eta(\varepsilon
  a_1y)$ because $M_1\mid\varepsilon a_1-ta_2$ and $\eta$ kills $M_1\Lambda$, while $\gamma_1$ carries
  $\eta(a_1\varepsilon y)$ to $\lambda_1(\varepsilon y)$, which $\delta$ carries to $\lambda_2(y)$.
\item Composing the witness with $(j_2,\gamma_2)$ gives a second such pair $(j_2\circ\pi,
  \gamma_2\circ\varphi\circ\psi)$: the Hecke isogeny $\psi$ for $1$ is an isomorphism compatible with
  the polarizations (\lref{UpperTheory.trace_unit_of_pullback_multiplier}); it carries
  $\eta(c_1\kappa y)$ to $\eta(c_2y)$ and $\eta(ta_2y)$ to $\eta(a_2y)$
  (\lref{UpperTheory.trace_eta_identity}: a Hecke isogeny for $h$ from level $N'u$ to level
  $N'u'$ carries $\eta(uy)$ to $\eta(u'y')$ whenever $h\cdot y\equiv y'\bmod N'\kstwoLaHat$); and
  $\gamma_2\circ\varphi$ carries these to the level-$N_0$ structure of $\mathcal B_2$ and to
  $\lambda_2(y)$.
\item The numbering $(c_1\kappa,ta_2)$ of $M_1$ with respect to $X_2$ is admissible:
  $M_1=N_0c_1\kappa$, and $(ta_2,M_1)\mathfrak b_2=t\,(a_2,M_2)\mathfrak b_2=(tM_2)=(M_1)$.
\item By the uniqueness of class maps (\ref{pf:ks2:S:cls}) the two pairs coincide: $d\circ
  j_1=j_2\circ\pi$ and $\delta\circ\gamma_1=\gamma_2\circ\varphi\circ\psi$.
\end{enumerate}
The morphism $\delta$ lies over $d$ with cartesian square, being a morphism of objects, and $d$ is
finite and \'etale; together with (6) this is the assertion.
\end{proof}

\begin{theorem}\label{pf:ks2:ksgood}
$\sfname{KSGood}$ holds. That is, in the situation above, if the Hecke data are unramified at
$\mathfrak w$, i.e.\ $\mathcal{U}(g,\theta,\mathfrak w)$, then $T.\mathrm{tower}(N,g)$ has good
reduction at $w$, i.e.\ $\mathcal{G}(T.\mathrm{tower}(N,g),\theta,w)$.
\formallabel{UpperTheory.Gal.FiniteEtaleModels.ksGood_holds}
\end{theorem}

\begin{proof}
By $\mathcal{U}$ choose $q,r,N_0,\mathfrak a,K_0$ as in \cref{pf:ks2:def:goodgrp}(5). By
\ref{pf:ks2:S:site} (with $\mathfrak n=N_0\mathfrak a$) there is a site, by \ref{pf:ks2:G:ext} an
extension datum $X_0$ at the level $N_0\mathfrak a$, by \ref{pf:ks2:S:num} a decomposition
$N=N_0c_0$ with $a_0$, by \ref{pf:ks2:S:cls} a class map $C_0$ for $M_N$, and by \ref{pf:ks2:S:cover}
a cover $(Y,\pi)=(\kstwoKS_{X_0}^k,\pi_{C_0})$ of $(T.\mathrm{tower}(N,g),\theta)$ over $O_w$. We show that
it is good. Let $j\in J$; $\mathcal{U}$ provides an integer $\kappa$, an auxiliary level
$N_0\kappa\mathfrak a'$ with $\mathfrak a'\subseteq\mathfrak a$ and level group $K_0'\supseteq
K_{M_{g_j}}$, $M_{g_j}\in N_0\kappa\mathfrak a'$, a decomposition $K_0=(K_{g_j}\cap K_0)K_0'$, and a
factorization $h_{g_j}=h_1h_2$ with the properties of \cref{pf:ks2:def:goodgrp}(5). By
\ref{pf:ks2:G:ext}, \ref{pf:ks2:S:num} and \ref{pf:ks2:S:cls} we get an extension datum $X_1$ at
$N_0\kappa\mathfrak a'$ and a class map $C_1$ for $M_{M_{g_j}}$; put $Y'=\kstwoKS_{X_1}^k$ and
$\pi'=\pi_{C_1}$.
\begin{enumerate}[label=(\alph*)]
\item (Level forgetting.) By \ref{pf:ks2:G:steps} (case $h=1$), the Hecke witness for $1$ from
  $M_{M_{g_j}}$ to $M_N$ extends to a finite \'etale morphism of \'etale-level models with cartesian
  square of universal abelian schemes; its $k$-th fibre power $s\colon Y'\to Y$ satisfies
  $\pi\circ\mathrm{tr}(M_{g_j},N,1)=s\circ\pi'$ (compatibility of class maps with witnesses,
  and of fibre powers with composition and cartesian squares). By \ref{pf:ks2:G:quot} $\pi'\circ s$ and $\pi'$ exhibit $Y$
  and $Y'$ as quotients by $K_0$ and $K_0'$, so by \ref{pf:ks2:G:deg} $s$ is finite \'etale of degree
  $[K_0:K_0']$ over $O_w$ ($Y$ is flat over $O_w$, being smooth).
\item (Translations.) For $x\in K_{g_j}$, the element $h_{g_j}x=h_1(h_2x)$ is integral; by
  \ref{pf:ks2:G:steps} (prime-to-$q$ case for the factor $h_2x$, which satisfies
  $\kappa\kstwoLaHat\subseteq h_2x\kstwoLaHat$ because $x\in K_{g_j}$ preserves $\kstwoLaHat$, and $T_i$
  case for $h_1$ when it is a $T_i$-element,
  nothing being needed when $h_1=1$), the Hecke witness for $h_{g_j}x$ extends to a morphism of
  Kuga--Sato powers $t\colon Y'\to Y$ over $O_w$ with $t\circ\pi'=\pi\circ\mathrm{tr}(M_{g_j},N,h_{g_j}x)$.
\end{enumerate}
Now \ref{pf:ks2:G:lies} applies and shows that $j$ lies over a correspondence $(s',t')$ on $Y$ over
$O_w$ with $s'$ finite \'etale. As $j$ was arbitrary, the cover is good, i.e.\
$\mathcal{G}(T.\mathrm{tower}(N,g),\theta,w)$.
\end{proof}

% ---------------------------------------------------------------------------
\subsection{Fixed points of twisted Hecke correspondences}\label{pf:ks2:sec:fix}

Throughout this subsection the ground field is $\C$, $D,G,L$ satisfy $\mathrm{Sc}(D,G)$ (so $a=m=1$
and there is a unique embedding $t_0=\tau_1\colon F\to\C$), the base $g_b$ over $\C$ is standard,
and $v$ is the archimedean place of $F$ at which the Gram matrix $H_v$ has exactly one positive
eigenvalue.

\subsubsection{Tangent data}

\begin{definition}[Tangent matrices]\label{pf:ks2:def:istan}
Let $s\colon Z\to\Spec\C$, $s'\colon Y\to\Spec\C$, $c_1,c_2\colon Y\to Z$ with $sc_1=sc_2=s'$, let
$y\in Y(\C)$ with $c_1(y)=c_2(y)$, $\delta\in\N$ and $A\in M_\delta(\C)$. Then
$\mathrm{Tan}(\delta,c_1,c_2,y,A)$ means: there are affine opens
$U\subseteq Z$ and $V\subseteq c_1^{-1}U\cap c_2^{-1}U$ with $y\in V$, functions
$t_1,\dots,t_\delta\in\Gamma(U,\OO_Z)$ and a $\C$-linear map
$\varphi\colon\Omega_{\Gamma(V,\OO_Y)/\C}\to\C^\delta$ with $\varphi(b\eta)=b(y)\varphi(\eta)$ and
$\ker\varphi=\mathfrak m_y\Omega_{\Gamma(V,\OO_Y)/\C}$, such that $\varphi(d\,c_1^*t_l)=e_l$ and
$\varphi(d\,c_2^*t_j)=(A_{jl})_l$ for all $j,l$. The \emph{Lefschetz term} of $A$ is
$\mathrm{lef}(A)(t)=\det(1-tA)/\det(1-A)$.
\formal{GlobalTheory.W1.IsTan}\formal{UpperTheory.Gal.W2.lefTerm}
\end{definition}

Thus $c_1^*t_1,\dots,c_1^*t_\delta$ are local coordinates at $y$ and $A$ is the cotangent matrix of
the local self-map ``$c_2\circ c_1^{-1}$'' at $y$.

\subsubsection{Forms, points and orbital sums}

\begin{definition}[Forms and their groups]\label{pf:ks2:def:forms}
Let $x\mapsto x^*$ be the involution of $C$.
\begin{enumerate}
\item A \emph{form} is $\beta\in C^\times$ with $\beta^*=\beta$ such
  that $x\beta x^*=t$ for some $x\in(\A_f^p\otimes C)^\times$ and $t\in(\A_f^p)^\times$, and
  $z\beta z^*=1$ for some $z\in(\R\otimes C)^\times$. Two forms are
  \emph{equivalent} if $\beta'=s\,y\beta y^*$ with $y\in C^\times$ and $s\in\Q_{>0}$; $\mathfrak{C}$
  denotes the set of classes. We fix a representative $\beta$ of
  every class and a trivialization $(x,t)$ for it.
\item $\Gamma_\beta=\{(y,s)\in C^\times\times\Q^\times:\ y\beta y^*=s\beta,\ s>0\}$. It maps to $\kstwoG$ by $y\mapsto xyx^{-1}$
  (denoted $y\mapsto y_f$), to $\kstwoGp$ by localization at $p$
  ($y\mapsto y_p$), and to $\GL_n(\C)$ by $\psi_v$
  ($y\mapsto y_v$).
\item $\mathcal{P}_\beta$ is the set of $z\in\C^n$ with $\Re(z^*H_v\psi_v(\beta^{-1})z)>0$, and $\mathcal{D}_\beta=(\kstwoG\times\kstwoGp)\times\mathcal{P}_\beta$.
  For subgroups $K\subseteq\kstwoG$ and $K'\subseteq\kstwoGp$, $\mathcal{X}_\beta(K,K')$ is the quotient of $\mathcal{D}_\beta$ by the equivalence
  relation $((g,g_p),z)\sim((y_fgk,y_pg_pk'),c\,y_vz)$ ($y\in\Gamma_\beta$, $k\in K$, $k'\in K'$,
  $c\in\C^\times$), and
  $\mathcal{X}(K,K')=\coprod_{[\beta]\in\mathfrak{C}}\mathcal{X}_\beta(K,K')$.
\item For $a=((g,g_p),z)\in\mathcal{D}_\beta$ and $y\in\Gamma_\beta$, $y$ \emph{fixes the line of $a$}
  if $y_vz=\lambda z$ for some $\lambda\in\C$; then
  $\lambda=\mathrm{eig}(y,a)$, and $\lambda\ne0$. For a
  subset $S\subseteq\kstwoG$, $y$ \emph{sees $a$ through $S$} if $y$
  fixes the line of $a$, $g^{-1}y_fg\in S$ and $g_p^{-1}y_pg_p\in\kstwoKp$. The stabilizer
  $\mathrm{Stab}(a)$ is the set of $y$ seeing $a$ through $K$.
\item $g\in\kstwoG$ is \emph{regular} if at some place $w\ne p$
  the centralizer of the $w$-component of $g$ in $\Q_w\otimes C$ is commutative.
\end{enumerate}
\formal{GlobalTheory.LKCx.Form}\formal{GlobalTheory.LKMid.Triv}\formal{GlobalTheory.LKCx.FCls}\formal{GlobalTheory.LKCx.grp}\formal{GlobalTheory.LKCx.toAfF}\formal{GlobalTheory.LKCx.toP1}\formal{GlobalTheory.LKCx.toV}\formal{GlobalTheory.LKCx.PosV}\formal{GlobalTheory.LKCx.Pts}\formal{GlobalTheory.LKCx.Same}\formal{GlobalTheory.LKCx.AllPts}\formal{GlobalTheory.LKCx.FixL}\formal{GlobalTheory.RootSumWeights.eigL}\formal{UpperTheory.Gal.W2.Sees}\formal{GlobalTheory.LKCx.Stab}\formal{GlobalTheory.LKCx.RegAf}
\end{definition}

\begin{definition}[Weights]\label{pf:ks2:def:weights}
For a polynomial $\chi\in F[X]$ and $\tau\colon F\to\C$ let $R_\tau(\chi)$ be the multiset of roots
of $\tau(\chi)$. For $\varepsilon\colon\Hom(F,\C)\to\{0,1\}$,
$x\in F$, $y,\mu\in\C$ put, with $R=R_{t_0}(\chi)\smallsetminus\{\mu\}$ (one copy of $\mu$
removed) and bars denoting complex conjugation,
\begin{align*}
\kstwoBall(y,\chi,\mu)&=\prod_{r\in R}\frac{r-y\mu}{r-\mu},\\
\kstwoFib(x;y,\chi,\mu)&=\Bigl[(1-t_0(x)y\mu)\prod_{r\in R}(1-t_0(x)r)\Bigr]
\Bigl[(1-\overline{t_0(x)}\,\bar\mu)\prod_{r\in R}(1-\overline{t_0(x)}\,y\,\bar r)\Bigr]\\
&\qquad\times\prod_{\tau\ne t_0,\bar t_0}\ \prod_{r\in R_\tau(\chi)}(1-\tau(x)y^{\varepsilon(\tau)}r),
\end{align*}
and, for $\alpha\in\kstwoOF^k$,
$\kstwoloc(\alpha;y,\chi,\mu)=\kstwoBall(y,\chi,\mu)\prod_{s=1}^k\kstwoFib(\alpha_s;y,\chi,\mu)$.
Let $\varepsilon_D(\tau)=0$ if $\tau(\omega)=\rho_c$ and $1$ otherwise,
and for $P\in F[X]$ with $P(0)\ne0$ let $P^\vee=P(0)^{-1}X^{\deg P}P(1/X)$ (the reversed polynomial). The \emph{fixed-point weight} is
\begin{equation}\label{pf:ks2:eq:cw}
c_W(\alpha,y)(P,\lambda)=\kstwoloc\bigl(\alpha;\,y,\,P^\vee,\,\kappa(\lambda)^{-1}\bigr),\qquad
\kappa(\lambda)=\begin{cases}\lambda&\text{if }v\text{ induces }t_0\text{ (via }L),\\ \bar\lambda&\text{otherwise,}\end{cases}
\end{equation}
where $\varepsilon=\varepsilon_D$ (it equals the weight of the statement of $\sfname{FixK}$ by
\lref{UpperTheory.Gal.FixZ.wtEq_of}, a direct comparison of definitions using that the two
roots of $X^2-\mathrm{tr}X+\mathrm{nm}$ are $\rho_c\ne\bar\rho_c$).
\formal{UpperTheory.Gal.W2.rts}\formal{UpperTheory.Gal.W2.ballW}\formal{UpperTheory.Gal.W2.fibW}\formal{UpperTheory.Gal.W2.locW}\formal{UpperTheory.Gal.W2.epsD}\formal{UpperTheory.Gal.W2.invPoly}\formal{UpperTheory.Gal.W2.wtFix}\formal{UpperTheory.Gal.LocalSmoothExtension.cW}
\end{definition}

\begin{definition}[Orbital sums]\label{pf:ks2:def:orbital}
For $c\colon F[X]\times\C\to\C$, $K\subseteq\kstwoG$ and $f\colon\kstwoG\to\C$ put
\begin{multline}\label{pf:ks2:eq:cxtotal}
\Sigma(c;K,f)=\sum_{[\beta]\in\mathfrak{C}}\ \sum_{q\in\mathcal{X}_\beta(K,\kstwoKp)}
\frac1{\#\mathrm{Stab}(a_q)}\\
\times\sum_{\substack{y\in\Gamma_\beta\\ y\text{ fixes the line of }a_q}}
f(g_q^{-1}y_fg_q)\,\one_{\kstwoKp}(g_{p,q}^{-1}y_pg_{p,q})\,c\bigl(\mathrm{minpoly}_F(y),\mathrm{eig}(y,a_q)\bigr),
\end{multline}
where $a_q=((g_q,g_{p,q}),z_q)$ is a chosen representative of $q$, $\mathrm{minpoly}_F(y)$ is the
minimal polynomial of $y\in C$ over $F$, and all sums are \emph{finite sums in Mathlib's sense}:
a sum over an infinite support is $0$. $\mathrm{Fin}(K,f)$ means that the set of triples $([\beta],q,y)$
with $y$ fixing the line of $a_q$, $f(g_q^{-1}y_fg_q)\ne0$ and $g_{p,q}^{-1}y_pg_{p,q}\in\kstwoKp$ is
finite.
\formal{GlobalTheory.RootSumWeights.cxTotalM}\formal{GlobalTheory.RootSumWeights.cxSumAtM}\formal{UpperTheory.Gal.LocalSmoothExtension.CxFin}
\end{definition}

\subsubsection{The statement}

For a Kuga--Sato tower $T$ over $\C$ and $\alpha\in\kstwoOF^k$ let
$[\alpha]\colon Z_N\to Z_N$ be the endomorphism acting on the $s$-th factor of $\kstwoKS_N^k$ by
$\alpha_s\in\kstwoOF\subseteq O$; for $N\mid N'$
and $h$ admissible from $N'$ to $N$ put $c_1=\mathrm{tr}(N',N,1)$ and
$c_2=[\alpha]\circ\mathrm{tr}(N',N,h)$.
Let $\Gamma_h=K_N\cap h^{-1}K_Nh$.

\begin{definition}[$\sfname{FixK}$]\label{pf:ks2:def:fixk}
$\sfname{FixK}$ is the statement: for every
Kuga--Sato tower $T$ over $\C$ for $(L,g_b,k)$ (with $\mathrm{Sc}(D,G)$, $g_b$ standard, $v$ as
above), levels $N\mid N'$, $\delta\in\N$ such that $Z_N\to\Spec\C$ is smooth of relative dimension
$\delta$, $h\in\kstwoG$ integral and admissible from $N'$ to $N$ (with $1$ admissible from $N'$ to $N$)
such that every element of $K_Nh^{-1}K_N$ is regular, and every $\alpha\in\kstwoOF^k$ with all
$\alpha_s$ nonzero and divisible by $2$:
\begin{enumerate}
\item $\mathrm{Fin}(K_N,\one_{K_Nh^{-1}K_N})$;
\item $c_2$ is finite and \'etale;
\item there are a finite set $\mathcal F\subseteq Z_{N'}(\C)$, equal to the set of $x$ with
  $c_1(x)=c_2(x)$, and matrices $A_x\in M_\delta(\C)$ with $\mathrm{Tan}(\delta,c_1,c_2,x,A_x)$ and
  $\det(1-A_x)\ne0$ for $x\in\mathcal F$, such that for all $y\in\C$
\begin{equation}\label{pf:ks2:eq:fixk}
\sum_{x\in\mathcal F}\frac{\det(1-yA_x)}{\det(1-A_x)}=[\Gamma_h:\Gamma_h\cap K_{N'}]\cdot
\Sigma\bigl(c_W(\alpha,y);K_N,\one_{K_Nh^{-1}K_N}\bigr).
\end{equation}
\end{enumerate}
\formal{UpperTheory.Gal.LocalSmoothExtension.FixK}
\end{definition}

\begin{theorem}\label{pf:ks2:fixk}
$\sfname{FixK}$ holds.
\formallabel{UpperTheory.Gal.FixZ.FixK_holds}
\end{theorem}

The proof occupies the rest of this subsection; it is completed after \cref{pf:ks2:fixz}.

\subsubsection{Legs and the reduction to moduli schemes}

For generic models $X'$, $X$ of levels $N'$, $N$ over $\C$, a \emph{leg} for $h$ is a morphism $c_q\colon X'\to X$ over $\C$ with a generic object
$o$ over $X'$, a Hecke isogeny $\psi$ for $h$ from the universal object of $X'$ to $o$, and a
morphism of objects $o\to\mathcal A_X$ over $c_q$; its \emph{total map} is the composite
$c_{\mathrm{tot}}\colon\mathcal A_{X'}\to\mathcal A_X$, its $k$-th power is the fibre power
$c^{(k)}\colon\mathcal A_{X'}^k\to\mathcal A_X^k$ over $c_q$, and the \emph{twisted power} is
$c^{(k)}_\alpha=[\alpha]\circ c^{(k)}$.

\begin{proposition}[Fixed points on moduli schemes]\label{pf:ks2:hfp}
Assume $a=m=1$, let $X,X'$ be generic models over $\C$ of levels $N\mid N'$ with $h$ admissible from
$N'$ to $N$, let $c_1$ be a leg for $1$ and $c_2$ a leg for $h$
from $X'$ to $X$, assume that every $g$ with $hg\in K_N$ is regular, let $\alpha\in\kstwoOF^k$ with all
$\alpha_s$ divisible by $2$, and assume that $c_1^{(k)}$ is \'etale. Then:
\begin{enumerate}
\item $\mathcal A_X^k\to\Spec\C$ is smooth of relative dimension $\delta=n-1+kg$;
\item there are a finite set $\mathcal F$ of complex points of $\mathcal A_{X'}^k$, equal to the set of
  fixed points of $(c_1^{(k)},c^{(k)}_{2,\alpha})$, and tangent matrices $A_x$ with $\det(1-A_x)\ne0$;
\item the set of $([\beta],q)\in\mathcal{X}(K_{N'},\kstwoKp)$ such that some $y$ sees $a_q$ through
  $h^{-1}K_N=\{g:hg\in K_N\}$ is finite;
\item if $y$ sees $a$ through $K_N$, then $y=1$;
\item for all $t\in\C$,
$\sum_{x\in\mathcal F}\mathrm{lef}(A_x)(t)=\Sigma\bigl(c_W(\alpha,t);K_{N'},\one_{h^{-1}K_N}\bigr)$.
\end{enumerate}
\formallabel{UpperTheory.Gal.W2.fixedPoints_int_base_holds}
\end{proposition}

The proof is given in \cref{pf:ks2:sec:fixproof}. We first change the level.

\begin{proposition}[Change of level]\label{pf:ks2:levelchange}
Let $K'\subseteq K$ be subgroups of $\kstwoG$ with $[K:K']$ finite and $K'\subseteq\Gamma_h=K\cap
h^{-1}Kh$, let $c$ be arbitrary, and assume: no $y\ne1$ sees any $a$ through $K$, and the set of
$([\beta],q)\in\mathcal{X}(K',\kstwoKp)$ such that some $y$ sees $a_q$ through $h^{-1}K$ is finite.
Then
\[
\Sigma(c;K',\one_{h^{-1}K})=[\Gamma_h:K']\cdot\Sigma(c;K,\one_{Kh^{-1}K}).
\]
\formallabel{UpperTheory.Gal.W2.LevelChange_holds}
\end{proposition}

\begin{proof}
It suffices to prove the identity for each class $[\beta]$ separately.
Fix $\beta$ and write $\Gamma=\Gamma_\beta$. We use the following invariance
(\lref{UpperTheory.Gal.W2.EigSame_holds}, a consequence of the definitions): if $a'=\gamma a k$
with $\gamma\in\Gamma$ and $k\in K''\times\kstwoKp$ (and a scalar on the $\C^n$-component), then
$y\mapsto\gamma y\gamma^{-1}$ is a bijection carrying the elements seeing $a$ through an
$K''$-conjugation invariant set $S$ to those seeing $a'$ through $S$, and preserving
$\mathrm{minpoly}_F(y)$ and $\mathrm{eig}(y,a)$. The sets $h^{-1}K$ and $Kh^{-1}K$ are invariant under
conjugation by $K'$, resp.\ $K$: for $u\in K'\subseteq K\cap h^{-1}Kh$, $u^{-1}gu\in h^{-1}K$ iff
$(huh^{-1})^{-1}(hg)u\in K$ iff $hg\in K$. So for $a\in\mathcal{D}_\beta$ and $S\in\{h^{-1}K,Kh^{-1}K\}$
the quantity
\[
\Phi_S(a)=\sum_{y\text{ sees }a\text{ through }S}c\bigl(\mathrm{minpoly}_F(y),\mathrm{eig}(y,a)\bigr)
\]
depends only on the class of $a$ at level $K'$, resp.\ $K$.

\emph{Stabilizers.} By hypothesis every stabilizer at level $K$ is trivial; since $K'\subseteq K$,
stabilizers at level $K'$ are trivial too. Hence all factors $1/\#\mathrm{Stab}$ in
\eqref{pf:ks2:eq:cxtotal} equal $1$.

\emph{Fibres of $\mathcal{X}_\beta(K',\kstwoKp)\to\mathcal{X}_\beta(K,\kstwoKp)$.} The fibre over the
class of $a=((g,g_p),z)$ consists of the classes of $a_k=((gk,g_p),z)$, $k\in K$, and $a_k\sim a_{k'}$
at level $K'$ iff $k\in k'K'$: if $(\gamma_fgk,\gamma_pg_p,c\gamma_vz)=(gk'u,g_pu_p,z)$ with $u\in K'$,
$u_p\in\kstwoKp$, then $\gamma$ sees $a$ through $K$, so $\gamma=1$ and $k=k'u$. Thus the fibre is in
bijection with $K/K'$.

\emph{Unfolding.} For $a=((g,g_p),z)$ and $y$ fixing the line of $a$ with $g_p^{-1}y_pg_p\in\kstwoKp$, put
$u=g^{-1}y_fg$. Then $y$ sees $a_k$ through $h^{-1}K$ iff $u\in kh^{-1}Kk^{-1}=kh^{-1}K$. The cosets
$kh^{-1}K$ ($k\in K$) depend only on $k\Gamma_h$, are pairwise distinct for distinct $k\Gamma_h$
($kh^{-1}K=k'h^{-1}K$ iff $k'^{-1}k\in h^{-1}Kh$), and their union is $Kh^{-1}K$. Hence
\[
\#\{k\in K/K':\ u\in kh^{-1}K\}=[\Gamma_h:K']\cdot\#\{k\in K/\Gamma_h:\ u\in kh^{-1}K\}
=[\Gamma_h:K']\cdot\one_{Kh^{-1}K}(u).
\]
Summing $c(\mathrm{minpoly}_F(y),\mathrm{eig}(y,a))$ over $y$ (note that $\mathrm{eig}(y,a_k)=\mathrm{eig}(y,a)$
and the line and $p$-conditions do not involve $k$) gives
$\sum_{k\in K/K'}\Phi_{h^{-1}K}(a_k)=[\Gamma_h:K']\,\Phi_{Kh^{-1}K}(a)$, and summing over the classes
of level $K$ gives the identity, provided all sums involved are finite. For a given $a_k$ at most one
$y$ sees it through $h^{-1}K$: if $y,y'$ do, then $y^{-1}y'$ sees $a_k$ through
$(h^{-1}K)^{-1}h^{-1}K=K$ hence $y=y'$. Therefore the inner sums
on the left are finite, the outer sum on the left has finite support by hypothesis, and the identity
just proved shows that the supports on the right are finite as well; with finite supports the
finite sums of Mathlib are ordinary sums and the rearrangements above are legitimate.
\end{proof}

\begin{corollary}\label{pf:ks2:fixz}
In the situation of \cref{pf:ks2:hfp}, the conclusions (1) and (2) hold, and for all $t\in\C$
\[
\sum_{x\in\mathcal F}\mathrm{lef}(A_x)(t)=[\Gamma_h:K_{N'}]\cdot\Sigma\bigl(c_W(\alpha,t);K_N,\one_{K_Nh^{-1}K_N}\bigr).
\]
\formallabel{UpperTheory.Gal.W2.FixZ_holds}
\end{corollary}

\begin{proof}
Since $N\mid N'$ we have $K_{N'}\subseteq K_N$, and admissibility
of $h$ gives $hK_{N'}h^{-1}\subseteq K_N$, so $K_{N'}\subseteq\Gamma_h$; $[K_N:K_{N'}]$ is finite. By
\cref{pf:ks2:hfp}(4) no $y\ne1$ sees any point through $K_N$, and by \cref{pf:ks2:hfp}(3) the finiteness
hypothesis of \cref{pf:ks2:levelchange} holds. Apply \cref{pf:ks2:levelchange} with $K=K_N$, $K'=K_{N'}$
and $c=c_W(\alpha,t)$ to the right-hand side of \cref{pf:ks2:hfp}(5).
\end{proof}

\begin{proof}[Proof of \cref{pf:ks2:fixk}]
The proof uses, besides \cref{pf:ks2:hfp} and
\cref{pf:ks2:fixz}, the following facts about a Kuga--Sato tower $T$ over $\C$:
\begin{enumerate}[label=(B\arabic*)]
\item\label{pf:ks2:B:tr} every transition map $\mathrm{tr}(N',N,h)$ is finite and \'etale
  (\lref{UpperTheory.Gal.FixZ.translation_finite_etale_of_isogeny_data}; it is deduced from the
  description of $\mathrm{tr}(N',N,h)$ as the fibre power of the total map of a Hecke witness, the
  uniqueness of witnesses, the finite \'etaleness of isogenies of abelian schemes whose degree is
  invertible on the base, and \cref{pf:ks2:gensm});
\item\label{pf:ks2:B:al} $[\alpha]$ is finite and \'etale if all $\alpha_s\ne0$
  (\lref{UpperTheory.Gal.FixZ.alFEt_of}: multiplication by a nonzero element of $\kstwoOF$ is an
  isogeny of degree invertible in $\C$);
\item\label{pf:ks2:B:leg} for $h$ admissible from $N'$ to $N$ there is a leg $c$ for $h$ from
  $M_{N'}$ to $M_N$ with $\mathrm{tr}(N',N,h)=c^{(k)}$ and $[\alpha]\circ\mathrm{tr}(N',N,h)=c^{(k)}_\alpha$
  under $Z\cong\kstwoKS^k$ (\lref{UpperTheory.Gal.FixZ.legOf_of}, from the definition of a
  Kuga--Sato tower);
\item\label{pf:ks2:B:tan} isomorphisms of schemes over $\C$ compatible with two pairs of maps transport
  fixed points and tangent matrices (\lref{UpperTheory.Gal.FixZ.tanIso_of}); a nonempty scheme smooth
  of relative dimensions $d$ and $d'$ over $\C$ has $d=d'$ (\lref{UpperTheory.Gal.FixZ.dimUniq_of});
\item\label{pf:ks2:B:cx} if $K'\subseteq K$, $hK'h^{-1}\subseteq K$, the finiteness of \cref{pf:ks2:hfp}(3)
  holds at level $K'$ and stabilizers at level $K$ are trivial, then $\mathrm{Fin}(K,\one_{Kh^{-1}K})$
  (\lref{UpperTheory.Gal.FixZ.cxFinOf_of}: if $y$ is counted for the point of $a=((g,g_p),z)$ at
  level $K$, write $g^{-1}y_fg=k_1h^{-1}k_2$ with $k_i\in K$; then $y$ sees $((gk_1,g_p),z)$ through
  $h^{-1}K$, and the resulting map from the counted triples to the points of level $K'$ seen through
  $h^{-1}K$ is injective, by the triviality of stabilizers and the argument of the proof of
  \cref{pf:ks2:levelchange}).
\end{enumerate}
Let the data of $\sfname{FixK}$ be given. Since $a=1$, $\tau_i=t_0$ for the unique $i$. By \ref{pf:ks2:B:leg}
choose legs $c_1$ for $1$ and $c_2$ for $h$ from $M_{N'}$ to $M_N$. If $hg\in K_N$ then
$g\in h^{-1}K_N\subseteq K_Nh^{-1}K_N$, so $g$ is regular. The map $c_1^{(k)}$ corresponds to
$\mathrm{tr}(N',N,1)$, which is \'etale. Hence \cref{pf:ks2:hfp} and \cref{pf:ks2:fixz} apply (with the
divisibility of the $\alpha_s$ by $2$).

(1) follows from \ref{pf:ks2:B:cx}, using \cref{pf:ks2:hfp}(3),(4). (2): $c_2=[\alpha]\circ
\mathrm{tr}(N',N,h)$ is a composite of finite \'etale maps by \ref{pf:ks2:B:tr},\ref{pf:ks2:B:al}. (3): If
$Z_N$ is empty, then $Z_{N'}$ has no complex point (its image under $c_1$ would be a point of $Z_N$), the
set $\mathcal F$ of \cref{pf:ks2:fixz} is empty, so the right-hand side of \cref{pf:ks2:fixz} vanishes,
and $\mathcal F=\emptyset$ works. Otherwise, $Z_N\cong\kstwoKS_N^k$ is smooth of relative dimension
$n-1+kg$ by \cref{pf:ks2:hfp}(1), so $\delta=n-1+kg$ by \ref{pf:ks2:B:tan}. Transport the set $\mathcal F$
and the matrices $A_x$ of \cref{pf:ks2:fixz} along the isomorphisms $Z_{N'}\cong\kstwoKS_{N'}^k$ and
$Z_N\cong\kstwoKS_N^k$ by \ref{pf:ks2:B:tan}; they have the required properties. Finally the right-hand side
of \cref{pf:ks2:fixz} is that of \eqref{pf:ks2:eq:fixk}: $[\Gamma_h:K_{N'}]=[\Gamma_h:\Gamma_h\cap
K_{N'}]$, the weight $c_W$ of \eqref{pf:ks2:eq:cw} equals the weight of the formal statement
(\lref{UpperTheory.Gal.FixZ.wtEq_of}), and the indicator of the double coset is the same function in
both formulations.
\end{proof}

\subsubsection{Proof of the fixed-point formula on moduli schemes}\label{pf:ks2:sec:fixproof}

The complex uniformization is formulated with the following notions.

\begin{definition}[Marked period lattices]\label{pf:ks2:def:periods}
Let $N\ge1$ and let $\rho\colon O\to M_g(\C)$ be a ring map.
\begin{enumerate}
\item A \emph{polarized $O$-lattice} is a free $\Z$-module $D'$ of finite rank with an action of $O$
  and a nondegenerate alternating form $E_{D'}$ with $E_{D'}(cx,y)=E_{D'}(x,c^*y)$; a \emph{marking}
  is a map $t\colon\Lambda\to D'$. A \emph{marked isometry} $w\colon(D',t)\to(D'',t'')$ is an
  $O$-linear isomorphism $D'\to D''$ preserving the forms with $w(t(\lambda))\equiv t''(\lambda)\bmod
  ND''$ for $\lambda\in\Lambda$.
\item $(D',t)$ lies in the class $\mathrm{Cls}_N$ if (i) for every $k\ge1$ prime to $p$ there are an
  additive map $f_k\colon\Lambda\to D'$ and an integer $\nu_k$ prime to $k$, with $f_N=t$ and
  $f_{kl}\equiv f_k$, $\nu_{kl}\equiv\nu_k$ modulo $k$, such that $f_k$ induces an $O$-linear
  isomorphism $\Lambda/k\Lambda\to D'/kD'$ and $E_{D'}(f_kx,f_ky)\equiv\nu_kE_\Lambda(x,y)\bmod k$; and
  (ii) $E_{D'}=2E_1$ with $E_1$ perfect up to an integer prime to $p$.
\item A \emph{period} of $D'$ is an injective additive map $P\colon D'\to\C^g$ with
  $P(cx)=\rho(c)P(x)$ for $c\in O$, whose image is discrete and spans $\C^g$ over $\R$, such that
  $\Im H(Px,Py)=E_{D'}(x,y)$ for some positive definite hermitian form $H$ on $\C^g$.
\item Let $\mathcal A$ be an abelian scheme over $S$ with an action of $O$, a polarization and a level
  structure $\eta\colon\Lambda\to\mathcal A(S)$ (for instance an object of level $N$), $y\in
  S(\C)$, and $e\colon\C^g\to\mathcal A_y(\C)$ an \emph{exponential}: a surjective holomorphic
  homomorphism whose derivative at $0$ identifies $\C^g$ with the tangent space at the origin. A
  marked lattice $(D',t)$ with a period $P$ \emph{uniformizes $\mathcal A$ at $y$ via $e$} if $\ker
  e=P(D')$, $e(\rho(c)z)=c\cdot e(z)$ for $c\in O$, the Weil pairing of $e(P(x)/k)$ and $e(P(x')/k)$
  is $\exp(-2\pi iE_{D'}(x,x')/k)$ for all $k\ge1$, and $\eta(\lambda)=e(P(t(\lambda))/N)$ for
  $\lambda\in\Lambda$.
\end{enumerate}
\formal{GlobalTheory.PELU.PLat}\formal{GlobalTheory.PELU.IsMarkedIsom}\formal{GlobalTheory.PELU.Genus}\formal{GlobalTheory.UPel.HT.HalfPerf}\formal{GlobalTheory.UPel.HT.Datum.Cls}\formal{GlobalTheory.PELU.IsPeriod}\formal{GlobalTheory.PELU.IsExp}\formal{GlobalTheory.PELU.UnifAt}
\end{definition}

We use the following statements, all proved in the formal development. Here $X$ is a
generic model over $\C$ of a level $N$, and the data are as in \cref{pf:ks2:hfp}.
\begin{enumerate}[label=(U\arabic*)]
\item\label{pf:ks2:U:rho} (\lref{UpperTheory.Gal.W2.RhoSig_holds}) There is a ring map
  $\rho\colon O\to M_g(\C)$ \emph{of signature $(t_0,\varepsilon_D)$}: $\kstwoOF\to O$ is onto (as $m=1$), $\rho$ is compatible with
  the involutions, every embedding of $F$ is non-real, and
  $\tr\rho(c)=\sum_\tau\mu(\tau)\tau(c)$ for $c\in\kstwoOF$, with $\mu(t_0)=1$, $\mu(\bar t_0)=n-1$ and
  $\mu(\tau)=n\varepsilon_D(\tau)$ otherwise; also
  $\varepsilon_D(\tau)+\varepsilon_D(\bar\tau)=1$ for $\tau\ne t_0,\bar t_0$.
\item\label{pf:ks2:U:unif} (\lref{UpperTheory.Gal.W2.GUnif_holds}, \emph{complex uniformization})
  $X$ is smooth of relative dimension $n-1$, quasi-compact and separated, the universal abelian scheme
  is smooth of relative dimension $g$, and its complex points are uniformized by marked period
  lattices (\cref{pf:ks2:def:periods}): there are finitely many marked polarized $O$-lattices
  $(\Lambda_i,t_i)$ in $\mathrm{Cls}_N$, pairwise not marked-isometric, such that every marked lattice
  in $\mathrm{Cls}_N$ which has a period is marked-isometric to one of them, and for every period $P$
  of $\Lambda_i$ a complex point $u_i(P)$ of $X$ and an exponential
  $\exp_{i,P}\colon\C^g\to\mathcal A_{u_i(P)}$ via which $(\Lambda_i,t_i,P)$ uniformizes the universal
  object at $u_i(P)$, such that every complex point is some $u_i(P)$, $u_i(P)=u_j(P')$ forces $i=j$,
  and $u_i(P)=u_i(P')$ if and only if $P'\circ w=h\circ P$ for a marked self-isometry $w$ of
  $(\Lambda_i,t_i)$ and a $\C$-linear automorphism $h$ of $\C^g$; in particular the fibre of
  $\mathcal A$ at $u_i(P)$ is $\C^g/P(\Lambda_i)$ via $\exp_{i,P}$.
\item\label{pf:ks2:U:dic} (\cref{pf:ks2:pdic,pf:ks2:rdic} below) For every form $\beta$ there are a
  $p$-adic dictionary and a real dictionary (\cref{pf:ks2:def:coord}): a $\Gamma_\beta$-equivariant
  bijection between $\kstwoGp/\kstwoKp$ and the lattices in $\Lambda\otimes\Q$ that are self-dual at
  $p$ for a positive multiple of the form attached to $\beta$, and a $\Gamma_\beta$-equivariant
  correspondence between periods of $\Lambda\otimes\Q$ for $\beta$ and lines $\C z$,
  $z\in\mathcal{P}_\beta$, under which an element $y$ acting on the periods by a matrix $H$ and on $z$
  by the eigenvalue $\lambda$ has $\kappa(\lambda)$ as an eigenvalue of $H$ on the $t_0$-isotypic part
  of $\C^g$.
\item\label{pf:ks2:U:unifg} (\cref{pf:ks2:unifg} below) There is a bijection
  $\Phi\colon X(\C)\to\mathcal{X}(K_N,\kstwoKp)$ such that for every complex point $x$ there are
  $([\beta],a)$ with $\Phi(x)=([\beta],[a])$, an index $i$, a period $P$ of $\Lambda_i$ with $x=u_i(P)$,
  and an embedding $j\colon\Lambda_i\to\Lambda\otimes\Q$ such that $(P,j,a)$ are coordinates
  (\cref{pf:ks2:def:coord}: compatible with the two dictionaries and the level structure).
\item\label{pf:ks2:U:stab} (\lref{UpperTheory.Gal.W2.StabTriv_holds}) If $N\ge3$ is prime to $p$,
  every $y$ seeing some $a$ through $K_N$ equals $1$.
\item\label{pf:ks2:U:fwd} (\lref{UpperTheory.Gal.W2.FixFwd_holds}) If $x=u_i(P)\in X'(\C)$ with
  coordinates $(P,j,a)$ satisfies $c_{1,q}(x)=c_{2,q}(x)$, then some $y\in\Gamma_\beta$ sees $a$ through
  $h^{-1}K_N$.
\item\label{pf:ks2:U:loc} (\cref{pf:ks2:fixloc} below) The local contribution over a fixed point of the
  base.
\item\label{pf:ks2:U:tan} (\lref{UpperTheory.Gal.W2.TanExists_holds}) If $Z\to\Spec\C$ is smooth of
  relative dimension $d$ and $c_1$ is \'etale, every complex fixed point of $(c_1,c_2)$ has a tangent
  matrix in $M_d(\C)$ (the pullbacks of local coordinates along the \'etale map $c_1$ are local
  coordinates).
\item\label{pf:ks2:U:fin} (\lref{UpperTheory.Gal.W2.FixFinite_holds}) If $Z$ is separated, $Y$ is
  quasi-compact and locally of finite type over $\C$, and every complex fixed point of $(c_1,c_2)$ has a
  tangent matrix $A$ with $\det(1-A)\ne0$, then the set of complex fixed points is finite (transversal
  fixed points are isolated, and the fixed locus is closed in the quasi-compact $Y$).
\item\label{pf:ks2:U:pow} (\lref{UpperTheory.Gal.W2.PowSm_holds}, \lref{UpperTheory.Gal.W2.PwOver_holds})
  The $k$-th fibre power of an abelian scheme of relative dimension $g$ over a smooth, quasi-compact,
  separated $\C$-scheme of relative dimension $d$ is smooth of relative dimension $d+kg$, quasi-compact,
  separated and locally of finite type; the powers $c^{(k)}$, $c^{(k)}_\alpha$ of a leg lie over $c_q$ and
  act on the factors through $c_{\mathrm{tot}}$, followed by $\alpha_s$ for $c^{(k)}_\alpha$.
\item\label{pf:ks2:U:eig} (\lref{UpperTheory.Gal.W2.EigSame_holds}) If $a,a'$ represent the same
  point at level $K$ and $S$ is invariant under conjugation by $K$, there is a bijection $E$ of
  $\Gamma_\beta$ with $\epsilon(1)=1$ such that $y$ sees $a$ through $S$ iff $\epsilon(y)$ sees $a'$ through $S$, and then
  $\mathrm{minpoly}_F(\epsilon(y))=\mathrm{minpoly}_F(y)$ and $\mathrm{eig}(\epsilon(y),a')=\mathrm{eig}(y,a)$.
\end{enumerate}

\begin{lemma}[Summation over points]\label{pf:ks2:sumpts}
Let $K\subseteq\kstwoG$, let $S\subseteq\kstwoG$ be invariant under conjugation by $K$, and assume that
no $y\ne1$ sees any $a$ through $K$ and that at most one $y$ sees a given $a$ through $S$. Let
$\Phi\colon\mathcal T\to\mathcal{X}(K,\kstwoKp)$ be a bijection, $\mathcal B\subseteq\mathcal T$ finite
and $w\colon\mathcal T\to\C$, and assume that for every $x\in\mathcal T$ there is $a$ with
$\Phi(x)=([\beta],[a])$ such that $x\in\mathcal B$ iff some $y$ sees $a$ through $S$, and
$w(x)=c(\mathrm{minpoly}_F(y),\mathrm{eig}(y,a))$ for such $y$. Then
$\sum_{x\in\mathcal B}w(x)=\Sigma(c;K,\one_S)$.
\formallabel{UpperTheory.Gal.W2.SumPts_holds}
\end{lemma}

\begin{proof}
For a point $s=([\beta],q)$ let $\sigma(s)$ be its term in
\eqref{pf:ks2:eq:cxtotal} with $f=\one_S$; since stabilizers are trivial, $\sigma(s)$ is the sum of
$c(\mathrm{minpoly}_F(y),\mathrm{eig}(y,a_q))$ over the $y$ seeing the chosen representative $a_q$
through $S$. For $x\in\mathcal T$ let $a$ be as in the hypothesis; $a$ and $a_q$ represent the same
point, so by \ref{pf:ks2:U:eig} the elements seeing $a_q$ through $S$ correspond bijectively to those
seeing $a$, with the same values of $c$. Hence $\sigma(\Phi(x))=w(x)$ if $x\in\mathcal B$ (there is
exactly one such $y$) and $\sigma(\Phi(x))=0$ otherwise. So $\sigma$ has finite support contained in
$\Phi(\mathcal B)$, and $\Sigma(c;K,\one_S)=\sum_s\sigma(s)=\sum_{x\in\mathcal B}w(x)$.
\end{proof}

\begin{proof}[Proof of \cref{pf:ks2:hfp}]
Choose $\rho$ by
\ref{pf:ks2:U:rho} and apply \ref{pf:ks2:U:unif} to $X$ and $X'$; fix a uniformization of $X'$. By
\ref{pf:ks2:U:pow}, $\mathcal A_X^k\to\Spec\C$ is smooth of relative dimension $\delta=n-1+kg$, which is
(1), and $\mathcal A_{X'}^k$ is quasi-compact and locally of finite type, $\mathcal A_X^k$ separated.
Since $c_1^{(k)}$ is \'etale, \ref{pf:ks2:U:tan} gives a tangent matrix $A_x$ at every fixed point $x$; we
fix such a choice. Choose dictionaries for every form (\ref{pf:ks2:U:dic}) and let $\Phi$ be the bijection
of \ref{pf:ks2:U:unifg} for $X'$. For a complex point $x$ of $\mathcal A_{X'}^k$ let $\pi(x)\in X'(\C)$ be
its image. By \ref{pf:ks2:U:pow}, if $x$ is a fixed point of $(c_1^{(k)},c_{2,\alpha}^{(k)})$, then
$\pi(x)$ is a fixed point of $(c_{1,q},c_{2,q})$.

\emph{Step 1: local data.} Let $x'=u_i(P)\in X'(\C)$ have coordinates $(P,j,a)$, $a\in\mathcal{D}_\beta$. If
$x'$ is a fixed point of $(c_{1,q},c_{2,q})$, \ref{pf:ks2:U:fwd} gives $y$ seeing $a$ through $h^{-1}K_N$.
Conversely, if $y$ sees $a$ through $h^{-1}K_N$, \cref{pf:ks2:fixloc} applies: the set $\mathcal F_{x'}$ of
fixed points of $(c_1^{(k)},c_{2,\alpha}^{(k)})$ over $x'$ is finite and nonempty, every fixed point $x$ over
$x'$ has $\det(1-A_x)\ne0$, and
$\sum_{x\in\mathcal F_{x'}}\mathrm{lef}(A_x)(t)=c_W(\alpha,t)(\mathrm{minpoly}_F(y),\mathrm{eig}(y,a))$.

\emph{Step 2: finiteness, (2).} Every fixed point $x$ lies over a fixed point $\pi(x)=u_i(P)$ of the base,
so Step 1 gives $\det(1-A_x)\ne0$. By \ref{pf:ks2:U:fin} the set $\mathcal F$ of fixed points is finite.
This proves (2).

\emph{Step 3: stabilizers, (4).} (4) is \ref{pf:ks2:U:stab} for the level $N$; the same statement for
$N'$ shows that stabilizers at level $K_{N'}$ are trivial. If $y,y'$ see $a$ through $h^{-1}K_N$, then
$y^{-1}y'$ fixes the line of $a$, satisfies the condition at $p$, and $g^{-1}(y^{-1}y')_fg=(g^{-1}y_fg)^{-1}
(g^{-1}y'_fg)\in K_N^{-1}hh^{-1}K_N=K_N$; so $y=y'$ by (4).

\emph{Step 4: (3).} Let $\mathcal B=\pi(\mathcal F)\subseteq X'(\C)$, a finite set. By Step 1, a point
$x'\in X'(\C)$ with coordinates $(P,j,a)$ lies in $\mathcal B$ iff some $y$ sees $a$ through $h^{-1}K_N$
($\Rightarrow$ by \ref{pf:ks2:U:fwd}, $\Leftarrow$ because $\mathcal F_{x'}\ne\emptyset$). If
$([\beta],q)\in\mathcal{X}(K_{N'},\kstwoKp)$ and some $y$ sees $a_q$ through $h^{-1}K_N$, let
$x'=\Phi^{-1}([\beta],q)$ with coordinates $(P,j,a)$; $a$ and $a_q$ represent the same point, and
$h^{-1}K_N$ is invariant under conjugation by $K_{N'}$ (proof of \cref{pf:ks2:levelchange}), so by
\ref{pf:ks2:U:eig} some element sees $a$, and $x'\in\mathcal B$. Hence the set in (3) is contained in
$\Phi(\mathcal B)$, which is finite.

\emph{Step 5: (5).} Group the fixed points according to their image: $\sum_{x\in\mathcal
F}\mathrm{lef}(A_x)(t)=\sum_{x'\in\mathcal B}w(x')$ with $w(x')=\sum_{x\in\mathcal F,\ \pi(x)=x'}\mathrm{lef}
(A_x)(t)$. For $x'\in\mathcal B$ with coordinates $(P,j,a)$ and $y$ seeing $a$ through $h^{-1}K_N$, the
inner set is $\mathcal F_{x'}$, so $w(x')=c_W(\alpha,t)(\mathrm{minpoly}_F(y),\mathrm{eig}(y,a))$ by Step 1.
Now \cref{pf:ks2:sumpts}, applied with $K=K_{N'}$, $S=h^{-1}K_N$, $\mathcal T=X'(\C)$, $\Phi$, $\mathcal B$,
$w$ and $c=c_W(\alpha,t)$ (its hypotheses hold by Steps 3 and 4), gives
$\sum_{x'\in\mathcal B}w(x')=\Sigma(c_W(\alpha,t);K_{N'},\one_{h^{-1}K_N})$.
\end{proof}

\subsubsection{The local contribution}

We keep the notation of \ref{pf:ks2:U:unif}: $\Lambda_i$ is a marked polarized lattice with alternating form $E_i$, $P$ a period
with positive definite hermitian matrix $M_0$, and for $z=(z_1,\dots,z_k)\in(\C^g)^k$ we write
$\mathrm{pt}_P(z)$ for the complex point of $\mathcal A_{X'}^k$ over $u_i(P)$ with coordinates
$\exp_{i,P}(z_s)$. Let $\mathfrak{t}_{M_0}$ be the space of matrices $C\in M_g(\C)$ such that $C^{\mathsf
T}M_0$ is symmetric and satisfies $C^{\mathsf T}M_0\,\overline{\rho(\bar b)}=\rho(b)^{*}C^{\mathsf T}M_0$
for $b\in O$; it is the space of first
order deformations $P\mapsto(1+C)P$ of the period, i.e.\ the tangent space of $X'$ at $u_i(P)$ in period
coordinates.

\begin{proposition}[Local contribution]\label{pf:ks2:fixloc}
In the situation of \cref{pf:ks2:hfp}, let $x'=u_i(P)\in X'(\C)$ have coordinates $(P,j,a)$ with
$a\in\mathcal{D}_\beta$, and let $y\in\Gamma_\beta$ see $a$ through $h^{-1}K_N$. Then the set $\mathcal F_{x'}$
of complex points $x$ of $\mathcal A_{X'}^k$ over $x'$ with $c_1^{(k)}(x)=c_{2,\alpha}^{(k)}(x)$ is finite and
nonempty, and for every choice of matrices $A_x$ with $\mathrm{Tan}(n-1+kg,c_1^{(k)},c_{2,\alpha}^{(k)},
x,A_x)$ ($x\in\mathcal F_{x'}$) one has $\det(1-A_x)\ne0$ and
\[
\sum_{x\in\mathcal F_{x'}}\frac{\det(1-tA_x)}{\det(1-A_x)}=c_W(\alpha,t)\bigl(\mathrm{minpoly}_F(y),\mathrm{eig}(y,a)\bigr)
\qquad(t\in\C).
\]
\formallabel{UpperTheory.Gal.W2.FixLoc_holds}
\end{proposition}

The proof rests on three lemmas: a count of fixed points on complex tori, the computation of the
local terms in period coordinates (\cref{pf:ks2:locsum}), and the computation of the eigenvalues
(\cref{pf:ks2:eigtop}).

\begin{lemma}[Fixed points on a complex torus]\label{pf:ks2:torus}
Let $W$ be a finite-dimensional complex vector space, $\Lambda\subseteq W$ a lattice (discrete, spanning
$W$ over $\R$) and $\Theta\colon W\to W$ a $\C$-linear map with $\Theta\Lambda\subseteq\Lambda$ and
$\det(1-\Theta)\ne0$. Then $\{z\in W:\Theta z-z\in\Lambda\}/\Lambda$ is finite of cardinality
$\det(1-\Theta)\,\overline{\det(1-\Theta)}=|\det(1-\Theta)|^2$.
\formallabel{UpperTheory.Gal.W2.TorusCount_holds}
\end{lemma}

\begin{proof}
$\Phi=1-\Theta$ is a $\C$-linear automorphism of $W$ with $\Phi\Lambda\subseteq\Lambda$. The set in
question is $\Phi^{-1}(\Lambda)/\Lambda$, and $\Phi$ induces an isomorphism $\Phi^{-1}(\Lambda)/\Lambda\cong
\Lambda/\Phi\Lambda$. The index of $\Phi\Lambda$ in $\Lambda$ is $|\det_\Z(\Phi|_\Lambda)|=|\det_\R\Phi|$, and
for a $\C$-linear map $\det_\R\Phi=\det_\C\Phi\cdot\overline{\det_\C\Phi}$ (in a complex basis, the real
matrix is conjugate over $\C$ to $\mathrm{diag}(\Phi,\bar\Phi)$).
\end{proof}

\begin{lemma}[Determinants from characteristic polynomials]\label{pf:ks2:detchar}
If $A\in M_d(\C)$, $T$ and $\Theta$ are endomorphisms of finite-dimensional complex vector spaces and
$\mathrm{charpoly}(A)=\mathrm{charpoly}(T)\,\mathrm{charpoly}(\Theta)$, then
$\det(1-tA)=\det(1-tT)\det(1-t\Theta)$ for all $t\in\C$.
\formallabel{UpperTheory.Gal.W2.DetOfCharpoly_holds}
\end{lemma}

\begin{proof}
$\det(1-tA)=\prod_\lambda(1-t\lambda)$ over the roots $\lambda$ of $\mathrm{charpoly}(A)$ with
multiplicity, and the multiset of roots of a product is the sum of the multisets.
\end{proof}

\begin{lemma}[Local sum]\label{pf:ks2:locsum}
In the situation of \cref{pf:ks2:fixloc}, let $\gamma$ be an $O$-linear endomorphism of $\Lambda_i$ and
$H\in M_g(\C)$ with $P\circ\gamma=H\circ P$ and $E_i(\gamma\cdot,\gamma\cdot)=mE_i$ for some integer $m>0$,
let $\mathrm{co}\colon\C^{n-1}\to\mathfrak{t}_{M_0}$ be a linear isomorphism, $T$ the endomorphism of
$\C^{n-1}$ with $\mathrm{co}(Tb)=H\,\mathrm{co}(b)\,(H^{*})^{-1}$, and $\Theta$ the endomorphism
$z\mapsto(H\rho(\alpha_s)z_s)_s$ of $(\C^g)^k$. Assume:
\begin{enumerate}[label=(\alph*)]
\item for all periods $P_1,P_2$ of $\Lambda_i$ with $P_2\circ\gamma=P_1$ and all $z\in(\C^g)^k$,
  $c_{2,\alpha}^{(k)}(\mathrm{pt}_{P_1}(z))=c_1^{(k)}(\mathrm{pt}_{P_2}(\rho(\alpha)z))$, where
  $\rho(\alpha)z=(\rho(\alpha_s)z_s)_s$;
\item $c_1^{(k)}(\mathrm{pt}_P(z))=c_1^{(k)}(\mathrm{pt}_P(z'))$ implies $\mathrm{pt}_P(z)=\mathrm{pt}_P(z')$;
\item $\det(1-T)\ne0$ and $\det(1-\Theta)\ne0$.
\end{enumerate}
Then $\mathcal F_{x'}$ is finite and nonempty, and for all tangent matrices $A_x$ as in
\cref{pf:ks2:fixloc}: $\det(1-A_x)=\det(1-T)\det(1-\Theta)\ne0$ and
\[
\sum_{x\in\mathcal F_{x'}}\frac{\det(1-tA_x)}{\det(1-A_x)}=\frac{\det(1-tT)}{\det(1-T)}\,
\det(1-t\Theta)\,\overline{\det(1-\Theta)}.
\]
\formallabel{UpperTheory.Gal.W2.LocSum_holds}
\end{lemma}

\begin{proof}
Let $\Lambda^k=P(\Lambda_i)^k\subseteq(\C^g)^k$, a lattice
(\lref{UpperTheory.Gal.W2.PiLattice_holds}). By the description of the fibre
(\lref{UpperTheory.Gal.W2.PtFibre_holds}), the complex points of $\mathcal A_{X'}^k$ over $x'$ are
exactly the $\mathrm{pt}_P(z)$, and $\mathrm{pt}_P(z)=\mathrm{pt}_P(z')$ iff $z-z'\in\Lambda^k$. $\Theta$
preserves $\Lambda^k$, because $H\rho(\alpha_s)P(u)=P(\gamma(\alpha_su))$. Applying (a) along the
deformation $P\mapsto(1+C)P$ of the period and evaluating at $C=0$
(\lref{UpperTheory.Gal.W2.CorrChart_holds}) gives
$c_{2,\alpha}^{(k)}(\mathrm{pt}_P(z))=c_1^{(k)}(\mathrm{pt}_P(\Theta z))$ for all $z$. Together with (b),
$\mathrm{pt}_P(z)$ is a fixed point iff $\mathrm{pt}_P(\Theta z)=\mathrm{pt}_P(z)$ iff $\Theta z-z\in\Lambda^k$.
By \cref{pf:ks2:torus} the fixed points over $x'$ are therefore in bijection with a set of cardinality
$|\det(1-\Theta)|^2$; in particular $\mathcal F_{x'}$ is finite, and nonempty since $\mathrm{pt}_P(0)$ is a
fixed point. At every fixed point $x=\mathrm{pt}_P(z_0)$ any tangent matrix satisfies
$\mathrm{charpoly}(A_x)=\mathrm{charpoly}(T)\,\mathrm{charpoly}(\Theta)$
(\lref{UpperTheory.Gal.W2.LocTan_holds}, which computes the tangent matrix at a fixed point
over $x'$ in period coordinates on the base and linear coordinates on the fibres, using (a) for
periods near $P$). By \cref{pf:ks2:detchar},
$\det(1-tA_x)=\det(1-tT)\det(1-t\Theta)$ for all $t$, in particular $\det(1-A_x)\ne0$ by (c). All local
terms are equal, so the sum is
\[
|\det(1-\Theta)|^2\cdot\frac{\det(1-tT)\det(1-t\Theta)}{\det(1-T)\det(1-\Theta)}
=\frac{\det(1-tT)}{\det(1-T)}\det(1-t\Theta)\overline{\det(1-\Theta)}.\qedhere
\]
\end{proof}

\begin{lemma}[Eigenvalues]\label{pf:ks2:eigtop}
Let $\rho$ be of signature $(t_0,\varepsilon)$, $P$ a period of a polarized lattice $\Lambda'$ with hermitian
matrix $M_0$, $\gamma$ and $H$ as in \cref{pf:ks2:locsum}, and $\chi\in F[X]$ monic of degree $n$ and
separable such that a polynomial $q\in\kstwoOF[X]$ satisfies $q(\gamma)=0$ iff $\chi$ divides $q$ in $F[X]$. Let $\mathrm{co}$, $T$ be as in \cref{pf:ks2:locsum} (they
exist by \lref{UpperTheory.Gal.W2.CspDim_holds} and \lref{UpperTheory.Gal.W2.AdExists_holds}, which
show that $\det H\ne0$, $H$ commutes with $\rho(O)$ and $C\mapsto HC(H^*)^{-1}$ preserves $\mathfrak{t}_{M_0}$),
and $\alpha\in(2\kstwoOF)^k$. Call $\mu\in\C$ a \emph{mark} of $H$ if $H$ has
an eigenvector with eigenvalue $\mu$ on which $\rho(c)$ acts by $t_0(c)$ for all $c\in\kstwoOF$. Then $H$ has a
mark, and for every mark $\mu$: $\mu\in R_{t_0}(\chi)$, $\det(1-T)\ne0$, $\det(1-\Theta)\ne0$, and for all
$t\in\C$
\begin{equation}\label{pf:ks2:eq:eigtop}
\frac{\det(1-tT)}{\det(1-T)}\,\det(1-t\Theta)\,\overline{\det(1-\Theta)}=\kstwoloc(\alpha;t,\chi,\mu).
\end{equation}
\formallabel{UpperTheory.Gal.W2.EigTop_holds}
\end{lemma}

\begin{proof}
Since $\chi$ is separable, $H$ and $\rho(\kstwoOF)$ are
simultaneously diagonalizable (\lref{UpperTheory.Gal.W2.DiagExists_holds}): there are a basis
$e_1,\dots,e_g$ of $\C^g$, embeddings $\sigma_i$ and complex numbers $\lambda_i$ with $\rho(c)e_i=\sigma_i(c)e_i$ and
$He_i=\lambda_ie_i$. For an embedding $\tau$ let $\mathrm{ev}(\tau)$ be the multiset of $\lambda_i$ with
$\sigma_i=\tau$. Then:
\begin{itemize}
\item $\#\mathrm{ev}(\tau)=\mu(\tau)$ (the multiplicity of \ref{pf:ks2:U:rho}), by comparing traces
  (\lref{UpperTheory.Gal.W2.TypeCount_holds});
\item $R_\tau(\chi)=\mathrm{ev}(\tau)+\overline{\mathrm{ev}(\bar\tau)}$ for every $\tau$
  (\lref{UpperTheory.Gal.W2.RootsFull_holds});
\item $\mu$ is a mark iff $\mathrm{ev}(t_0)=\{\mu\}$ (\lref{UpperTheory.Gal.W2.MarkDiag_holds});
\item $\dim\mathfrak{t}_{M_0}=n-1$ and $\det(1-tT)=\prod_{r\in\mathrm{ev}(\bar t_0)}(1-t\mu/\bar r)$
  (\lref{UpperTheory.Gal.W2.BallChar_holds});
\item $\det(1-t\Theta)=\prod_{s=1}^k\prod_{i=1}^g(1-\sigma_i(\alpha_s)\,t\,\lambda_i)$
  (\lref{UpperTheory.Gal.W2.FibDet_holds});
\item $1-\sigma_i(2x)\lambda_i\ne0$ for $x\in\kstwoOF$ (\lref{UpperTheory.Gal.W2.EvIntegral_holds}; indeed
  $H$ represents the lattice endomorphism $\gamma$, so $\lambda_i$ is a root of the characteristic polynomial
  of $\gamma$ on $\Lambda'$, a monic polynomial with integer coefficients, hence an algebraic integer, and
  $\sigma_i(2x)\lambda_i$ lies in $2\overline{\Z}$, which does not contain $1$);
\item $\prod_i(1-\sigma_i(x)t\lambda_i)\cdot\overline{\prod_i(1-\sigma_i(x)\lambda_i)}=\kstwoFib(x;t,\chi,\mu)$ when
  $\mathrm{ev}(t_0)=\{\mu\}$ (\lref{UpperTheory.Gal.W2.FibW_holds}).
\end{itemize}
As $\mu(t_0)=1$, $\mathrm{ev}(t_0)$ is a single element, which is a mark; and a mark $\mu$ satisfies
$\mathrm{ev}(t_0)=\{\mu\}$, so $\mu\in R_{t_0}(\chi)=\{\mu\}+\overline{\mathrm{ev}(\bar t_0)}$. Since $\chi$ is
separable, $\mu\ne\bar r$ for $r\in\mathrm{ev}(\bar t_0)$, so $\det(1-T)\ne0$, and
$\det(1-tT)/\det(1-T)=\prod_{r}(\bar r-t\mu)/(\bar r-\mu)=\kstwoBall(t,\chi,\mu)$ because
$R_{t_0}(\chi)\smallsetminus\{\mu\}=\overline{\mathrm{ev}(\bar t_0)}$. Since the $\alpha_s$ are even,
$\det(1-\Theta)\ne0$. Finally $\det(1-t\Theta)\overline{\det(1-\Theta)}=\prod_s\kstwoFib(\alpha_s;t,\chi,\mu)$
by the last two items, and \eqref{pf:ks2:eq:eigtop} follows. (The last item is the bookkeeping identity
behind the definition of $\kstwoFib$: on the $t_0$-line $H\rho(x)$ acts by $t_0(x)\mu$; on the
$\bar t_0$-part by $\overline{t_0(x)}\lambda$ with $\bar\lambda$ running over $R_{t_0}(\chi)\smallsetminus\{\mu\}$;
on the $\tau$-part, $\tau\ne t_0,\bar t_0$, by $\tau(x)r$ with $r\in R_\tau(\chi)$ if $\varepsilon(\tau)=1$ and
not at all if $\varepsilon(\tau)=0$; conjugating the second factor exchanges $\tau$ and $\bar\tau$.)
\end{proof}

\begin{proof}[Proof of \cref{pf:ks2:fixloc}]
Let $\lambda=\mathrm{eig}(y,a)$; $\lambda\ne0$ since $y_v$ is
invertible and $z\ne0$, and $y^{-1}$ fixes the line of $a$ with eigenvalue $\lambda^{-1}$.
\begin{enumerate}
\item By \lref{UpperTheory.Gal.W2.SheetOf_holds} the uniformization gives period coordinates around
  $u_i(P)$ (a \emph{sheet}), in which $\mathrm{pt}_P$ is defined.
\item By \lref{UpperTheory.Gal.W2.FixGa_holds}, since $y$ sees $a$, there are an $O$-linear
  endomorphism $\gamma$ of $\Lambda_i$ with $j(\gamma u)=y^{-1}\cdot j(u)$ and a matrix $H$ with
  $P\circ\gamma=H\circ P$ and $E_i(\gamma\cdot,\gamma\cdot)=mE_i$ ($m>0$), such that $Q(y^{-1}w)=HQ(w)$ for the
  periods $Q$ of $\Lambda\otimes\Q$ extending $P$ through $j$. By the real dictionary
  (\ref{pf:ks2:U:dic}) applied to $y^{-1}$, $\kappa(\lambda^{-1})=\kappa(\lambda)^{-1}$ is a mark of $H$.
\item By \lref{UpperTheory.Gal.W2.MinBridge_holds}, the minimal polynomial of $\gamma$ in the sense of
  \cref{pf:ks2:eigtop} is $\mathrm{minpoly}_F(y^{-1})$ (transported to $F_D$ by $L$), which equals
  $\mathrm{minpoly}_F(y)^\vee$, of the same degree and separable if $\mathrm{minpoly}_F(y)$ is
  (\lref{UpperTheory.Gal.W2.InvMin_holds}: the roots of $\mathrm{minpoly}(u^{-1})$ are the inverses of
  those of $\mathrm{minpoly}(u)$). Since $g^{-1}y_fg\in h^{-1}K_N$ is regular and $y$ fixes a line,
  $\mathrm{minpoly}_F(y)$ has degree $n$ and is separable (\lref{UpperTheory.Gal.W2.RegSep_holds}).
\item By \cref{pf:ks2:eigtop} with $\chi=\mathrm{minpoly}_F(y)^\vee$ and the mark $\mu=\kappa(\lambda)^{-1}$:
  $\det(1-T)\ne0$, $\det(1-\Theta)\ne0$ and \eqref{pf:ks2:eq:eigtop} holds.
\item The correspondence acts in period coordinates as required in \cref{pf:ks2:locsum}(a)
  (\lref{UpperTheory.Gal.W2.PowCorr_holds} together with \lref{UpperTheory.Gal.W2.FixEq_holds}, and
  \lref{UpperTheory.Gal.W2.CorrChart_holds} at the base point), and $c_1^{(k)}$ is injective on the fibre
  over $x'$ (\cref{pf:ks2:locsum}(b)): by \lref{UpperTheory.Gal.W2.OneAt_holds}, $c_{1,\mathrm{tot}}$ maps
  the fibre at $u_i(P)$ to the fibre at $c_{1,q}(u_i(P))=u_i(P)$ at level $N$ compatibly with $\exp_{i,P}$, so
  $c_{1,\mathrm{tot}}(\exp_{i,P}z)=c_{1,\mathrm{tot}}(\exp_{i,P}z')$ forces $z-z'\in P(\Lambda_i)$.
\end{enumerate}
\Cref{pf:ks2:locsum} now gives finiteness and nonemptiness of $\mathcal F_{x'}$, $\det(1-A_x)\ne0$, and
\[
\sum_{x\in\mathcal F_{x'}}\mathrm{lef}(A_x)(t)=\kstwoloc\bigl(\alpha;t,\mathrm{minpoly}_F(y)^\vee,\kappa(\lambda)^{-1}\bigr)
=c_W(\alpha,t)\bigl(\mathrm{minpoly}_F(y),\lambda\bigr)
\]
by \eqref{pf:ks2:eq:eigtop} and the definition \eqref{pf:ks2:eq:cw}.
\end{proof}

This completes the proof of \cref{pf:ks2:hfp}, hence of \cref{pf:ks2:fixz} and \cref{pf:ks2:fixk}.

\subsubsection{Period coordinates}\label{pf:ks2:sec:coord}

We prove the statements \ref{pf:ks2:U:dic} and \ref{pf:ks2:U:unifg}. Throughout, $a=m=1$, $v$ is the
archimedean place of $F$ at which $H_v$ has exactly one positive eigenvalue (signature $(1,n-1)$),
$t_0$ is the embedding of the datum, and $\rho$ is of signature $(t_0,\varepsilon_D)$
(\ref{pf:ks2:U:rho}). Put $V_\Q=\Lambda\otimes\Q$; the algebra $C$ acts on $V_\Q$ through the link
$L$, and $y\in\Gamma_\beta$ acts by $w\mapsto y\cdot w$. For a form $\beta$ let
$E_\beta(u,w)=E_\Lambda(\beta^{-1}u,w)$, and for an archimedean place $v'$ let $h^\beta_{v'}(x,x')=
x'^*H_{v'}\psi_{v'}(\beta^{-1})x$ on $\C^n$, a hermitian form, linear in the first variable, with
$\mathcal P_\beta=\{z:h^\beta_v(z,z)>0\}$. Recall the trivializations $(x_\beta,t_\beta)$ of the fixed
representatives $\beta$ of the classes of forms (\cref{pf:ks2:def:forms}).

\begin{definition}[Dictionaries and coordinates]\label{pf:ks2:def:coord}
Let $\beta$ be a form.
\begin{enumerate}
\item A \emph{$p$-lattice} is an $O$-stable subgroup $M\subseteq V_\Q$ for which there is a finite set
  $S$ spanning $V_\Q$ over $\Q$ such that $M$ consists of the $w$ with $mw\in\Z S$ for some integer
  $m$ prime to $p$. For $s\in\Q$, $M$ is \emph{self-dual at $p$ for $sE_\beta$} if $M$ consists of the
  $w\in V_\Q$ such that $sE_\beta(w,u)$ has denominator prime to $p$ for all $u\in M$.
\item A \emph{$p$-adic dictionary} for $\beta$ is a map $g\mapsto M_g$ from $\kstwoGp$ to the set of
  $p$-lattices such that every $M_g$ is self-dual at $p$ for $sE_\beta$ for some rational $s>0$;
  $M_{y_pg}=y\cdot M_g$ for $y\in\Gamma_\beta$; $M_g=M_{g'}$ if and only if $g^{-1}g'\in\kstwoKp$; and
  every $p$-lattice which is self-dual at $p$ for $sE_\beta$ for some rational $s>0$ is some $M_g$.
\item A \emph{period of $V_\Q$} (for $\beta$ and $\rho$) is an additive map $Q\colon V_\Q\to\C^g$ with
  $Q(c\cdot w)=\rho(c)Q(w)$ for $c\in O$, whose image spans $\C^g$ over $\R$, such that
  $\Im H(Qu,Qw)=E_\beta(u,w)$ for some positive definite hermitian form $H$ on $\C^g$.
\item A \emph{real dictionary} for $\beta$ is a relation $Q\sim z$ between periods $Q$ of $V_\Q$ and
  vectors $z\in\C^n$ such that: (line) every period is related to some $z\in\mathcal P_\beta$; (per)
  every $z\in\mathcal P_\beta$ is related to some period; (uniq) if $Q\sim z$ and $Q'\sim z'$ with
  $z,z'\in\mathcal P_\beta$, then $Q'=h\circ Q$ for a $\C$-linear automorphism $h$ of $\C^g$ if and
  only if $z'\in\C^\times z$; (equiv) if $y\in\Gamma_\beta$, $Q'(y\cdot w)=Q(w)$ for all $w$ and
  $Q\sim z$, then $Q'\sim y_vz$; (mark) if $Q\sim z$ with $z\in\mathcal P_\beta$, $y\in\Gamma_\beta$,
  $Q(y\cdot w)=HQ(w)$ for all $w$ and a matrix $H$, and $y_vz=\lambda z$, then $H$ has an eigenvector
  $u\ne0$ with eigenvalue $\kappa(\lambda)$ (\cref{pf:ks2:def:weights}) such that $\rho(c)u=t_0(c)u$
  for all $c\in\kstwoOF$.
\item Let dictionaries $(M^\beta,\sim_\beta)$ be given for the fixed representatives $\beta$, and let
  $N\ge1$. For a polarized $O$-lattice $D'$ with marking $t$ (\cref{pf:ks2:def:periods}), a period $P$
  of $D'$, an additive map $j\colon D'\to V_\Q$ and $a=((g,g_p),z)\in\mathcal D_\beta$, the triple
  $(P,j,a)$ is a system of \emph{coordinates} if: (place) $j$ is injective and $O$-linear, every
  $w\in V_\Q$ has a nonzero integer multiple in $j(D')$, and $E_{D'}(u,u')=sE_\beta(ju,ju')$ for some
  rational $s>0$; (adelic) for $w\in V_\Q$, some $p^rw$ lies in $j(D')$ if and only if
  $g^{-1}x_\beta\cdot w\in\kstwoLaHat$, and $g^{-1}x_\beta\cdot j(t(\lambda))\equiv\lambda\bmod
  N\kstwoLaHat$ for $\lambda\in\Lambda$; ($p$-adic) the $w$ with $mw\in j(D')$ for some $m$ prime to
  $p$ form the lattice $M^\beta_{g_p}$; (real) $P=Q\circ j$ for a period $Q$ of $V_\Q$ with
  $Q\sim_\beta z$.
\end{enumerate}
\formal{UpperTheory.Gal.W2.IsPLat}\formal{UpperTheory.Gal.W2.PerfAt}\formal{UpperTheory.Gal.W2.formB}\formal{UpperTheory.Gal.W2.PDic}\formal{UpperTheory.Gal.W2.IsPerQ}\formal{UpperTheory.Gal.W2.RDic}\formal{UpperTheory.Gal.W2.IsMark}\formal{UpperTheory.Gal.W2.kap}\formal{UpperTheory.Gal.W2.IsPlace}\formal{UpperTheory.Gal.W2.IsCoordA}\formal{UpperTheory.Gal.W2.IsLocP}\formal{UpperTheory.Gal.W2.IsCoord}
\end{definition}

\begin{proposition}[$p$-adic dictionaries]\label{pf:ks2:pdic}
Assume $a=m=1$. For every form $\beta$ there is a $p$-adic dictionary.
\formallabel{UpperTheory.Gal.W2.PDicEx_holds}
\end{proposition}

\begin{proof}
Let $V_p=\prod_wF_w^n$, the product over the local fields $F_w$ of the data of $G$ at $p$, with its
standard $\Q_p$-valued pairing $\langle\cdot,\cdot\rangle$ and its integral elements
$\prod_w\OO_{F_w}^n$. As $m=1$, $\kstwoGp$ is identified with $\Q_p^\times\times\prod_w\GL_n(F_w)$,
and $\kstwoKp$ with the set of $g=(g_1,(g_w))$ such that $|g_1|_p=1$ and the matrices $g_w$ and
$g_w^{-1}$ are integral (this is the definition of $\kstwoKp$, transported along the
identification). By
\lref{UpperTheory.Gal.W2.PFrameEx_holds} there is a \emph{$p$-adic frame} for $\beta$: a pair of
additive maps $T,T'\colon V_\Q\to V_p$, equivariant for the action of $C$ (through its local
components, and twisted by $\beta$ and the involution for $T'$), with
$E_\beta(u,x)=\langle T'x,Tu\rangle-\langle T'u,Tx\rangle$, such that $(T,T')$ has dense image and the
integral points of $(T,T')$ are, up to a fixed power of $p$, the localization at $p$ of a lattice. For
$g\in\kstwoGp$ let
\[
M_g=\{x\in V_\Q:\ g^{-1}T(x)\text{ is integral and }|g_1^{-1}\langle T'(x),gz\rangle|_p\le1\text{ for
all integral }z\in V_p\}.
\]
The following properties are proved in the formal development: each $M_g$ is a $p$-lattice
(\lref{UpperTheory.Gal.W2.PkLat_holds}); if $s\in\Q$ and $|sg_1|_p=1$, then $M_g$ is self-dual at
$p$ for $sE_\beta$ (\lref{UpperTheory.Gal.W2.PkPerf_holds}); $y\cdot M_g=M_{y_pg}$ for
$y\in\Gamma_\beta$ (\lref{UpperTheory.Gal.W2.PkEquiv_holds}); $M_g=M_{g'}$ if and only if
$g^{-1}g'=(u_1,(u_w))$ has $|u_1|_p=1$ and integral $u_w$, $u_w^{-1}$
(\lref{UpperTheory.Gal.W2.PkInj_holds}); and every $p$-lattice which is self-dual at $p$ for some
$sE_\beta$ with $s>0$ is some $M_g$ (\lref{UpperTheory.Gal.W2.PkOnto_holds}).

We check that $g\mapsto M_g$ is a $p$-adic dictionary. Self-duality: given $g$, let $\nu$ be the
$p$-adic valuation of $g_1$ and $s=p^{-\nu}>0$; then $|sg_1|_p=p^{\nu}p^{-\nu}=1$, so $M_g$ is self-dual
at $p$ for $sE_\beta$. Equivariance and surjectivity are two of the listed properties. Injectivity
modulo $\kstwoKp$: $M_g=M_{g'}$ if and only if $g^{-1}g'$ satisfies the integrality condition of
\lref{UpperTheory.Gal.W2.PkInj_holds}, which by the description of $\kstwoKp$ means
$g^{-1}g'\in\kstwoKp$.
\end{proof}

\begin{lemma}[Positive complex structures]\label{pf:ks2:jline}
Let $h(x,y)=y^*Mx$ be the hermitian form of a hermitian matrix $M\in M_n(\C)$, let $\epsilon=\pm1$,
and let $J\in M_n(\C)$ be \emph{positive for $(h,\epsilon)$}: $J^2=-1$, $J^*MJ=M$ (i.e.\ $J$ preserves
$h$), and $\Re(-\epsilon i\,h(Jx,x))>0$ for $x\ne0$. Then:
\begin{enumerate}
\item every $z\ne0$ with $Jz=\epsilon iz$ satisfies $h(z,z)>0$;
\item if $h$ is negative definite, then $J=-\epsilon i$;
\item if $h(z_0,z_0)>0$ for some $z_0$, then $Jz=\epsilon iz$ for some $z\ne0$;
\item if $h(z,z)>0$, $h$ is negative definite on $z^\perp=\{x:h(x,z)=0\}$ and $Jz=\epsilon iz$, then
  $Jx=\epsilon i\bigl(2\frac{h(x,z)}{h(z,z)}z-x\bigr)$ for all $x$.
\end{enumerate}
\formallabel{UpperTheory.Gal.W2.JLine_holds}
\end{lemma}

\begin{proof}
(1) $h(Jz,z)=\epsilon i\,h(z,z)$, so $0<\Re(-\epsilon i\,h(Jz,z))=\Re(\epsilon^2h(z,z))=h(z,z)$. Since
$(J-\epsilon i)(J+\epsilon i)=J^2+1=0$, the image of $J+\epsilon i$ lies in $\ker(J-\epsilon i)$; so if
$J$ has no eigenvector with eigenvalue $\epsilon i$, then $J=-\epsilon i$. (2) By (1), an eigenvector
with eigenvalue $\epsilon i$ would be a positive vector; so there is none, and $J=-\epsilon i$. (3)
Otherwise $J=-\epsilon i$, and then $0<\Re(-\epsilon i\,h(Jx,x))=\Re(\epsilon^2i^2h(x,x))=-h(x,x)$ for
$x\ne0$, so $h$ would be negative definite. (4) As $h(z,z)\ne0$, $\C^n=\C z\oplus z^\perp$. Since
$z=-\epsilon iJz$ and $J$ preserves $h$, for $y\in z^\perp$ we get $h(Jy,z)=h(Jy,-\epsilon
iJz)=\epsilon i\,h(y,z)=0$; so $J$ preserves $z^\perp$, and its restriction is positive for the
restriction of $(h,\epsilon)$, which is negative definite. The argument of (2) on $z^\perp$ gives
$Jy=-\epsilon iy$ for $y\in z^\perp$. Writing $x=cz+y$ with $c=h(x,z)/h(z,z)$ and $y\in z^\perp$,
$Jx=\epsilon icz-\epsilon iy=\epsilon i(2cz-x)$.
\end{proof}

\begin{proposition}[Real dictionaries]\label{pf:ks2:rdic}
Assume $a=m=1$, and let $v,t_0,\rho$ be as above. For every form $\beta$ there is a real dictionary.
\formallabel{UpperTheory.Gal.W2.RDicEx_holds}
\end{proposition}

\begin{proof}
Let $\mathrm{Arch}=\prod_{v'}\C^n$, the product over the archimedean places $v'$ of $F$; for $v'$ let
$\epsilon_{v'}=1$ if $v'$ maps $\omega$ (via $L$) to $\rho_c$, and $\epsilon_{v'}=-1$ otherwise; and
write $e_v(z)\in\mathrm{Arch}$ for the vector with component $z$ at $v$ and $0$ elsewhere. Positivity
of complex structures is meant as in \cref{pf:ks2:jline}. We use the following
statements, proved in the formal development.
\begin{enumerate}[label=(R\arabic*)]
\item\label{pf:ks2:R:arch} (\lref{UpperTheory.Gal.W2.ArchEx_holds}) There is an \emph{archimedean
  frame}: an additive map $\theta\colon V_\Q\to\mathrm{Arch}$ with $\theta(c\cdot
  w)_{v'}=\psi_{v'}(c)\theta(w)_{v'}$ for $c\in C$, whose image spans $\mathrm{Arch}$ over $\R$, with
  $\dim_\R\mathrm{Arch}=\dim_\Q V_\Q=2g$, such that $E_\Lambda$ is a positive linear combination of the
  forms $(u,w)\mapsto\Re(-\epsilon_{v'}i\,w_{v'}^*H_{v'}u_{v'})$ evaluated at $\theta(u),\theta(w)$.
\item\label{pf:ks2:R:sig} (\lref{UpperTheory.Gal.W2.SigB_holds}) The forms $h^\beta_{v'}$ are
  hermitian; at $v$, $h^\beta_v$ has a positive vector and is negative definite on the orthogonal
  complement of every positive vector; at every place $v'$ at which $H_{v'}$ does not have exactly one
  positive eigenvalue, $h^\beta_{v'}$ is negative definite.
\item\label{pf:ks2:R:ext} (\lref{UpperTheory.Gal.W2.PerExt_holds}) Every period $Q$ of $V_\Q$ is
  $R\circ\theta$ for an $\R$-linear isomorphism $R\colon\mathrm{Arch}\to\C^g$.
\item\label{pf:ks2:R:J} (\lref{UpperTheory.Gal.W2.JStruct_holds}) For such $Q$ and $R$ there are
  matrices $J_{v'}$, positive for $(h^\beta_{v'},\epsilon_{v'})$, with $R(Jx)=iR(x)$ for
  $x\in\mathrm{Arch}$, where $J$ acts on $\mathrm{Arch}$ by $J_{v'}$ on the component $v'$.
\item\label{pf:ks2:R:line} (\cref{pf:ks2:jline}) Let $J$ be positive for $(h,\epsilon)$. Then (i) every
  $z\ne0$ with $Jz=\epsilon iz$ has $h(z,z)>0$; (ii) if $h$ is negative definite, $J=-\epsilon i$; (iii)
  if $h$ has a positive vector, some $z\ne0$ has $Jz=\epsilon iz$; (iv) if $h(z,z)>0$, $h$ is negative
  definite on $z^\perp$ and $Jz=\epsilon iz$, then $Jx=\epsilon i\bigl(2\frac{h(x,z)}{h(z,z)}z-x\bigr)$
  for all $x$.
\item\label{pf:ks2:R:per} (\lref{UpperTheory.Gal.W2.RPer_holds}) For every $z\in\mathcal P_\beta$
  there are a period $Q$ and an $\R$-linear $R$ with $Q=R\circ\theta$ and
  $R(e_v(iz))=\epsilon_viR(e_v(z))$.
\item\label{pf:ks2:R:place} (\lref{UpperTheory.Gal.W2.PlaceT_holds}) $v$ is the only place at which
  $H_{v'}$ has exactly one positive eigenvalue; and either $\epsilon_v=1$ and $v$ induces $t_0$ (via
  $L$), or $\epsilon_v=-1$ and $v$ induces $\bar t_0\ne t_0$.
\end{enumerate}
Define $Q\sim z$ if $Q=R\circ\theta$ for an $\R$-linear map $R\colon\mathrm{Arch}\to\C^g$ with
$R(e_v(iz))=\epsilon_viR(e_v(z))$. If $Q$ is a period, such an $R$ is the isomorphism of
\ref{pf:ks2:R:ext} (it is determined by $Q$, as $\theta(V_\Q)$ spans $\mathrm{Arch}$); with $J$ as in
\ref{pf:ks2:R:J}, $R(Je_v(z))=iR(e_v(z))$, so $Q\sim z$ if and only if $J_vz=\epsilon_viz$ (as
$\epsilon_v^2=1$ and $R$ is injective). We verify the five conditions.

(line) Let $Q$ be a period, with $R$ and $J$. By \ref{pf:ks2:R:sig} $h^\beta_v$ has a positive vector,
so by \ref{pf:ks2:R:line}(iii) there is $z\ne0$ with $J_vz=\epsilon_viz$; it is positive by
\ref{pf:ks2:R:line}(i), so $z\in\mathcal P_\beta$ and $Q\sim z$.

(per) This is \ref{pf:ks2:R:per}.

(uniq) Let $Q\sim z$ and $Q'\sim z'$, with $R,J$ and $R',J'$. If $Q'=h\circ Q$ with $h$ $\C$-linear,
then $R'=h\circ R$ (both sides agree on $\theta(V_\Q)$), hence $J'=J$. As $J_vz=\epsilon_viz$, $z$ is
positive and $h^\beta_v$ is negative definite on $z^\perp$ (\ref{pf:ks2:R:sig}),
\ref{pf:ks2:R:line}(iv) gives $J_vx=\epsilon_vi(2\frac{h(x,z)}{h(z,z)}z-x)$ for all $x$, with
$h=h^\beta_v$. Applied to $x=z'$, for which $J_vz'=J'_vz'=\epsilon_viz'$, this yields
$z'=\frac{h(z',z)}{h(z,z)}z$, and the coefficient is nonzero as $z'\ne0$. Conversely, let $z'=cz$ with
$c\ne0$; then $J'_vz=\epsilon_viz$. By \ref{pf:ks2:R:line}(iv), $J_v$ and $J'_v$ are both given by the
formula in terms of $z$, hence equal; for $v'\ne v$, $h^\beta_{v'}$ is negative definite by
\ref{pf:ks2:R:sig} and \ref{pf:ks2:R:place}, so $J_{v'}=-\epsilon_{v'}i=J'_{v'}$ by
\ref{pf:ks2:R:line}(ii). Thus $J=J'$, and $h=R'\circ R^{-1}$ commutes with multiplication by $i$, i.e.\
is $\C$-linear, with $Q'=h\circ Q$.

(equiv) Let $y\in\Gamma_\beta$ act on $\mathrm{Arch}$ by $y_\infty=(\psi_{v'}(y))_{v'}$, so that
$\theta(y\cdot w)=y_\infty\theta(w)$. If $Q'(y\cdot w)=Q(w)$ for all $w$ and $Q=R\circ\theta$, then
$Q'=R'\circ\theta$ with $R'=R\circ y_\infty^{-1}$, and
$R'(e_v(iy_vz))=R(e_v(iz))=\epsilon_viR(e_v(z))=\epsilon_viR'(e_v(y_vz))$; so $Q'\sim y_vz$.

(mark) Let $Q\sim z$ with $z\in\mathcal P_\beta$, $Q(y\cdot w)=HQ(w)$ and $y_vz=\lambda z$, and put
$u=R(e_v(z))$, which is nonzero. The $\R$-linear maps $H\circ R$ and $R\circ y_\infty$ agree on
$\theta(V_\Q)$, hence everywhere, so
\[
Hu=R(e_v(\lambda z))=\Re(\lambda)R(e_v(z))+\Im(\lambda)R(e_v(iz))=(\Re\lambda+\epsilon_vi\Im\lambda)u .
\]
Similarly, for $c\in\kstwoOF$ we have $\rho(c)\circ R=R\circ c_\infty$, and $c$ acts on the component
$v$ by the scalar $v(c)$, so $\rho(c)u=(\Re v(c)+\epsilon_vi\Im v(c))u$. By \ref{pf:ks2:R:place},
$\Re\lambda+\epsilon_vi\Im\lambda=\kappa(\lambda)$ and $\Re v(c)+\epsilon_vi\Im v(c)=t_0(c)$. So $u$
is an eigenvector of $H$ with eigenvalue $\kappa(\lambda)$ on which $\kstwoOF$ acts through $t_0$.
\end{proof}

\begin{proposition}[Period coordinates]\label{pf:ks2:unifg}
Assume $a=m=1$, let $v,t_0,\rho$ be as above, fix a $p$-adic and a real dictionary for every form, and
let $N\ge1$ be prime to $p$. Let $\mathcal A\to X$ be a family over $\C$ with a uniformization by
marked period lattices $(\Lambda_i,t_i)$ in $\mathrm{Cls}_N$ and points $u_i(P)$ as in
\ref{pf:ks2:U:unif}. Then there is a bijection $\Phi\colon X(\C)\to\mathcal X(K_N,\kstwoKp)$ such that
for every complex point $x$ there are a class $[\beta]$, an element $a\in\mathcal D_\beta$, an index
$i$, a period $P$ of $\Lambda_i$ and an additive map $j\colon\Lambda_i\to V_\Q$ with
$\Phi(x)=([\beta],[a])$, $x=u_i(P)$ and $(P,j,a)$ coordinates.
\formallabel{UpperTheory.Gal.W2.UnifG_holds}
\end{proposition}

\begin{proof}
We use four statements proved in the formal development. (Existence) Every marked lattice $(D',t)$
in $\mathrm{Cls}_N$ with a period $P$ admits coordinates $(P,j,a)$ for some class $[\beta]$,
$a\in\mathcal D_\beta$ and $j$ (\lref{UpperTheory.Gal.W2.CoordEx_holds}). (Occurrence) For every
form $\beta$ and $a\in\mathcal D_\beta$ there are a marked lattice $(D',t)$ in $\mathrm{Cls}_N$, a
period $P$ of $D'$ and a map $j$ such that $(P,j,a)$ are coordinates
(\lref{UpperTheory.Gal.W2.CoordOcc_holds}). (Comparison) If $(P,j,a)$ and $(P'',j'',a'')$ are
coordinates for $(D',t)$ and $(D'',t'')$ in $\mathrm{Cls}_N$, with $a\in\mathcal D_\beta$ and
$a''\in\mathcal D_{\beta''}$, then $([\beta],[a])=([\beta''],[a''])$ in $\mathcal X(K_N,\kstwoKp)$ if
and only if there are a marked isometry $w\colon(D'',t'')\to(D',t)$ and a $\C$-linear automorphism $h$
of $\C^g$ with $P\circ w=h\circ P''$ (\lref{UpperTheory.Gal.W2.CoordIso_holds}). (Transport) If
$w\colon(D'',t'')\to(D',t)$ is a marked isometry and $(P,j,a)$ are coordinates for $(D',t)$, then
$(P\circ w,j\circ w,a)$ are coordinates for $(D'',t'')$ (\lref{UpperTheory.Gal.W2.CoordMove_holds}).

For a complex point $x$ choose $i$ and a period $P$ of $\Lambda_i$ with $x=u_i(P)$, then by (Existence)
coordinates $(P,j,a)$ with $a\in\mathcal D_\beta$, and put $\Phi(x)=([\beta],[a])$. The asserted
properties of $\Phi(x)$ hold by construction.

$\Phi$ is injective: let $\Phi(x)=\Phi(x')$, where $x=u_i(P)$ and $x'=u_{i'}(P')$. By (Comparison)
there are a marked isometry $w\colon\Lambda_{i'}\to\Lambda_i$ and $h$ with $P\circ w=h\circ P'$. As
distinct indices are not marked-isometric, $i=i'$; then $P'\circ w^{-1}=h^{-1}\circ P$ with $w^{-1}$ a
marked self-isometry, so $u_i(P)=u_i(P')$ by \ref{pf:ks2:U:unif}, i.e.\ $x=x'$.

$\Phi$ is surjective: let $([\beta],[a])\in\mathcal X(K_N,\kstwoKp)$ with $\beta$ the fixed
representative. By (Occurrence) there are $(D',t)$ in $\mathrm{Cls}_N$, a period $P'$ and $j'$ with
$(P',j',a)$ coordinates. By \ref{pf:ks2:U:unif}, $(D',t)$ is marked-isometric to some
$(\Lambda_{i_0},t_{i_0})$, say by $w\colon D'\to\Lambda_{i_0}$. Then $P_0=P'\circ w^{-1}$ is a period of
$\Lambda_{i_0}$ and $(P_0,j'\circ w^{-1},a)$ are coordinates by (Transport). Let $x=u_{i_0}(P_0)$, and let
$\Phi(x)=([\beta_x],[a_x])$ be defined through $x=u_{i_x}(P_x)$ and coordinates $(P_x,j_x,a_x)$. As
$u_{i_0}(P_0)=u_{i_x}(P_x)$, \ref{pf:ks2:U:unif} gives $i_x=i_0$, a marked self-isometry $w'$ of
$\Lambda_{i_0}$ and $h'$ with $P_x\circ w'=h'\circ P_0$; by (Comparison),
$([\beta],[a])=([\beta_x],[a_x])=\Phi(x)$.
\end{proof}

% ---------------------------------------------------------------------------
\subsection{Classical inputs}\label{pf:ks2:sec:classical}

The following statements are classical; they are proved in the formal development, and we cite
the literature for them rather than reproducing their (long) formal proofs. They are stated in the
form in which they are used above. For two of them we also give the top layer of the formal proof:
the extension of objects over discrete valuation rings (\cref{pf:ks2:extobj}) and the
classification of complex points of models (\cref{pf:ks2:finecx}).

\begin{littheorem}[Properness in the compact case; Kottwitz, Lan]\label{pf:ks2:classical:proper}
Let $R$ be a Noetherian domain of characteristic $0$ with a base $b$ of a datum $D$ of compact
type, $N\ge1$, and $M$ a model of level $N$ over $R$. Then $M\to\Spec R$ is proper.
\formallabel{GlobalTheory.T.harrisTaylor_integralModel_isProper_of_compact}
\end{littheorem}

For PEL moduli problems of this kind the properness of the moduli space when the associated group
is anisotropic modulo its centre is classical \cite{Kottwitz1992}; see \cite{Lan2013} for a general
treatment. The formal proof checks the valuative criterion for complete discrete valuation rings
with fraction field of characteristic $0$; potential good reduction of the abelian varieties
parametrized by $M$ comes from Grothendieck's semistable reduction theorem
\cite[Exp.~IX]{Grothendieck1972}, a nontrivial toric part of the reduction being excluded by the
compact-type condition together with Kottwitz's condition; then all the additional structures
extend (\cref{pf:ks2:extobj}).

\begin{proposition}[Extension of objects over discrete valuation rings]\label{pf:ks2:extobj}
Let $R$ be a ring with a base $b$ of $D$, $N\ge1$, $W$ a discrete valuation ring with fraction field
$L$ and $R\to W$ a ring map. Let $o$ be an object of level $N$ over $\Spec L$, $G$ an abelian scheme
over $\Spec W$, and $\phi\colon A_o\to G$ a homomorphism identifying $A_o$ with the generic fibre of
$G$. Then there are an object $o'$ of level $N$ over $\Spec W$ with abelian scheme $G$ and a morphism of
objects $o\to o'$ over $\Spec L\to\Spec W$.
\formallabel{GlobalTheory.HTObjExt.extends_obj}
\end{proposition}

\begin{proof}
The morphism $o\to o'$ will be $\phi$; we extend the remaining structures of $o$ to $G$, each
uniquely determined by its restriction to the generic fibre.
\begin{enumerate}
\item ($O$-action.) Endomorphisms of the generic fibre of an abelian scheme over a discrete
  valuation ring extend uniquely, so the action of $O$ extends to an action on $G$ compatible with
  $\phi$ (\lref{GlobalTheory.T.exists_ring_action_extension_compatible}).
\item (Polarization bundle.) The bundle $\mathcal P$ of $o$ on $A_o\times_LA_o$, trivial along both
  axes, extends to a line bundle $\mathcal P'$ on $G\times_WG$, trivial along both axes, with an
  isomorphism $\mathcal P'_L\cong\mathcal P$
  (\lref{GlobalTheory.AxisTrivExt.exists_extension_square_of_axes_trivial}). Isomorphisms between
  line bundles on $G\times_WG$ which are trivial along the zero sections extend from the generic fibre
  (\lref{GlobalTheory.T.lineBundleIso_pullback_extend_DVR}); this gives the rigidifications, the
  symmetry of $\mathcal P'$ and the Rosati condition from those of $\mathcal P$.
\item (Ampleness.) The restriction of $\mathcal P'$ to the diagonal is relatively ample over $W$,
  because that of $\mathcal P$ is ample
  (\lref{GlobalTheory.T.relAmple_diagonal_extendedBundle_of_generic}).
\item (Kernel of the polarization.) If the kernel of the polarization of $o$ is killed by an integer
  $e$ prime to $p$, so is that of $\mathcal P'$
  (\lref{GlobalTheory.T.polarisation_kernel_killed_over_dvr_of_generic}).
\item (Level structure.) As $G$ is proper over $W$, restriction to the generic fibre is a group
  isomorphism from the sections of $G$ over $W$ onto those of $A_o$ over $L$
  (\lref{GlobalTheory.T.properGroupScheme_dvr_sections_mulEquiv_generic}, the valuative
  criterion). So $\eta$ extends uniquely to $\eta'\colon\Lambda\to G(W)$, which is additive,
  $O$-linear and kills $N\Lambda$ because $\eta$ is.
\item (Kottwitz's condition.) The determinant condition on $\Lie G$ holds because it holds on the
  generic fibre (\lref{GlobalTheory.T.harrisTaylor_lie_condition_extends_abelianScheme_dvr}); here
  one uses that the discriminant $\mathrm{tr}^2-4\mathrm{nm}$, being prime to $p$, is a unit in $W$.
\item (Tate condition.) At every geometric point of $\Spec W$, $\eta'$ is the reduction of a
  prime-to-$p$ trivialization of the Tate module compatible with $O$ and with the Weil pairing up to a
  unit (\lref{GlobalTheory.T.tateSimilitude_geometric_points_of_generic_dvr}: integers prime to
  $p$ are units in $R$, hence in $W$, so the torsion of $G$ of order prime to $p$ is finite \'etale over
  $W$, and trivializations at geometric points of the special fibre are obtained from those of the
  generic fibre by specialization).
\end{enumerate}
These data form an object $o'$ over $\Spec W$, and $\phi$, with the induced map on $A_o\times_LA_o$, is
a morphism of objects $o\to o'$.
\end{proof}

\begin{littheorem}[Isogenies of type $T_j$]\label{pf:ks2:classical:tj}
Let $R$ be Noetherian with a base $b$, $N$ a level, $p$ invertible in $R$, $m=1$, $\mathfrak w\supseteq
(p,\omega-r)$ maximal, $0\le j\le n$ and $c=p^e$ with $e>0$ if $j>0$. For every object $o$ of level
$N$ over a locally Noetherian $R$-scheme $S$ there is a finite \'etale $S$-scheme $Y$ carrying a
$p$-isogeny of degree $c$ and type $T_j$ at $\mathfrak w$ out of the pullback of $o$, which is
universal: every such isogeny out of the pullback of $o$ along $x\colon T\to S$ is the pullback
along a unique $T\to Y$ over $x$.
\formallabel{GlobalTheory.T.harrisTaylor_typeTj_isogenies_representable}
\end{littheorem}

Where $p$ is invertible, isogenies of degree a power of $p$ out of $A_o$ correspond to their kernels,
which are finite \'etale subgroup schemes of the $p$-power torsion of the required type, and
quotients of abelian varieties by finite subgroup schemes exist \cite{Mumford1970}; the
representability of such kernels by a finite \'etale scheme is standard.

\begin{littheorem}[Generic ordinariness; Wedhorn]\label{pf:ks2:classical:ord}
Let $R$ be a discrete valuation domain of characteristic $0$ with $p$ in its maximal ideal, $b$ a
base over $R$ with $m=1$, $N$ a level, $\mathfrak w$ the inverse image of the maximal ideal of $R$
under some $t_i$, and $M$ a model of level $N$ over $R$. If $z\in M$ lies over the closed point of
$\Spec R$ and has no proper generization in the special fibre, then every geometric point $y$ of $M$
with image $z$ is \emph{$\mathfrak w$-ordinary}: the group of $\mathfrak w$-torsion points of the abelian
variety at $y$ has exactly $\Nm\mathfrak w^{\,n-1}$ elements.
\formallabel{UpperTheory.Gal.Ord.ordGen}
\end{littheorem}

This is a form of the density of the ordinary locus in the special fibre of PEL-type Shimura
varieties with good reduction at primes split in the reflex field \cite{Wedhorn1999}. The formal
proof works with the action of Frobenius on the de Rham cohomology of the special fibre and the
Hodge filtration, and uses the following classical fact.

\begin{littheorem}[Characteristic polynomials on de Rham homology]\label{pf:ks2:classical:charpoly}
Let $A$ be an abelian variety over an algebraically closed field $k$, $f$ an endomorphism of $A$, and
$\ell\ne\mathrm{char}\,k$ a prime. Let $(L,\varphi)$ be a free $\Z$-module of finite rank with an
endomorphism, together with compatible surjections $L\to A[\ell^j](k)$ with kernel $\ell^jL$
intertwining $\varphi$ and $f$ (a model of the Tate module $T_\ell A$ with the action of $f$). If the
first de Rham homology of $A$ is finite-dimensional over $k$, then the characteristic polynomial of
$f$ on it is the image in $k[X]$ of the characteristic polynomial of $\varphi$.
\formallabel{GlobalTheory.T.charpoly_deRhamHomology_endomorphism_eq_map_ladic}
\end{littheorem}

The characteristic polynomial of $f$ on $T_\ell A$ has integer coefficients independent of $\ell$
\cite{Mumford1970}; in characteristic $0$ the statement follows from the comparison of de Rham and
Betti cohomology, and in characteristic $p$ from the relation between de Rham cohomology and
Dieudonn\'e modules \cite{Oda1969} (compare \cite{KatzMessing1974}).

\begin{littheorem}[Complex uniformization]\label{pf:ks2:classical:unif}
Assume $a=m=1$, let $t_0$ be the embedding of the datum, $\rho$ as in \ref{pf:ks2:U:rho}, $g_b$ a
standard base over $\C$, $N$ a level and $X$ a generic model of level $N$ over $\C$. Then $X$ is smooth
of relative dimension $n-1$, quasi-compact and separated over $\C$, the universal abelian scheme is
smooth of relative dimension $g$, and the complex points of $X$ and of the universal abelian scheme are
uniformized by marked period lattices as described in \ref{pf:ks2:U:unif}.
\formallabel{UpperTheory.Gal.W2.GUnif_holds}
\end{littheorem}

This is the classical description of the complex points of PEL moduli spaces by lattices with
additional structure \cite{Shimura1971,Kottwitz1992}; the formal proof
(\cref{pf:ks2:finecx} below, which classifies the complex points of a model by marked period
lattices, and \lref{GlobalTheory.T.complexUniformisation_of_markedPeriods_family}, which
constructs the period coordinates in families) proceeds through the analytic theory of complex tori.

\begin{proposition}[Complex points of models]\label{pf:ks2:finecx}
Let $R$ be a ring with a base $b$ of $D$, $N\ge1$ prime to $p$, $\sigma\colon R\to\C$ a ring map and
$\rho\colon O\to M_g(\C)$ a ring map such that $b$ is \emph{of type $(\sigma,\rho)$}:
\[
\tr\rho(c)=m\Bigl(\sum_i\sigma(t_i(c))+n\sum_{\tau(\omega)\ne\sigma(\rho_b)}\tau(c)
-\sum_i\overline{\sigma(t_i(c))}\Bigr)\qquad(c\in\kstwoOF),
\]
the middle sum over the embeddings $\tau\colon F\to\C$ with $\tau(\omega)\ne\sigma(\rho_b)$. Let $M$
be a model of level $N$ over $R$, with universal object $\mathcal A$. Then, in the sense of
\cref{pf:ks2:def:periods}:
\begin{enumerate}
\item for every complex point $y$ of $M$ over $\sigma$ there are $(D',t)$ in $\mathrm{Cls}_N$, a period
  $P$ of $D'$ and an exponential $e$ such that $(D',t,P)$ uniformizes $\mathcal A$ at $y$ via $e$;
\item for every $(D',t)$ in $\mathrm{Cls}_N$ and every period $P$ of $D'$ there are a complex point $y$
  of $M$ over $\sigma$ and an exponential $e$ such that $(D',t,P)$ uniformizes $\mathcal A$ at $y$ via
  $e$;
\item if $(D',t,P)$ uniformizes $\mathcal A$ at $y$ and $(D'',t'',P'')$ at $y'$ (via some
  exponentials), and there are a marked isometry $w\colon(D',t)\to(D'',t'')$ and a $\C$-linear
  automorphism $h$ of $\C^g$ with $P''\circ w=h\circ P$, then $y=y'$.
\end{enumerate}
\formallabel{GlobalTheory.T.complexPoints_classify_markedPeriods}
\end{proposition}

\begin{proof}
For a complex point $y$ over $\sigma$, an object $o$ over $\Spec\C$ (via $\sigma$) and a morphism of
objects $o\to\mathcal A$ over $y$, the uniformization conditions for $\mathcal A$ at $y$ and for $o$
at the identity point are equivalent, for the same $(D',t,P)$: all conditions only involve the fibre,
which the morphism identifies with $A_o$ compatibly with the $O$-action, the polarization and the
level structure.

(1) Pull $\mathcal A$ back along $y$ to an object $o$ over $\C$. By
\lref{GlobalTheory.T.exists_complexUniformisation_harrisTaylor} (an object of level $N$ over
$\C$ for a base of type $(\sigma,\rho)$ is uniformized by a marked period lattice in
$\mathrm{Cls}_N$: the lattice of periods $H_1(A_o(\C),\Z)$ with the Riemann form of the polarization,
the period map, and the marking given by the level structure) $o$ is uniformized at the identity by
some $(D',t,P)$, hence so is $\mathcal A$ at $y$.

(2) By \lref{GlobalTheory.T.exists_harrisTaylor_object_of_period} (the complex torus
$\C^g/P(D')$, with $O$ acting through $\rho$, the polarization given by the Riemann form $E_{D'}$ and
the level structure $\lambda\mapsto P(t(\lambda))/N$, is an object of level $N$ over $\C$, uniformized
by $(D',t,P)$) there is an object $o$ over $\C$ uniformized at the identity by $(D',t,P)$. The
universal property of $M$ gives $y\colon\Spec\C\to M$ over $\sigma$ with a morphism $o\to\mathcal A$
over $y$, so $\mathcal A$ is uniformized by $(D',t,P)$ at $y$.

(3) Pull $\mathcal A$ back along $y$ and $y'$ to objects $o_1,o_2$ over $\C$, uniformized at the
identity by $(D',t,P)$ and $(D'',t'',P'')$. By
\lref{GlobalTheory.T.nonempty_moduliMap_of_marked_period_isometry} (the $\C$-linear map $h$,
which carries $P(D')$ onto $P''(D'')$, descends to an isomorphism of complex tori respecting the
$O$-actions, the polarizations and the level structures) there is a morphism of objects $o_1\to o_2$
over $\mathrm{id}_{\Spec\C}$. Composing it with $o_2\to\mathcal A$ over $y'$ gives a morphism
$o_1\to\mathcal A$ over $y'$, and $o_1\to\mathcal A$ over $y$ is another; by the uniqueness in the
universal property of $M$, $y=y'$.
\end{proof}

% ---------------------------------------------------------------------------
\subsection{Remarks on the formal statements}\label{pf:ks2:sec:remarks}

\begin{remark}[Names]\label{pf:ks2:rem:names}
In the formal names of the main theorem and of $\sfname{KSAll}$, ``KS'' and ``Kodaira--Spencer''
refer to Kuga--Sato varieties: no Kodaira--Spencer map occurs in this section. The words ``weight
one'' in the name of the main theorem do not refer to modular forms of weight one; the hypothesis of
the theorem is the holomorphic Lefschetz formula.
\end{remark}

\begin{remark}[An unused hypothesis]\label{pf:ks2:rem:unused}
In the formal proof, $\sfname{StdEx}$ is obtained from a lemma whose hypothesis is the integral input
$\sfname{MIn}$ of \cref{pf:ks2:min}; this hypothesis is not used, and the lemma is
\cref{pf:ks2:stdex}.
\formal{UpperTheory.Gal.Int.Geo.CovMain.stdEx_of}
\end{remark}

\begin{remark}[Hypotheses of $\sfname{FixK}$]\label{pf:ks2:rem:fixk}
(1) The hypothesis that $1$ is admissible from $N'$ to $N$ follows from $N\mid N'$.
(2) The divisibility of the $\alpha_s$ by $2$ is used only to guarantee that the endomorphism
$\Theta$ of the fibres has $\det(1-\Theta)\ne0$, i.e.\ that the fixed points in the fibres are
transversal (\cref{pf:ks2:eigtop}); their nonvanishing is used only for the finiteness and
\'etaleness of $c_2$. The moduli version \cref{pf:ks2:hfp} does not require $\alpha_s\ne0$.
(3) The integer $\delta$ is arbitrary in the statement; the proof shows $\delta=n-1+kg$ when $Z_N$ is
nonempty.
(4) The weight $c_W$ in the formal statement of $\sfname{FixK}$ is defined with the conjugation
convention attached to the embedding of $\omega$ at $v$, and the weight used in the proof with the
convention attached to the comparison of $v$ with $t_0$; they agree under $\mathrm{Sc}(D,G)$
(\lref{UpperTheory.Gal.FixZ.wtEq_of}, which states that the two weight functions coincide when
$a=1$ and $v$ has signature $(1,n-1)$).
\end{remark}

\begin{remark}[Finite sums]\label{pf:ks2:rem:finsum}
The orbital sums \eqref{pf:ks2:eq:cxtotal} are finite sums in the sense of Mathlib, which assign the
value $0$ to a sum with infinite support. The finiteness statements \cref{pf:ks2:def:fixk}(1) and
\cref{pf:ks2:hfp}(3) show that, in the situations where they are used, the supports are finite and
the sums are ordinary finite sums.
\end{remark}

\begin{remark}[Pointwise congruence]\label{pf:ks2:rem:ptid}
The congruence relation \cref{pf:ks2:congks} is an equality of multisets of residue points over the
generic points of the special fibre, for every valuation ring with algebraically closed fraction field;
it is not stated as an identity of correspondences or of cycles. Ordinariness is defined by counting
$\mathfrak w$-torsion points ($\Nm\mathfrak w^{\,n-1}$ of them), and the line $C$ of the ordinary reading is
the set of $\mathfrak w$-torsion points of the generic fibre reducing to the origin.
\end{remark}

\begin{remark}[Generic models]\label{pf:ks2:rem:generic}
A generic model is required only to be a fine moduli scheme; its smoothness, quasi-compactness,
separatedness and (in the compact case) projectivity are proved in \cref{pf:ks2:gensm} by comparison
with an integral model.
\end{remark}

